% ================================================================
% APUNTES INTEGRADOS DE DERECHO DE LA INTEGRACIÓN
% Documento único, autónomo y compilable, resultado de la
% integración editorial de diez clases/apuntes originales.
% Compilar con pdflatex (dos pasadas, para índice y referencias).
% ================================================================
\documentclass[12pt,oneside]{book}

% ----------------------------------------------------------------
% METADATOS
% ----------------------------------------------------------------
\newcommand{\universidad}{Universidad de Buenos Aires (UBA)}
\newcommand{\facultad}{Facultad de Derecho}
\newcommand{\carrera}{Carrera de Abogacía}
\newcommand{\materia}{Derecho de la Integración}
\newcommand{\titulo}{Derecho de la Integración}
\newcommand{\subtitulo}{Apuntes integrados: de los fundamentos teóricos al Mercosur y la Unión Europea}
\newcommand{\autor}{Isaac Edgar Camacho Ocampo}
\newcommand{\anio}{2026}

% ----------------------------------------------------------------
% PAQUETES
% ----------------------------------------------------------------
\usepackage[utf8]{inputenc}
\usepackage[T1]{fontenc}
\usepackage[spanish,provide=*]{babel}

\usepackage[
    top=2.5cm,
    bottom=2.5cm,
    left=3cm,
    right=2.5cm,
    headheight=15pt
]{geometry}

\usepackage{amsmath}
\usepackage{amssymb}
\usepackage{mathtools}

\usepackage{graphicx}
\usepackage{float}
\usepackage{caption}
\usepackage{subcaption}

\usepackage{booktabs}
\usepackage{array}
\usepackage{longtable}
\usepackage{tabularx}

\usepackage{enumitem}

\usepackage{xcolor}
\usepackage[most]{tcolorbox}

\usepackage{fancyhdr}
\usepackage{ragged2e}

\usepackage[
    colorlinks=true,
    linkcolor=blue!50!black,
    citecolor=green!40!black,
    urlcolor=blue!60!black,
    pdftitle={\titulo},
    pdfauthor={\autor},
    pdfsubject={Apuntes universitarios integrados},
    bookmarks=true,
    plainpages=false
]{hyperref}

% ----------------------------------------------------------------
% CONFIGURACIÓN GENERAL
% ----------------------------------------------------------------
\graphicspath{{./img/}{./graficos/}{./figuras/}}

\addto\captionsspanish{\renewcommand{\contentsname}{Índice general}}

\setcounter{tocdepth}{2}

\setlength{\parindent}{0pt}
\setlength{\parskip}{0.5em}

\setlist[itemize]{topsep=0.3em, itemsep=0.2em}
\setlist[enumerate]{topsep=0.3em, itemsep=0.2em}

\clubpenalty=10000
\widowpenalty=10000

\pagestyle{fancy}
\fancyhf{}
\fancyhead[L]{\nouppercase{\leftmark}}
\fancyhead[R]{\materia}
\fancyfoot[C]{\thepage}
\renewcommand{\headrulewidth}{0.4pt}
\renewcommand{\footrulewidth}{0pt}
\renewcommand{\chaptermark}[1]{\markboth{\thechapter.\ #1}{}}

\fancypagestyle{plain}{
    \fancyhf{}
    \fancyfoot[C]{\thepage}
    \renewcommand{\headrulewidth}{0pt}
}

% ----------------------------------------------------------------
% COMANDOS MATEMÁTICOS (previstos por la especificación; de uso
% eventual en tablas de aranceles o fórmulas de reparto de bancas)
% ----------------------------------------------------------------
\newcommand{\abs}[1]{\left|#1\right|}
\newcommand{\norm}[1]{\left\lVert#1\right\rVert}
\newcommand{\vect}[1]{\mathbf{#1}}
\newcommand{\deriv}[2]{\frac{d#1}{d#2}}
\newcommand{\pderiv}[2]{\frac{\partial#1}{\partial#2}}

% ----------------------------------------------------------------
% ENTORNOS / CAJAS ACADÉMICAS SEMÁNTICAS
% ----------------------------------------------------------------
\newtcolorbox{definicion}[1]{
    enhanced, breakable, title={Definición: #1}, fonttitle=\bfseries,
    colback=blue!3, colframe=blue!50!black, boxrule=0.8pt, arc=2mm,
    before upper={\setlength{\parskip}{0.5em}}
}
\newtcolorbox{importante}{
    enhanced, breakable, title={Importante}, fonttitle=\bfseries,
    colback=yellow!8, colframe=orange!70!black, boxrule=0.8pt, arc=2mm,
    before upper={\setlength{\parskip}{0.5em}}
}
\newtcolorbox{ejemplo}[1]{
    enhanced, breakable, title={Ejemplo: #1}, fonttitle=\bfseries,
    colback=green!4, colframe=green!40!black, boxrule=0.8pt, arc=2mm,
    before upper={\setlength{\parskip}{0.5em}}
}
\newtcolorbox{formula}[1]{
    enhanced, breakable, title={Fórmula / Regla: #1}, fonttitle=\bfseries,
    colback=purple!4, colframe=purple!50!black, boxrule=0.8pt, arc=2mm,
    before upper={\setlength{\parskip}{0.5em}}
}
\newtcolorbox{ejercicio}[1]{
    enhanced, breakable, title={Ejercicio: #1}, fonttitle=\bfseries,
    colback=gray!5, colframe=gray!60!black, boxrule=0.8pt, arc=2mm,
    before upper={\setlength{\parskip}{0.5em}}
}
\newtcolorbox{nota}{
    enhanced, breakable, title={Nota}, fonttitle=\bfseries, coltitle=black,
    colback=white, colframe=gray!50, boxrule=0.6pt, arc=2mm,
    before upper={\setlength{\parskip}{0.5em}}
}
\newtcolorbox{normacitada}[1]{
    enhanced, breakable, title={Norma citada: #1}, fonttitle=\bfseries,
    colback=red!3, colframe=red!45!black, boxrule=0.8pt, arc=2mm,
    before upper={\setlength{\parskip}{0.5em}}
}

% ----------------------------------------------------------------
% CUADRO CORNELL (tabla real: columna izquierda de preguntas,
% columna derecha de respuestas, fila final con el resumen)
% ----------------------------------------------------------------
\newlength{\cornizq}
\newlength{\cornancho}
\AtBeginDocument{
    \setlength{\cornizq}{0.28\textwidth}
    \setlength{\cornancho}{\dimexpr\textwidth-2\tabcolsep-13pt\relax}
}
\newcounter{cornellrow}
\newenvironment{cornell}{%
    \begingroup\small\setcounter{cornellrow}{0}%
    \renewcommand{\arraystretch}{1.3}%
    \begin{longtable}{>{\raggedright\arraybackslash}p{\cornizq}!{\vrule width 0.4pt}>{\raggedright\arraybackslash}p{\dimexpr\textwidth-\cornizq-4\tabcolsep-12pt\relax}}
    \toprule
    \textbf{Preguntas / palabras clave} & \textbf{Notas / respuestas} \\
    \midrule
    \endfirsthead
    \toprule
    \textbf{Preguntas / palabras clave} & \textbf{Notas / respuestas} \\
    \midrule
    \endhead
    \midrule
    \multicolumn{2}{r}{\footnotesize\itshape continúa en la página siguiente} \\
    \endfoot
    \bottomrule
    \endlastfoot
}{%
    \end{longtable}\endgroup
}
\newcommand{\cpregunta}[2]{%
    \ifnum\value{cornellrow}>0 \midrule[\lightrulewidth]\fi
    \stepcounter{cornellrow}%
    \textbf{#1} & #2 \\[0.3em]
}
\newcommand{\cresumen}[1]{%
    \midrule
    \multicolumn{2}{p{\cornancho}}{\vspace{0.2em}\textbf{Resumen del capítulo:} #1\vspace{0.2em}} \\
}

\hypersetup{breaklinks=true}

% ================================================================
\begin{document}
% ================================================================

\frontmatter

% ----------------------------------------------------------------
% PORTADA
% ----------------------------------------------------------------
\begin{titlepage}
\thispagestyle{empty}
\begin{center}

\vspace*{1cm}
{\large \textsc{\universidad}}

\vspace{0.2cm}
{\large \textsc{\facultad} --- \carrera}

\vspace{3cm}
{\Huge \textbf{\titulo}}

\vspace{1cm}
{\LARGE \subtitulo}

\vspace{0.8cm}
\textit{Estudiar con constancia convierte la comprensión en conocimiento.}

\vspace{1.2cm}
\rule{0.70\textwidth}{0.5pt}

\vspace{0.5cm}
{\large Apuntes de estudio integrados}

\vspace{0.5cm}
\rule{0.70\textwidth}{0.5pt}

\vfill
{\large \autor}

\vspace{0.3cm}
{\large \anio}

\end{center}

\vfill
\begin{flushleft}
\textit{Este material integra, reorganiza y depura en clave editorial diez apuntes de clase correspondientes a la materia ``\materia''. Es un aporte de estudio elaborado por un estudiante y no pretende sustituir el material oficial de la cátedra.}
\end{flushleft}
\end{titlepage}

% ----------------------------------------------------------------
% ÍNDICE
% ----------------------------------------------------------------
\tableofcontents

% ----------------------------------------------------------------
% INTRODUCCIÓN GENERAL
% ----------------------------------------------------------------
\chapter*{Introducción general}
\addcontentsline{toc}{chapter}{Introducción general}
\markboth{Introducción general}{Introducción general}

\section*{Qué es este libro}
\addcontentsline{toc}{section}{Qué es este libro}

Este volumen reúne, en un único cuerpo de estudio, el desarrollo completo de la materia \materia{}: desde sus fundamentos teóricos hasta el funcionamiento concreto de los dos procesos de integración regional que sirven de referencia constante a lo largo del curso, el Mercosur y la Unión Europea. No es una sucesión de clases: es un recorrido pensado para leerse de principio a fin, en el que cada capítulo se apoya en lo explicado antes y prepara lo que viene después.

La materia se organiza tradicionalmente en dos grandes partes. La primera estudia las \textbf{relaciones públicas de integración}: qué tipo de integración crea cada bloque, qué estructura institucional adopta, qué facultades tienen sus órganos, cómo se toman las decisiones y a quiénes obligan. Resulta más cercana a quienes ya se sintieron atraídos por el Derecho Internacional Público. La segunda parte estudia las \textbf{relaciones privadas dentro de un área de integración}: lo que ocurre con los particulares --personas y empresas-- cuando circulan mercaderías, servicios, capitales y trabajadores dentro del bloque. Este libro sigue, en términos generales, esa misma lógica, ampliada con los desarrollos institucionales y jurisprudenciales que la cursada fue incorporando.

\section*{Cómo está organizado}
\addcontentsline{toc}{section}{Cómo está organizado}

El libro se divide en nueve partes. Las Partes I a III sientan las bases conceptuales: qué es la integración, qué la distingue del derecho internacional común, qué grados de profundidad puede alcanzar y cómo se desarrolló históricamente en América Latina. Las Partes IV y V estudian en profundidad los dos casos de referencia --Mercosur y Unión Europea-- en su origen, su estructura institucional y su funcionamiento político. La Parte VI compara los sistemas de solución de controversias de ambos bloques, con especial atención al Tribunal de Justicia de la Unión Europea. La Parte VII desarrolla la inmunidad del Estado y de los órganos de integración frente a la jurisdicción de terceros. La Parte VIII cierra el andamiaje teórico con la jerarquía de fuentes propia del derecho de la integración y su relación con el derecho interno de cada Estado. La Parte IX sintetiza el recorrido completo y ofrece pautas de estudio para el examen final.

Cada capítulo se cierra, cuando la densidad conceptual lo justifica, con un cuadro Cornell: una tabla de preguntas y respuestas seguida de un resumen breve, pensada para la recuperación activa de lo estudiado. No se generó un cuadro para cada subtema menor, sino uno por unidad temática completa, de modo que sirva efectivamente como herramienta de repaso y no como una repetición mecánica del texto.

\section*{Advertencia sobre las fuentes}
\addcontentsline{toc}{section}{Advertencia sobre las fuentes}

Todo el contenido académico de este libro proviene exclusivamente de los apuntes de clase originales. Cuando una fecha, una cifra o una referencia normativa resultaba ambigua o inconsistente entre los distintos apuntes, se conservó tal como figuraba en la fuente y se señaló con una nota de verificación (\texttt{\% TODO}) en el código fuente, en lugar de corregirla con información externa. El criterio seguido en toda la obra es el mismo: fidelidad a lo efectivamente explicado en clase, por encima de cualquier otra consideración editorial.

\mainmatter
% ================================================================
\part{Fundamentos del Derecho de la Integración}
% ================================================================

% ----------------------------------------------------------------
\chapter{El Derecho de la Integración como rama autónoma}
\chaptermark{Autonomía del Derecho de la Integración}
% ----------------------------------------------------------------

Todo estudio del Derecho de la Integración debe partir de una pregunta previa: ¿de qué rama del derecho se trata? La respuesta no es evidente, porque la materia se ubica en un territorio intermedio entre dos ramas clásicas del derecho internacional, sin coincidir plenamente con ninguna de ellas.

\section{Distinción con el Derecho Internacional Público y Privado}

\begin{definicion}{Derecho Internacional Público}
Regula las relaciones entre \textbf{Estados soberanos}, actuando en su calidad de tales; es decir, ejerciendo actos de soberanía (\emph{imperium}).
\end{definicion}

\begin{definicion}{Derecho Internacional Privado}
Se ocupa de \textbf{relaciones particulares} que se desarrollan en un ámbito internacional. El Estado puede estar involucrado, pero actuando como persona privada; por ejemplo, al celebrar un contrato de compraventa internacional de mercadería.
\end{definicion}

\begin{ejemplo}{La compra de vacunas durante la pandemia}
Cuando durante la pandemia un Estado compró vacunas a otro país, se trataba de un contrato de compraventa internacional de mercadería con un Estado extranjero. En esa relación no estaba en juego la soberanía, ni el mar territorial, ni la plataforma submarina, ni el espacio aéreo: el Estado actuaba como una parte contractual más, en su calidad de persona privada.
\end{ejemplo}

Ambos tipos de relaciones jurídicas se desarrollan en un ámbito territorial ampliado, pero se diferencian en \textbf{la calidad en que actúa el Estado}: como soberano en un caso, y como persona privada en el otro. A estas dos ramas clásicas se agrega una tercera: el Derecho de la Integración.

\section{Un tercer ámbito: ni derecho público ni derecho privado puro}

El Derecho de la Integración regula un tipo de relaciones jurídicas que no pertenecen ni al Derecho Internacional Público ni al Derecho Internacional Privado. Existen reglas de relación entre los Estados, pero distintas de las que rigen entre Estados soberanos con actos de imperio; y también existen relaciones que afectan a las personas privadas, pero de naturaleza distinta a las del Derecho Internacional Privado clásico.

Se trata de un \textbf{híbrido}: no es derecho privado ni derecho público, porque se está integrando un área en la que se acuerdan con los socios determinadas reglas de compromiso o de comportamiento en distintos ámbitos --el cultural, el económico, el político-- según el tipo de integración que se establezca.

\section{El debate sobre la autonomía de la materia}

La ubicación del Derecho de la Integración dentro del plan de estudios fue objeto de un debate académico real. Existió un proyecto de reforma del plan de estudios de la Facultad que proponía denominar a las materias «Derecho Internacional Público y de la Integración» y «Derecho Internacional Privado y de la Integración». Ello implicaba tratar las relaciones de integración entre Estados junto con las relaciones internacionales públicas, y las relaciones privadas de integración junto con el derecho internacional privado. Siguiendo el modelo de la Unión Europea, otras universidades todavía mantienen esa fusión en sus programas, bajo denominaciones como «Derecho Internacional Privado y Comunitario».

\begin{importante}
La Facultad de Derecho tomó la postura de que se trata de una \textbf{materia distinta}: un derecho autónomo, que tiene relaciones entre Estados y relaciones con particulares, pero que no son de carácter internacional, sino \textbf{de integración}.
\end{importante}

\section{División de la materia en dos grandes partes}

A partir de esta autonomía, el programa de la materia --y, en consecuencia, este libro-- se organiza en dos grandes partes.

\begin{table}[H]
\centering
\small
\renewcommand{\arraystretch}{1.3}
\begin{tabularx}{\textwidth}{>{\raggedright\arraybackslash}p{0.27\textwidth}>{\raggedright\arraybackslash}X}
\toprule
\textbf{Parte} & \textbf{Contenido} \\
\midrule
\textbf{Relaciones públicas de integración} & Vinculadas con las relaciones entre Estados. Analizan qué tipo de integración se creó (política, económica, cultural), qué estructura institucional y qué órganos se establecieron, qué facultades tienen esos órganos, cómo están integrados y a quiénes obligan. Es la más afín al Derecho Internacional Público. \\
\addlinespace
\textbf{Relaciones privadas de integración} & Se ocupan de lo que ocurre con los particulares dentro de un área de integración: circulación de mercaderías, servicios, capitales y personas. Es la más afín al Derecho Internacional Privado. \\
\bottomrule
\end{tabularx}
\end{table}

\begin{ejemplo}{Un trabajador dentro del Mercosur}
Un trabajador argentino que va a trabajar a Uruguay dentro del Mercosur: qué ley se aplica a sus derechos laborales, dónde se jubila y con qué ley. Si aportó años de trabajo en dos países distintos del bloque, ¿con qué ley se jubilaría? Este interrogante --que la materia desarrolla en detalle al estudiar el mercado común-- ilustra por qué las relaciones privadas de integración constituyen un objeto de estudio propio.
\end{ejemplo}

\section{Cuadro Cornell}

\begin{cornell}
\cpregunta{¿Qué diferencia al Derecho Internacional Público del Privado?}{El Público regula relaciones entre Estados soberanos actuando como tales (actos de imperio). El Privado regula relaciones particulares en un ámbito internacional, incluso cuando interviene el Estado, pero actuando como persona privada.}
\cpregunta{¿Por qué el Derecho de la Integración es una materia autónoma?}{Porque regula relaciones que no son ni de Derecho Público ni de Derecho Privado clásico: reglas propias entre Estados socios y reglas propias para los particulares dentro de un área de integración. Es un híbrido.}
\cpregunta{¿Qué postura tomó la Facultad frente a los proyectos que fusionaban la materia con el derecho internacional?}{Sostuvo que se trata de una materia distinta y autónoma, frente a denominaciones como «Derecho Internacional Público y de la Integración» o «Derecho Internacional Privado y Comunitario».}
\cpregunta{¿En qué dos partes se divide el programa?}{Relaciones públicas de integración (entre Estados) y relaciones privadas de integración (particulares dentro del área).}
\cresumen{El Derecho Internacional Público regula relaciones entre Estados soberanos que actúan como tales, y el Privado, relaciones particulares en un ámbito internacional. El Derecho de la Integración constituye un tercer ámbito, híbrido, con reglas propias entre Estados y para particulares, y la Facultad lo consideró una materia autónoma. Su programa se divide en relaciones públicas de integración y relaciones privadas de integración, división que también organiza este libro.}
\end{cornell}

% ----------------------------------------------------------------
\chapter{Concepto de integración y principios propios}
\chaptermark{Concepto y principios}
% ----------------------------------------------------------------

\section{Por qué no existe una definición única de integración}

Los conceptos de integración que aparecen en los libros no están «mal», pero muchas veces no representan a un espacio de integración concreto. Una afirmación habitual sostiene que la integración implica que los Estados ceden soberanía en órganos supranacionales; sin embargo, no todos los espacios de integración delegan soberanía a un organismo internacional. Existen distintos tipos de integración, y cada órgano recibe facultades limitadas o amplias según lo que los Estados le hayan cedido específicamente, lo cual depende del tipo de integración de que se trate (política, económica, cultural, de defensa, etc.).

\begin{importante}
No es posible construir un concepto general y cerrado de integración: la definición dependerá siempre del ámbito de integración concreto que se tenga enfrente.
\end{importante}

\begin{ejemplo}{El Mercosur y la adopción}
Existen tesis doctorales sobre la adopción en el Mercosur. Frente a la postura de que sería conveniente regularla, cabe responder que no es conveniente porque excede las facultades del órgano de integración: el Mercosur es un mercado común del sur y no tiene competencia para regular minoridad ni filiación. Mientras no se modifiquen los objetivos del área de integración en el tratado fundacional, el órgano carece de facultades para regular esas materias. Este límite --desarrollado en el Capítulo~\ref{cap:atribucion} bajo el nombre de principio de atribución de facultades-- es una consecuencia directa de la idea de que no existe una integración «total» y automática.
\end{ejemplo}

\section{Cooperación versus integración}

Se distingue entre dos tipos de relaciones entre Estados que suelen confundirse: la cooperación y la integración.

\begin{definicion}{Cooperación}
Es «yo te ayudo»: una colaboración voluntaria, sin que exista una obligación estructural entre las partes ni un objetivo común permanente.
\end{definicion}

\begin{ejemplo}{El matrimonio y las tareas domésticas}
Es lo que ocurre cuando uno de los integrantes de un matrimonio afirma «yo te ayudé, puse el lavarropas». No se trata de ayuda: es su obligación (quien ensució, lava el baño). Cooperar es esto de «yo te ayudo».
\end{ejemplo}

La relación de integración no es solamente cooperación: hay un \textbf{objetivo común} y un \textbf{compromiso}, no simplemente ayuda mutua ocasional. Dentro del ámbito de integración pueden además existir distintos niveles de compromiso: no todos los Estados que se organizan en un área de integración se obligan al mismo nivel de responsabilidad, y pueden coexistir distintos modelos de integración según cuántos compromisos se asuman.

\section{La metáfora del matrimonio}
\label{sec:matrimonio}

Para explicar el concepto de integración se recurre a una analogía extendida entre la integración y el matrimonio, contrastándola con un simple acuerdo o negocio transitorio entre dos partes.

\begin{importante}
La integración es un matrimonio, no una unión de hecho, porque implica \textbf{obligaciones y compromiso recíprocos}. No se trata de dos partes que se unieron transitoriamente: hay un proyecto común, y cada uno puede establecerse con requisitos propios y facultades propias, pero todo para un objetivo común y no para un objetivo propio.
\end{importante}

Esta idea se ilustra con una distribución de roles dentro de un matrimonio en la que uno de los integrantes trabaja fuera de casa mientras el otro se encarga del cuidado de los hijos: aunque esa distribución pueda generar comentarios de terceros, responde a una organización funcional orientada a un objetivo común. A partir de esta base se desarrolla el contraste central de la metáfora:

\begin{itemize}
    \item \textbf{En un negocio}, se participa por el beneficio que reporta; no se mira si a la contraparte le va bien o mal, porque es problema de ella.
    \item \textbf{En una integración}, existen otro tipo de relaciones: cada Estado debe tener conciencia de que si hace algo que afecte a otro Estado del área, tarde o temprano le repercute a él.
\end{itemize}

\begin{ejemplo}{La economía doméstica y el corte de luz}
Supóngase que ambos integrantes de una pareja trabajan y comparten las tareas, pero uno gasta todo su sueldo en compras personales y al otro le queda pagar los impuestos. Si no se pagó la luz porque el dinero se gastó en otra cosa, se la cortan a toda la casa, y no solamente al integrante despilfarrador. Trasladada esta lógica al mercado común, una medida económica que en lo inmediato beneficia a un socio del bloque puede a largo plazo perjudicar a la unión completa, restringiendo la economía y las relaciones del conjunto --idea que reaparece en la Parte~IV al analizar las tensiones entre política interna y compromisos de integración en el Mercosur.
\end{ejemplo}

\section{Principio de reciprocidad presumida}

La reciprocidad funciona de manera distinta en el Derecho Internacional clásico y dentro de un área de integración.

\begin{table}[H]
\centering
\small
\renewcommand{\arraystretch}{1.3}
\begin{tabularx}{\textwidth}{>{\raggedright\arraybackslash}p{0.27\textwidth}>{\raggedright\arraybackslash}X}
\toprule
\textbf{Ámbito} & \textbf{Funcionamiento de la reciprocidad} \\
\midrule
Internacional & Se reconoce un derecho «en tanto y en cuanto» la ley del otro Estado reconozca a los ciudadanos propios ese mismo derecho. Es una \textbf{reciprocidad condicionada}, que debe verificarse caso por caso. \\
\addlinespace
De integración & Se da por \textbf{sobreentendida}. Si un Estado celebró un acuerdo de integración con su socio, es porque ese socio va a aplicar exactamente la misma norma que se está aplicando: no es necesario someter el derecho a la condición de reciprocidad, porque esta rige desde el momento mismo en que se decidió asociarse. \\
\bottomrule
\end{tabularx}
\end{table}

Esta idea se relaciona con el funcionamiento de la norma de conflicto propia del Derecho Internacional Privado, que implica un «salto al vacío», porque no se sabe a dónde se irá a parar. La diferencia en el ámbito de integración es que se conoce quiénes son los socios: cuando se elaboran las normas de conflicto y se establece que se aplicará la ley del lugar de ejecución, y que la ejecución será, por ejemplo, en Brasil, ello es porque ya se sabe qué dice la norma brasileña.

\section{Igualdad de trato, no discriminación y preferencia nacional}

Algunos principios que parecen lógicos en el ámbito del Derecho Internacional Público, como la preferencia de los nacionales frente a los extranjeros en determinadas situaciones, resultan \textbf{inadmisibles dentro de un área de integración}. De la lógica de integración descrita en la Sección~\ref{sec:matrimonio} se deriva una consecuencia normativa concreta: en un mercado único, un Estado miembro \textbf{no puede aplicar preferencias nacionales frente a sus socios} del bloque, aunque sí podría hacerlo frente a terceros Estados ajenos a la integración.

\begin{ejemplo}{La licitación de obra pública}
Si hay una obra pública y se presentan una empresa argentina y una brasileña, no puede decirse «prefiero la argentina por compra nacional», porque se estaría violando la igualdad de trato con el socio. En cambio, si se presenta una empresa chilena (y Chile no integra el bloque), sí puede decirse: a igualdad de precio, se prefiere la argentina, porque es compra nacional.
\end{ejemplo}

\begin{definicion}{Arraigo procesal}
Exigencia procesal por la cual un extranjero que pretende litigar en un país debe fijar una garantía --por ejemplo, un domicilio, un bien o un seguro-- para responder eventualmente por las costas y honorarios del juicio en caso de resultar perdidoso y no contar con bienes en el país.
\end{definicion}

En el ámbito del Derecho Internacional clásico se presume que el extranjero debe fijar arraigo, salvo que un convenio bilateral específico lo exceptúe expresamente.

\begin{importante}
En el ámbito de integración es al revés: rige el \textbf{principio de igualdad de trato}. No se puede exigir arraigo a un extranjero proveniente de un Estado socio, ni limitar su acceso a la justicia; debe tratárselo en las mismas condiciones que a un nacional. Por la misma razón, así como un nacional sin bienes puede solicitar el beneficio de litigar sin gastos, un extranjero proveniente de un país socio del bloque podría, en principio, solicitar ese mismo beneficio.
\end{importante}

\section{Cuadro Cornell}

\begin{cornell}
\cpregunta{¿Por qué no existe una definición única de integración?}{Porque depende del tipo de integración (política, económica, cultural, de defensa) y de las facultades específicas que cada bloque haya cedido a sus órganos.}
\cpregunta{¿Qué diferencia a la cooperación de la integración?}{La cooperación es ayuda voluntaria, sin obligación estructural ni objetivo común permanente. La integración implica un compromiso recíproco orientado a un objetivo común.}
\cpregunta{¿Por qué la integración se parece más a un matrimonio que a un negocio?}{Porque implica obligaciones y compromiso recíprocos y un proyecto común, y no una unión transitoria orientada solo al beneficio propio: lo que afecta a un Estado del área tarde o temprano repercute en los demás.}
\cpregunta{¿Por qué la reciprocidad se presume en el ámbito de integración?}{Porque al asociarse con un socio del bloque ya se sabe que ese socio aplicará exactamente la misma norma; no hace falta condicionar el reconocimiento de un derecho, como sí ocurre en el ámbito internacional clásico (reciprocidad condicionada).}
\cpregunta{¿Qué implica la igualdad de trato respecto de la preferencia nacional y el arraigo procesal?}{Un Estado miembro no puede aplicar preferencia nacional frente a sus socios del bloque, aunque sí frente a terceros. Del mismo modo, no puede exigirse arraigo procesal a un extranjero socio del bloque: debe recibir el mismo trato que un nacional.}
\cresumen{No existe un concepto único de integración: depende del tipo de integración y de las facultades cedidas a los órganos. Se distingue de la cooperación por el compromiso recíproco y el objetivo común, comparable a un matrimonio y no a un negocio transitorio. Sus principios propios incluyen la reciprocidad presumida y la igualdad de trato, que impide la preferencia nacional frente a socios del bloque y el arraigo procesal.}
\end{cornell}

% ----------------------------------------------------------------
\chapter{La técnica de los tratados marco: objetivos y medios}
\chaptermark{Tratados marco: objetivos y medios}
% ----------------------------------------------------------------

\section{La crítica a la técnica legislativa}

Una advertencia metodológica atraviesa todo el estudio de la materia: los tratados deben leerse con \textbf{sentido crítico}, sin conformarse con el título ni con la enunciación general del texto. Con frecuencia los legisladores redactan cláusulas ambiguas, que admiten interpretaciones contradictorias; ese déficit de técnica legislativa obliga al intérprete a «leer entre líneas» para establecer qué quiso decir realmente la norma. Un texto puede sonar comprometido y, sin embargo, no contener ninguna propuesta concreta: una situación semejante puede presentarse en los tratados de integración cuando las declaraciones de buenas intenciones no se traducen en obligaciones exigibles.

\section{Reconocer, garantizar y promover: la clave de lectura}

\begin{importante}
No es lo mismo que un Estado \textbf{reconozca} un derecho, que lo \textbf{garantice} o que lo \textbf{promueva}. Estas tres categorías implican compromisos jurídicos de intensidad completamente distinta, y buena parte del trabajo de interpretación de un tratado consiste en detectar cuál de las tres está efectivamente en juego.
\end{importante}

Para ilustrar la distinción con un ejemplo del derecho constitucional argentino, puede compararse el tratamiento que la Constitución Nacional da a la libertad de culto en general y al culto católico en particular. El artículo 2 establece que «el Gobierno federal sostiene el culto católico apostólico romano», mientras que el artículo 14 reconoce a todos los habitantes el derecho de «profesar libremente su culto». Ambas cláusulas tienen igual jerarquía constitucional, pero generan obligaciones estatales de naturaleza distinta: el artículo 14 \textbf{garantiza} la libertad general de culto, mientras que el artículo 2 impone al Estado federal una obligación de \textbf{sostenimiento económico} específica respecto del culto católico. Garantizar la libertad de culto no es lo mismo que sostener un culto determinado, y ese tipo de distinciones es el que debe aprenderse a identificar al analizar un tratado de integración. La misma lógica separa a una \textbf{declaración} de derechos humanos, que expresa un reconocimiento sin generar necesariamente una obligación de garantía exigible, de una \textbf{convención}, que implica un compromiso jurídico concreto de protección.

\section{Integración formal e integración informal}

\begin{definicion}{Integración formal}
Existe cuando el Estado es parte de un tratado, es decir, cuando el compromiso se formaliza jurídicamente y se le da una estructura institucional.
\end{definicion}

\begin{definicion}{Integración informal}
Existe cuando, de hecho, un conjunto de Estados actúa como bloque --por ejemplo, votando de manera conjunta frente a organismos internacionales o sosteniendo posiciones y propuestas de defensa comunes-- sin que exista un tratado que los obligue jurídicamente a conducirse de ese modo.
\end{definicion}

\section{El tratado marco: la «Constitución» del área de integración}

\begin{definicion}{Tratado marco}
Dentro de las integraciones formales, el tratado marco es el instrumento que cumple, dentro del área de integración, una función equivalente a la que cumple una Constitución dentro de un Estado: establece cuál es el proyecto común, qué aspiran los Estados parte, qué buscan y cuáles son sus fines y objetivos.
\end{definicion}

\section{Objetivos y medios: la relación instrumental}

\begin{definicion}{Objetivo (o fin) y medio}
El \textbf{objetivo} (o fin) es lo que se persigue alcanzar. El \textbf{medio} es el instrumento del que se vale el Estado para alcanzar ese objetivo.
\end{definicion}

\begin{ejemplo}{Llegar a la facultad}
El objetivo es llegar a la facultad; el medio es el trayecto elegido. Si el estudiante toma el subte y el servicio se interrumpe, puede pedir un vehículo de aplicación y de todos modos cumplir el objetivo de llegar, aunque haya cambiado el medio. Esta lógica es trasladable al área de integración: si el objetivo está claro, el medio puede modificarse en tanto siga permitiendo alcanzar ese mismo objetivo.
\end{ejemplo}

\begin{ejemplo}{Objetivo y medio en una convención de derechos humanos}
Al hablar de una convención ya no se trata de meras declaraciones: es un compromiso, que consiste en el reconocimiento y la protección de los derechos humanos; ese es el \textbf{fin}. Respecto del \textbf{medio}, un Estado podría sancionar una ley que castigue penalmente los atentados contra la vida de otra persona, o crear un cuerpo policial especializado en su protección. El objetivo --proteger el derecho a la vida-- permanece estable, mientras que los medios elegidos pueden variar de un Estado a otro.
\end{ejemplo}

\textbf{Criterio práctico de lectura.} Los objetivos de un tratado suelen encontrarse en sus palabras introductorias o \textbf{preámbulo}, en fórmulas del tipo «nosotros, convencidos de...» o «conscientes de la necesidad de...». Antes de leer la totalidad del articulado, conviene prestar atención a esas cláusulas introductorias, porque allí suele anticiparse el rumbo general del instrumento.

\begin{importante}
El medio no puede absorber ni desplazar al objetivo: si el medio se lleva puesto el objetivo, cambió el área de integración, porque el objetivo pasa a ser otro. Se está cambiando de área de integración cuando se cambia el objetivo perseguido, no simplemente cuando se cambia el medio empleado para conseguirlo.
\end{importante}

\section{Tipos de objetivos: político, económico, cultural y de seguridad}

Un área de integración puede perseguir distintos tipos de objetivos: políticos, culturales, económicos, o de seguridad y defensa común. Si bien es habitual asociar la palabra «integración» únicamente con la integración económica, existen áreas cuyo objetivo principal no es económico. El Mercosur (Mercado Común del Sur) es un ejemplo de objetivo económico; la Unión de Naciones Suramericanas (UNASUR), en cambio, fue un organismo de integración con objetivos predominantemente políticos --la región llegó a contar simultáneamente con ambos organismos, lo que también puede explicarse por disputas de liderazgo político regional--. Este comportamiento se contrasta con el proceso europeo, que partió de un objetivo inicial acotado (el carbón, el acero y la energía atómica, según se desarrolla en el Capítulo~\ref{cap:origenes-ue}) y fue sumando progresivamente nuevas dimensiones.

\begin{ejemplo}{La defensa común frente a un enemigo compartido}
En las áreas de integración cuyo objetivo principal es la defensa común, los Estados que se sienten amenazados por un enemigo compartido se asocian primordialmente para garantizar su protección, con independencia de que compartan o no la misma religión, el mismo régimen de gobierno o los mismos criterios sobre otras cuestiones internas.
\end{ejemplo}

\section{Todo acuerdo de integración es, en el fondo, un acuerdo de paz}

Según la posición de un sector de la doctrina, todo acuerdo de integración lleva implícito, aunque no se explicite, un \textbf{acuerdo de paz}: cuando dos o más Estados comparten un proyecto común, en principio no debería haber una controversia de fondo entre ellos respecto de ese proyecto, del mismo modo en que dos personas que deciden formar una familia acuerdan de antemano sus fines compartidos, aun cuando después deban ajustar los medios concretos para alcanzarlos.

\section{La integración como proceso progresivo}

La integración es un \textbf{proceso} que avanza desde un compromiso menor hacia compromisos progresivamente mayores, hasta alcanzar, en su forma más completa, una integración que abarca todos los aspectos de la vida y de las relaciones entre los Estados. Frente a la objeción de que un bloque «no hace integración» por carecer de ciertos elementos presentes en otros modelos, lo determinante no es la ausencia de determinados elementos institucionales sino el \textbf{cambio del objetivo perseguido}: mientras el objetivo se mantenga, los medios pueden variar o ampliarse progresivamente sin que ello implique abandonar el proyecto original. Por último, el objetivo de un área de integración debe ser un \textbf{objetivo común} que beneficie a todos los Estados parte por igual; ese carácter común es lo que legitima, frente a cada Estado, la cesión de determinadas facultades a favor del proyecto compartido.

\section{Cuadro Cornell}

\begin{cornell}
\cpregunta{¿Qué diferencia a una integración formal de una informal?}{La formal existe cuando el Estado es parte de un tratado, con estructura institucional. La informal existe cuando un conjunto de Estados actúa de hecho como bloque, sin tratado que los obligue.}
\cpregunta{¿Qué es el tratado marco y dónde suelen estar sus objetivos?}{Es el instrumento que cumple la función de una Constitución dentro del área de integración. Sus objetivos suelen encontrarse en el preámbulo, antes que en el articulado.}
\cpregunta{¿Por qué deben leerse los tratados con sentido crítico?}{Porque los legisladores redactan a veces cláusulas ambiguas que obligan a «leer entre líneas»; un texto puede sonar comprometido sin contener ninguna propuesta concreta.}
\cpregunta{¿Qué diferencia hay entre reconocer, garantizar y promover un derecho?}{Son tres categorías de compromiso jurídico de intensidad distinta (ejemplo: arts.\ 2 y 14 de la Constitución Nacional respecto del culto), del mismo modo en que una declaración de derechos no equivale a una convención.}
\cpregunta{¿Cuál es la relación entre objetivo y medio?}{El medio puede modificarse en tanto siga permitiendo alcanzar el objetivo; pero el medio no puede desplazar al objetivo, porque entonces cambió el área de integración.}
\cpregunta{¿Qué tipos de objetivos puede perseguir un área de integración?}{Políticos, culturales, económicos y de seguridad o defensa común; no toda integración es económica (Mercosur frente a UNASUR).}
\cresumen{Un tratado de integración debe leerse críticamente, distinguiendo si el Estado reconoce, garantiza o promueve un derecho. La integración es formal cuando existe un tratado marco --la «Constitución» del área-- e informal cuando los Estados actúan como bloque sin obligación jurídica. Lo central es la relación entre objetivo (político, económico, cultural o de defensa) y medio: el medio puede variar mientras se preserve el objetivo, pero no puede desplazarlo, y la integración es un proceso progresivo con un objetivo común para todos los Estados parte.}
\end{cornell}

% ----------------------------------------------------------------
\chapter{Modelos de integración y estructura institucional}
\chaptermark{Modelos e instituciones}
\label{cap:atribucion}
% ----------------------------------------------------------------

\section{La integración como proceso evolutivo: la metáfora de la escalera}

Una forma habitual de caracterizar a las áreas de integración es como un proceso progresivo, comparable a una \textbf{escalera}: un bloque inicia en el escalón más elemental y va perfeccionándose y profundizándose con el tiempo hasta alcanzar una integración más completa. El caso europeo es el ejemplo paradigmático de este modelo evolutivo: partió de una comunidad limitada al carbón y al acero, y fue incorporando progresivamente energía atómica, luego la economía en su conjunto, después aspectos políticos y más adelante materias tan alejadas del objeto original como el medio ambiente o el trato a la ancianidad (véase el Capítulo~\ref{cap:origenes-ue}).

\begin{ejemplo}{El límite político del proceso evolutivo}
En algún momento la Unión Europea intentó avanzar hacia una Constitución propia. Al someterse la propuesta a referéndum en los Estados miembros, la población la rechazó, por entender que se perdía identidad nacional (este episodio se retoma con más detalle en el Capítulo~\ref{cap:evolucion-ue}). El proceso evolutivo de profundización, entonces, no es ilimitado ni automático: encuentra un techo en la voluntad política y popular de cada Estado.
\end{ejemplo}

\section{Modelos estáticos: México, Estados Unidos y Canadá}

Frente al modelo evolutivo europeo se presenta un contraejemplo deliberado: el bloque conformado por México, Estados Unidos y Canadá, organizado en torno a una zona de libre comercio con un objetivo deliberadamente acotado y estable en el tiempo, sin vocación de ir sumando facultades. Si se aplica la noción europea de integración a este modelo, algunos autores podrían objetar que «eso tampoco es integración», precisamente por carecer de la proyección de avance progresivo propia del modelo europeo. Sin embargo, esto \textbf{no invalida al modelo}: se trata de un tipo de integración distinto, con objetivos definidos y estables.

\section{Principio de atribución (o especialidad) de facultades}

\begin{definicion}{Principio de atribución (o especialidad) de facultades}
Los órganos de integración solo tienen las facultades que los Estados miembros les otorgaron específicamente, y no pueden expandirlas por sí mismos.
\end{definicion}

\begin{ejemplo}{Las provincias y el Estado Nacional argentino}
Las provincias precedieron al Estado Nacional; por lo tanto, el Estado Nacional no puede legislar sobre aquellas cuestiones que no fueron específicamente cedidas por las provincias a la Nación. La misma idea vale para los órganos de integración: la provincia de Corrientes no puede sancionar un Código Civil, porque esa facultad fue transferida de manera exclusiva a la Nación.
\end{ejemplo}

Aplicando esta lógica al Mercosur, puede regular sobre matrimonio internacional dentro del área, pero no sobre minoridad, a pesar de que existen tesis doctorales sobre la adopción en el Mercosur: no son facultades del órgano de integración. La misma idea se extiende al tráfico de menores: aunque desde el punto de vista de la política legislativa pudiera parecer deseable una regulación conjunta, ese ámbito excede el objeto del Mercosur --que carece de competencia penal-- y, por lo tanto, no corresponde otorgarle esa facultad sin antes modificar el tratado fundacional.

\begin{nota}
Se hizo referencia, en este sentido, a un protocolo del Mercosur sobre medidas cautelares en cuestiones penales y administrativas. Si bien se trata de un instrumento valioso, corresponde en rigor al ámbito del Derecho Internacional --un tratado entre los Estados miembros-- y no estrictamente al Derecho de la Integración, por exceder el objeto específico del mercado común.
\end{nota}

\begin{importante}
El órgano de integración no puede ir más allá de las facultades que le fueron otorgadas en el tratado fundacional. Las facultades no se amplían solas: sería necesario cambiar el tratado fundamental, fijando un objetivo distinto. Este límite funciona como una \textbf{garantía} frente al traspaso de facultades de los Estados al órgano comunitario sin el debido respaldo del tratado fundacional.
\end{importante}

\section{Supranacionalidad versus intergubernamentalidad}
\label{sec:supra}

Existen dos maneras estructuralmente distintas de organizar los órganos de un bloque de integración, según el grado de independencia que tengan respecto de los gobiernos que los designan.

\begin{definicion}{Órgano supranacional}
Una vez designado el funcionario, este no responde al Estado que lo propuso, sino a la comunidad de integración en su conjunto. Las decisiones se toman por \textbf{mayoría} y resultan obligatorias para todos los Estados miembros, incluso para aquellos que votaron en contra.
\end{definicion}

\begin{definicion}{Órgano intergubernamental}
Está integrado por \textbf{representantes de los gobiernos}, que siguen las instrucciones de sus respectivos Estados. Las decisiones no se imponen por mayoría simple, sino que requieren un mecanismo más cercano al consenso entre todos los socios.
\end{definicion}

La elección entre uno u otro modelo institucional --qué tipo de órgano se crea y qué facultades se le otorgan-- depende del \textbf{tipo de relaciones que existan entre los socios} y del grado de confianza mutua y de profundidad al que aspire el bloque. Cuando no existe una confianza plena entre los socios, resulta más difícil otorgar a un órgano común facultades supranacionales que puedan imponerse por mayoría incluso en contra de la propia voluntad de un Estado. El Capítulo~\ref{cap:origenes-ue} desarrolla el caso de la CECA como ejemplo histórico de estructura supranacional, en contraste con el modelo intergubernamental que adoptó el Mercosur (Parte~IV).

\section{Cuadro Cornell}

\begin{cornell}
\cpregunta{¿Por qué algunos autores describen la integración como una escalera?}{Porque la ven como un proceso evolutivo: el bloque parte del escalón más elemental y se profundiza con el tiempo (caso europeo), encontrando un techo en la voluntad política y popular.}
\cpregunta{¿Todo modelo de integración aspira a profundizarse con el tiempo?}{No. El bloque de México, Estados Unidos y Canadá es una zona de libre comercio con objetivo acotado y estable, sin vocación de sumar materias: un tipo de integración distinto, no inválido.}
\cpregunta{¿Qué es el principio de atribución (o especialidad) de facultades?}{Los órganos de integración solo pueden ejercer las facultades específicamente cedidas por los Estados en el tratado fundacional, y no pueden ampliarlas por sí mismos.}
\cpregunta{¿Puede el Mercosur regular adopción o tráfico de menores?}{No: excede su objeto fundacional (mercado común) y carece de competencia penal. Regularlo requeriría modificar el tratado fundacional.}
\cpregunta{¿Qué diferencia a un órgano supranacional de uno intergubernamental?}{El supranacional actúa con independencia de los Estados que lo designaron, decide por mayoría y sus decisiones obligan a todos. El intergubernamental representa a los gobiernos y decide mediante un mecanismo cercano al consenso.}
\cresumen{La integración puede concebirse como un proceso evolutivo (modelo europeo) o como un modelo estático con objetivo acotado (México, Estados Unidos y Canadá). En todos los casos rige el principio de atribución de facultades: los órganos solo tienen las facultades que el tratado fundacional les otorgó. Los órganos pueden ser supranacionales (mayoría, independencia de los Estados) o intergubernamentales (representantes de gobiernos, consenso), según las relaciones entre los socios.}
\end{cornell}

% ----------------------------------------------------------------
\chapter{Características del proceso de integración}
\chaptermark{Características del proceso}
\label{cap:caracteristicas}
% ----------------------------------------------------------------

La integración puede analizarse a partir de un objetivo dominante --político, económico, cultural o de defensa--, pero nunca se agota exclusivamente en ese aspecto. Ponerse de acuerdo en materia económica presupone siempre, en algún grado, un acuerdo político previo, porque implica crear órganos y otorgarles facultades que son, en principio, soberanas del Estado.

\begin{importante}
Ninguna integración es exclusivamente cultural, exclusivamente política ni exclusivamente económica: para ponerse de acuerdo en lo económico tuvo que haber un acuerdo político.
\end{importante}

\section{Carácter pluridimensional y plurifacético}

\begin{definicion}{Integración pluridimensional y plurifacética}
La integración es \textbf{pluridimensional} porque no se limita a una única dimensión de las relaciones entre los Estados: impacta simultáneamente sobre distintos planos, como el político, el económico, el legislativo y el cultural. Es \textbf{plurifacética} porque afecta distintas fases y dimensiones de las relaciones entre los Estados y entre los particulares.
\end{definicion}

Dentro del plano legislativo, al decidir crear normas comunes, los Estados deben definir, en primer lugar, una política acerca de cómo se van a dictar: por mayoría simple, por mayoría calificada, por consenso o por unanimidad. Esa decisión, aunque pertenece al plano normativo, es en sí misma una decisión política que condiciona el funcionamiento de todo el sistema.

\subsection{Uniformidad y armonización legislativa}

\begin{definicion}{Uniformar y armonizar}
\textbf{Uniformar} implica dictar una única norma, con una única solución aplicable por igual a todos los Estados parte. \textbf{Armonizar} implica acordar principios generales comunes, dejando librada a cada Estado la regulación específica de los detalles de implementación.
\end{definicion}

\begin{ejemplo}{Inversiones extranjeras: un caso de armonización}
Varios Estados de la región acuerdan un régimen de inversiones extranjeras en el que se establece, de manera común, qué tipo de inversiones se aceptarán --por ejemplo, podría excluirse de manera conjunta a las inversiones mineras contaminantes y priorizarse las vinculadas con tecnología o energía sustentable--. Superado ese acuerdo de principios generales, cada Estado regula libremente los requisitos concretos de inscripción y los beneficios aplicables. El resultado es que una misma empresa puede recibir un trato distinto según decida invertir en un país o en otro dentro del mismo bloque, porque solo se armonizaron los criterios generales y no la solución de fondo.
\end{ejemplo}

Esta distinción es relevante también para entender por qué distintos autores discuten, muchas veces, si un determinado bloque «realmente hace integración»: la respuesta depende, en parte, de qué modelo de integración se utilice como parámetro de comparación, ya que uniformidad y armonización son dos caminos jurídicos válidos, aunque de distinta profundidad. El caso concreto de los derechos del consumidor en el Mercosur, desarrollado en el Capítulo~\ref{cap:mercado-comun}, ilustra por qué la armonización no siempre resulta alcanzable.

\section{Interrelación e interdependencia}

\begin{definicion}{Interrelación e interdependencia}
Hay \textbf{interrelación} cuando la conducta de uno de los Estados provoca una reacción en los demás. Hay \textbf{interdependencia} cuando todos los Estados dependen, en algún grado, de las decisiones de los otros.
\end{definicion}

\begin{ejemplo}{Preferencias a la industria nacional y su efecto en cadena}
Un Estado sanciona una ley de preferencia a la industria nacional. Esa decisión afecta a los productos de otro Estado miembro del bloque, lo que a su vez repercute en el funcionamiento del mercado común en su conjunto, y ese efecto retorna, de un modo u otro, sobre el propio Estado que dictó la medida, en tanto forma parte del mismo mercado.
\end{ejemplo}

\section{Estructura supranacional y delegación de soberanía}

Las definiciones doctrinarias de integración suelen incluir, como característica, la existencia de una \textbf{estructura orgánica supranacional} y la \textbf{delegación de soberanía} estatal a favor de esa estructura. Sin embargo, en la mayoría de las áreas de integración esa estructura no existe; sí existe, con matices, en la Unión Europea (véase la Parte~V).

Técnicamente, el Estado nunca cede su soberanía en sentido pleno: \textbf{delega el ejercicio exclusivo de determinadas facultades} a un órgano de integración, absteniéndose de regular por sí mismo en esa área específica --la misma lógica de facultades no delegadas que sostiene, en el orden interno argentino, la relación entre las provincias y el Estado Nacional (Capítulo~\ref{cap:atribucion})--.

\begin{importante}
Si los Estados delegaron determinadas facultades a un órgano y este se pone a regular sobre cosas que no le fueron transferidas, excede su ámbito de competencia. Del mismo modo en que, dentro de un Estado, un decreto reglamentario no puede modificar la ley que reglamenta, un órgano de integración no puede legislar sobre materias que los Estados no le delegaron.
\end{importante}

\section{Cuadro Cornell}

\begin{cornell}
\cpregunta{¿Por qué ninguna integración es puramente económica, política o cultural?}{Porque ponerse de acuerdo en lo económico presupone un acuerdo político previo: implica crear órganos y otorgarles facultades que son, en principio, soberanas del Estado.}
\cpregunta{¿Qué significa que la integración es pluridimensional y plurifacética?}{Que no se limita a una única dimensión de las relaciones entre Estados: impacta sobre los planos político, económico, legislativo y cultural, y afecta distintas fases de las relaciones entre Estados y particulares.}
\cpregunta{¿Qué diferencia hay entre uniformar y armonizar una legislación?}{Uniformar es dictar una única norma con una única solución para todos los Estados parte. Armonizar es acordar principios generales comunes, dejando que cada Estado regule los detalles de implementación.}
\cpregunta{¿Qué significa que la integración se caracteriza por interrelación e interdependencia?}{Interrelación: la conducta de un Estado provoca una reacción en los demás. Interdependencia: todos los Estados dependen, en algún grado, de las decisiones de los otros.}
\cpregunta{¿Qué implica la delegación de soberanía en una estructura supranacional?}{El Estado no cede su soberanía en sentido pleno: delega el ejercicio exclusivo de determinadas facultades a un órgano, absteniéndose de regular por sí mismo en esa materia. Lo no delegado permanece como competencia exclusiva del Estado.}
\cresumen{Toda integración es pluridimensional, plurifacética, interrelacionada e interdependiente, cualquiera sea su objetivo dominante. En el plano normativo, uniformar y armonizar son caminos válidos de distinta profundidad. Solo algunos modelos, como el europeo, presentan una estructura supranacional; en ella el Estado delega el ejercicio exclusivo de ciertas facultades, y lo no delegado sigue siendo suyo.}
\end{cornell}

% ----------------------------------------------------------------
\chapter{Factores que condicionan la integración}
\chaptermark{Factores condicionantes}
% ----------------------------------------------------------------

Más allá de la voluntad política declarada, existen aspectos que pueden facilitar o dificultar un proceso de integración entre Estados: factores culturales y religiosos; ideológicos y políticos; territoriales y climáticos; y factores vinculados con la estabilidad democrática y la continuidad de la política exterior.

\section{Factor cultural y religioso}

Las diferencias de religión, de costumbres y de organización social influyen de manera directa sobre la viabilidad de un proyecto de integración, aunque con distinta intensidad según el tipo de integración de que se trate: la religión puede no ser determinante para una integración estrictamente económica, pero sí puede serlo para una integración política o cultural.

\subsection{El caso latinoamericano: una base común subestimada}

A diferencia de lo que ocurre entre bloques culturalmente heterogéneos, los países latinoamericanos comparten una base religiosa, cultural y jurídica común, heredada del proceso de conquista y colonización española --las leyes de Indias y las Siete Partidas de Alfonso el Sabio estuvieron regulando durante siglos--, de modo que lo que es un contrato, una sociedad o un matrimonio para un país es también un contrato, una sociedad o un matrimonio para los demás. Esa base debería facilitar, y no dificultar, los procesos de integración regional. Se reconoce, dentro de esta explicación, la preexistencia de comunidades y pueblos originarios con organizaciones sociales propias, que conviven dentro del ordenamiento jurídico común --la Constitución y el Código Civil-- en un marco de reconocimiento de su preexistencia étnica y cultural. Si la base religiosa, social y jurídica de la región es relativamente homogénea, la dificultad para integrarse debe buscarse, entonces, en el plano político.

\section{Factor ideológico y político}

Los cambios de orientación ideológica de los gobiernos de turno inciden directamente sobre la continuidad de los compromisos de integración: cada nueva gestión puede modificar sus prioridades respecto de con quiénes conviene asociarse y en qué términos. Si un gobierno prefiere un mercado libre y considera que las preferencias con determinados socios no convienen a su Estado, no propiciará el avance de un área de integración en la que deba otorgar forzosamente privilegios a sus socios.

\begin{nota}
\textbf{La falta de un proyecto país de largo plazo.} Esta inestabilidad ideológica se vincula con una característica recurrente de la región: la ausencia de un proyecto de país sostenido en el tiempo, más allá del signo político del gobierno de turno. Con frecuencia, cada nueva gestión tiende a desestimar lo construido por la anterior y a comenzar de nuevo, en lugar de sostener acuerdos básicos de largo plazo. Este mismo problema, con ejemplos concretos de la historia económica argentina, se desarrolla en profundidad en el Capítulo~\ref{cap:tensiones}.
\end{nota}

\subsection{La noción de nación según la teoría de Mancini}

Se retoma aquí un concepto trabajado en Derecho Internacional Privado: la teoría de la nacionalidad de Pascuale Stanislao Mancini, jurista italiano del siglo~XIX. Según esta teoría, no constituye una nación un grupo de personas que están en un territorio hablando un mismo idioma y con las mismas costumbres: la \textbf{nación} implica tener conciencia de un pasado en común y un objetivo propio en común. Trasladada al plano de la integración regional, esta noción sugiere que lo que permite que un conjunto de Estados conforme efectivamente un proyecto compartido no es solamente compartir un idioma o un origen histórico común, sino tomar conciencia de esa tradición compartida y fijarse, a partir de ella, un objetivo común hacia el futuro --elemento que, según este planteo, suele faltar en los procesos de integración latinoamericanos.

\section{Factor territorial y climático}

La extensión territorial y las condiciones climáticas de los Estados involucrados condicionan de manera directa la viabilidad y el costo de implementación de un área de integración, particularmente en lo relativo a infraestructura y transporte.

\begin{ejemplo}{El costo de la infraestructura regional}
En Europa las distancias entre países son comparativamente pequeñas. Un área de integración como el Mercosur, en cambio, debe conectar territorios de Argentina, Brasil, Paraguay y Uruguay, atravesando selvas, distancias considerablemente mayores y condiciones climáticas dispares. Para lograr el mismo objetivo de libre tránsito de mercaderías, un área de integración latinoamericana requiere entonces una inversión en infraestructura sustancialmente mayor que un área de integración europea, con independencia de las diferencias en capacidad económica entre ambas regiones.
\end{ejemplo}

Por eso una zona de libre comercio en África no es lo mismo que una en Europa ni que una en Latinoamérica: la falta de avance de un proceso de integración nunca responde a una única causa aislada, sino a un conjunto de condiciones políticas, culturales y geográficas que operan de manera combinada. Como ejemplo de cómo la decisión política incide sobre el diseño institucional, cabe adelantar aquí el contraste --desarrollado en profundidad en la Parte~IV-- entre el sistema de decisión de la Unión Europea, basado en distintos tipos de mayorías, y el del Mercosur, que exige \textbf{consenso} entre todos los Estados parte.

\section{Estabilidad democrática y política exterior}

Un último factor es la estabilidad institucional y democrática de cada Estado, y su capacidad de sostener una política exterior coherente a lo largo del tiempo, con independencia del signo político del gobierno de turno.

\begin{ejemplo}{La autonomía del banco central como política de Estado}
En el sistema estadounidense, la presidencia del banco central es ejercida por un funcionario de carrera que no es removido por el cambio de gobierno, lo cual otorga estabilidad a la política monetaria con independencia de qué administración se encuentre en el poder. Este esquema se contrapone con la práctica argentina, en la que el cambio de gobierno suele traer aparejado el reemplazo del titular del Banco Central y del Ministerio de Economía, lo que dificulta sostener una política monetaria de largo plazo.
\end{ejemplo}

Como ejemplos adicionales de continuidad de la política exterior más allá del signo de gobierno se mencionan la guerra de Vietnam, que se extendió durante distintas administraciones estadounidenses de distinto signo partidario (retomada en el Capítulo~\ref{cap:tensiones}), y el caso del entonces presidente Barack Obama, quien, luego de recibir el Premio Nobel de la Paz en 2009, autorizó en 2011 la operación militar que culminó con la muerte de Osama bin Laden.

\begin{importante}
La falta de continuidad institucional y de políticas de Estado sostenidas en el tiempo constituye uno de los principales obstáculos para consolidar proyectos de integración duraderos en la región, aunque las condiciones culturales, religiosas y jurídicas de base sean, en principio, relativamente favorables.
\end{importante}

\subsection{El límite de la integración frente al odio interétnico}

Para ilustrar el extremo opuesto --contextos en los que las diferencias entre comunidades resultan insalvables para cualquier proyecto de integración-- puede citarse el genocidio ocurrido en Ruanda en 1994, conflicto interétnico de extrema violencia entre las comunidades hutu y tutsi en el que murieron cientos de miles de personas en un período muy breve, sin que las fuerzas internacionales desplegadas lograran evitar la masacre. El punto central es que, en contextos de odio interétnico profundamente arraigado, ningún modelo de integración económica, social o cultural resulta viable, porque el rechazo hacia el otro impide incluso las relaciones comerciales o laborales más básicas.

\section{Cuadro Cornell}

\begin{cornell}
\cpregunta{¿Qué factores culturales y religiosos condicionan un proceso de integración?}{Las diferencias de religión, costumbres y organización social influyen sobre la viabilidad de la integración, con intensidad variable según el tipo de integración de que se trate.}
\cpregunta{¿Qué base común comparten los países latinoamericanos?}{Una base religiosa, cultural y jurídica heredada de la conquista y colonización española, que debería facilitar la integración; la dificultad se explica en el plano político.}
\cpregunta{¿Por qué los cambios de ideología afectan los compromisos de integración?}{Porque cada gobierno puede modificar con quiénes conviene asociarse y en qué términos, a lo que se suma la ausencia de un proyecto de país de largo plazo.}
\cpregunta{¿Qué noción de nación desarrolla Mancini y cómo se aplica a la integración?}{La nación implica conciencia de un pasado común y un objetivo propio en común, no solo compartir territorio e idioma; según este planteo, eso suele faltar en América Latina.}
\cpregunta{¿Qué factores territoriales y climáticos condicionan la integración?}{La extensión territorial y las condiciones climáticas inciden en el costo de la infraestructura necesaria para el libre tránsito de mercaderías (Mercosur frente a Europa).}
\cpregunta{¿Qué límite impone el odio interétnico a la integración?}{Cuando existe un odio interétnico profundamente arraigado (ejemplo: Ruanda, 1994), ningún modelo de integración resulta viable.}
\cresumen{La viabilidad de una integración depende, además de la voluntad política declarada, de factores culturales y religiosos, ideológicos y políticos, territoriales y climáticos, y de la estabilidad democrática y la continuidad de la política exterior. América Latina comparte una base cultural, religiosa y jurídica relativamente homogénea, pero la falta de continuidad de las políticas de Estado constituye uno de los principales obstáculos para la integración.}
\end{cornell}

% ================================================================
\part{Grados de integración económica}
% ================================================================

Así como un bloque de integración puede describirse según su objetivo dominante y sus características generales, la dimensión específicamente \textbf{económica} de la integración admite una escala propia de niveles, cada uno de mayor profundidad y complejidad institucional que el anterior. Esta parte recorre esa escala --de la zona franca a la unión económica-- y analiza el caso del Mercosur como ejemplo práctico permanente.

% ----------------------------------------------------------------
\chapter{Zona franca, unión fronteriza y zona de libre comercio}
\chaptermark{Primeros escalones de integración económica}
\label{cap:zlc}
% ----------------------------------------------------------------

\section{La zona franca: una decisión unilateral, no un acuerdo de integración}

Existe un error conceptual frecuente: confundir una zona de libre comercio con una zona franca.

\begin{definicion}{Zona franca}
Espacio dentro del territorio de un único Estado soberano, creado por decisión unilateral de ese Estado, con el propósito de promover el desarrollo económico de una determinada región de su propio territorio --mediante beneficios fiscales, incentivos salariales o exenciones de aranceles de importación--, sin que su existencia dependa del acuerdo o de la voluntad de otros Estados.
\end{definicion}

\begin{ejemplo}{Beneficios unilaterales para poblar una zona desfavorable}
Dos mecanismos típicos de una zona franca son el pago de un adicional salarial para incentivar que trabajadores se radiquen en una zona geográfica considerada desfavorable, y la exención de aranceles de importación para determinados bienes, como ocurre con ciertos regímenes de compra de vehículos en zonas francas. En ambos casos se trata de decisiones que el Estado adopta de manera unilateral para promover una parte específica de su propio territorio.
\end{ejemplo}

\begin{importante}
La zona franca no es un acuerdo de integración: es un espacio dentro del territorio que el Estado soberano quiere promover, sin importar lo que piensen los demás Estados. No compromete a ningún otro Estado.
\end{importante}

\section{La unión fronteriza}

\begin{definicion}{Unión fronteriza}
Primer nivel de integración económica propiamente dicha. Establece políticas específicas de tránsito de personas, bienes o servicios limitadas exclusivamente a las áreas de frontera entre los Estados involucrados --por ejemplo, determinadas provincias o regiones limítrofes--, sin extenderse a la totalidad del territorio de cada Estado.
\end{definicion}

La delimitación concreta del área de frontera alcanzada por este tipo de acuerdo depende de la organización interna de cada Estado --puede tratarse de una provincia entera, como Formosa y Misiones en el caso argentino, o de una región más acotada-- y en general responde a una decisión política consensuada con las provincias o regiones involucradas, ya que estas deben estar de acuerdo con la apertura del tránsito en su territorio.

\begin{ejemplo}{Unión fronteriza entre Chile y Argentina}
Existe un tratado entre Chile y Argentina orientado a facilitar la integración regional de las provincias fronterizas. Una unión fronteriza no equivale a establecer una zona de libre comercio entre las capitales de ambos países, sino que responde a una necesidad de regular una realidad de hecho: la circulación cotidiana de personas y mercaderías que ya ocurre entre las poblaciones fronterizas. La regulación de la triple frontera responde a esta misma lógica.
\end{ejemplo}

\section{Zona de libre comercio}

\begin{definicion}{Zona de libre comercio}
Espacio de integración económica que comprende la \textbf{totalidad del territorio} de cada uno de los Estados que la integran --y no solamente sus zonas limítrofes--, cuyo objetivo es el libre comercio de mercaderías: que el arancel a su circulación entre los Estados miembros llegue a ser \textbf{cero}, mediante reducciones progresivas.
\end{definicion}

Ese resultado no se alcanza de manera inmediata a partir de la firma del tratado, sino a través de una \textbf{escala de desgravación arancelaria}. En el Mercosur, esta reducción no puede hacerse de un día para el otro, porque afecta la balanza presupuestaria de cada Estado; por eso se requiere un organismo que se reúna, negocie, evalúe los efectos y actualice periódicamente la tabla de aranceles --que en el Boletín Oficial se publica aproximadamente cada tres meses--.

\begin{importante}
El objetivo de la zona de libre comercio es que el arancel a la circulación de mercaderías entre los Estados miembros llegue a ser \textbf{cero}. Por eso se habla de \textbf{levantamiento de las barreras arancelarias}: se busca eliminar todo obstáculo al intercambio, es decir, el proteccionismo.
\end{importante}

La desgravación arancelaria no puede aplicarse de manera uniforme entre todos los Estados ni sobre todos los productos por igual, porque la sensibilidad económica de cada rubro varía según el peso que tiene en el presupuesto de cada país: una misma rebaja puede representar una pérdida presupuestaria considerable para un Estado y ser insignificante para otro. Por eso los órganos técnicos de un bloque de libre comercio suelen establecer tablas de desgravación diferenciadas, ajustadas según la repercusión real en cada economía. Ningún Estado puede eliminar de un día para el otro la totalidad de sus ingresos aduaneros, porque el Estado se sostiene en gran medida gracias a esos impuestos --incluida la educación y la salud pública--; por eso el proceso es necesariamente gradual y está sujeto a permanente revisión.

\subsection{El problema del origen de los productos: «industria nacional»}

Si un Estado no cobra aranceles a los demás miembros del área, surge la siguiente cuestión: ¿qué ocurre si se compran bienes a un tercer país a un arancel más conveniente y luego, como los demás miembros no pueden cobrar aranceles entre sí, esos bienes circulan a costo cero dentro del área? La respuesta radica en dos elementos: existe un \textbf{arancel externo} que se cobra a los terceros países (desarrollado en el Capítulo~\ref{cap:union-aduanera}), y para que se dé la libre circulación los bienes deben ser \textbf{industria nacional}, es decir, \textbf{productos originarios del área}.

\begin{importante}
Para circular libremente en el área de integración, los bienes deben ser \textbf{originarios del área}. Sin embargo, casi ningún producto es cien por ciento de fabricación nacional. Por eso la normativa establece \textbf{qué porcentaje del valor final} del producto puede corresponder a componentes extranjeros sin que pierda la calificación de «industria nacional». Esto se determina \textbf{producto por producto}, y el criterio no es la cantidad de piezas sino el \textbf{valor} que aporta el componente extranjero.
\end{importante}

\begin{ejemplo}{La vaca y el automóvil}
Una vaca en pie que nació en la Argentina y se alimentó de pasto argentino es cien por ciento argentina; pero si el empaque proviene de otro Estado o si fue vacunada con insumos importados, cabe cuestionar si sigue siendo cien por ciento argentina. Del mismo modo, el motor de un automóvil puede provenir de Italia mientras que los neumáticos, los frenos y el chasis se fabrican en el país: lo relevante es qué porcentaje del \textbf{valor final} del auto representa el componente extranjero, no la cantidad de piezas importadas.
\end{ejemplo}

\subsection{Caso práctico: el conflicto de las bicicletas entre Argentina y Uruguay}

\begin{ejemplo}{Conflicto de las bicicletas (Argentina y Uruguay)}
Argentina tenía una gran industria nacional de elaboración de bicicletas. Uruguay comenzó a importar piezas separadas desde China y otros países y las ensamblaba localmente. El valor de las piezas sumadas representaba una fracción reducida del valor final de la bicicleta armada, pintada y funcionando, de modo que el valor agregado por Uruguay era muy superior al valor de la importación; según la norma, no se estaría vulnerando entonces el principio de industria nacional. Argentina argumentaba que Uruguay traía todas las piezas del exterior y solo las armaba; Uruguay se defendía sosteniendo que la norma fija el porcentaje sobre el valor agregado al producto, no sobre el valor de las piezas. El resultado fue que Uruguay vendía bicicletas a Argentina sin aranceles y más baratas que las de producción nacional, lo que afectó a la industria argentina del sector. Se inició un proceso de \textbf{arbitraje} ante el Mercosur, y el caso se resolvió finalmente mediante acuerdo entre las partes.
\end{ejemplo}

Estas cuestiones, que pueden parecer menores, se vinculan directamente con el libre tránsito de mercaderías: las condiciones en que se produce ese tránsito importan, porque el levantamiento de las barreras puede hacer quebrar la economía en un determinado rubro.

\section{Cuadro Cornell}

\begin{cornell}
\cpregunta{¿Por qué una zona franca no constituye integración?}{Porque es una decisión unilateral de un Estado que otorga beneficios fiscales para instalar personas e industrias en un área de su propio territorio, sin comprometer a ningún otro Estado.}
\cpregunta{¿Qué caracteriza a la unión fronteriza?}{Es el primer nivel de integración económica. Establece políticas de circulación de personas, bienes o servicios limitadas a las áreas de frontera entre los Estados, no a la totalidad de sus territorios.}
\cpregunta{¿Qué es la zona de libre comercio y cómo se alcanza el arancel cero?}{Espacio que comprende la totalidad del territorio de los Estados parte y aspira al arancel intrazona cero, mediante una escala de desgravación arancelaria progresiva y diferenciada según el impacto en cada economía.}
\cpregunta{¿Qué es un producto originario del área?}{El que cumple con el porcentaje de valor nacional exigido por la normativa del bloque, medido por el valor final del producto y no por la cantidad de piezas, y determinado producto por producto.}
\cpregunta{¿Qué fue el conflicto de las bicicletas?}{Uruguay importaba piezas, las ensamblaba y vendía bicicletas a Argentina sin aranceles, amparándose en que el valor agregado superaba el umbral legal. El caso se resolvió mediante arbitraje y acuerdo.}
\cresumen{La escala de integración económica comienza en la zona franca --decisión unilateral que no compromete a otros Estados--, sigue con la unión fronteriza --limitada a las áreas de frontera-- y continúa con la zona de libre comercio, que abarca todo el territorio de los Estados parte y aspira al arancel intrazona cero mediante desgravación progresiva. La libre circulación exige además que los bienes sean originarios del área, criterio medido por valor y no por cantidad de piezas, como ilustra el conflicto de las bicicletas entre Argentina y Uruguay.}
\end{cornell}

% ----------------------------------------------------------------
\chapter{Unión aduanera}
\chaptermark{Unión aduanera}
\label{cap:union-aduanera}
% ----------------------------------------------------------------

Una aduana se ocupa de las importaciones y exportaciones. En la zona de libre comercio el análisis se centraba en el ámbito intrazona; en la unión aduanera se incluyen además las importaciones y exportaciones con \textbf{terceros países}, y se establecen los aranceles que se cobrarán.

\begin{definicion}{Unión aduanera}
Nivel que supera la zona de libre comercio al establecer un \textbf{Arancel Externo Común (AEC)}: todos los Estados miembros cobran la misma tasa arancelaria a los productos provenientes de terceros países. Si todos cobran lo mismo al exterior y el arancel intrazona es cero, la disparidad competitiva entre los Estados miembros ya no estará en los aranceles sino en los costos productivos internos (salarios, impuestos nacionales, eficiencia productiva).
\end{definicion}

\section{Necesidad de un órgano específico}

En la unión aduanera se requiere una estructura institucional más compleja. Existen dos posturas entre los autores: algunos sostienen que en este nivel se necesita un \textbf{órgano concreto, específico y totalmente independiente} --no una reunión de representantes de distintos países, sino una aduana del bloque que no dependa de las directrices de ningún Estado en particular--; otros consideran que basta con una \textbf{oficina de aduana} dentro de cada aduana nacional, sin necesidad de un edificio propio, siempre que cuente con una estructura que funcione con independencia de la aduana nacional.

\begin{table}[H]
\centering
\renewcommand{\arraystretch}{1.3}
\begin{tabularx}{\textwidth}{>{\raggedright\arraybackslash}p{0.24\textwidth}>{\raggedright\arraybackslash}X>{\raggedright\arraybackslash}X}
\toprule
\textbf{Característica} & \textbf{Zona de libre comercio} & \textbf{Unión aduanera} \\
\midrule
Arancel intrazona & Cero (objetivo) & Cero (objetivo) \\
Arancel a terceros países & Cada Estado lo fija libremente & AEC: todos cobran lo mismo \\
Estructura institucional & Órgano negociador común & Órgano más independiente (aduana) \\
\bottomrule
\end{tabularx}
\caption{Zona de libre comercio y unión aduanera.}
\end{table}

\section{Evolución histórica y estado actual del Mercosur}

En la época de la ALALC (véase el Capítulo~\ref{cap:aladi}), el tratamiento del comportamiento frente a terceros países se dejaba para después de lograr el arancel cero intrazona; la ALADI, en cambio, partió de la necesidad de avanzar en ambos frentes a la vez, porque no es posible estabilizar el mercado interno y la economía si no se sale de las crisis económicas.

\begin{importante}
El Mercosur es definido como una \textbf{zona de libre comercio incompleta} y una \textbf{unión aduanera incompleta}: no se llegó al arancel cero intrazona en la totalidad de los productos, ni se logró establecer el AEC para la totalidad de los bienes. Faltan por cubrir áreas y aspectos en ambos niveles.
\end{importante}

\section{Cuadro Cornell}

\begin{cornell}
\cpregunta{¿Qué es el Arancel Externo Común?}{La tasa arancelaria única que todos los Estados miembros de una unión aduanera cobran a los productos provenientes de terceros países.}
\cpregunta{¿En qué se diferencian la zona de libre comercio y la unión aduanera?}{En ambas el objetivo intrazona es el arancel cero. En la zona de libre comercio cada Estado fija libremente el arancel a terceros países; en la unión aduanera todos cobran el mismo (AEC), y se requiere además un órgano aduanero más independiente.}
\cpregunta{¿Por qué el Mercosur es «incompleto»?}{Porque no se alcanzó el arancel intrazona cero en todos los productos ni el AEC para la totalidad de los bienes.}
\cresumen{La unión aduanera agrega a la zona de libre comercio el Arancel Externo Común frente a terceros países y exige una estructura institucional más independiente. El Mercosur es, a la vez, una zona de libre comercio y una unión aduanera incompletas.}
\end{cornell}

% ----------------------------------------------------------------
\chapter{Mercado común}
\chaptermark{Mercado común}
\label{cap:mercado-comun}
% ----------------------------------------------------------------

Al ascender en la escala se llega al mercado común. Ya no se trata solamente de la parte comercial: se habla de \textbf{economía de mercado}.

\begin{definicion}{Mercado común}
Aspira a la \textbf{libre circulación de los cuatro factores de producción económica}: (1) mercaderías o bienes; (2) personas, en su faz laboral (el factor trabajo); (3) servicios, incluidos los profesionales y los turísticos; y (4) capitales, tanto de los socios como de terceros países.
\end{definicion}

\section{Libre circulación de mercaderías}

Un mercado común aspira a que las mercaderías circulen sin adicionarles impuestos, aranceles, retenciones ni otras cargas. Es \textbf{mercadería} aquello que puede comprarse o venderse, es decir, lo que es susceptible de apreciación económica --por ejemplo, un monumento público no sería objeto de una norma de libre circulación, porque no está en el comercio.

\section{Libre circulación de personas: el factor trabajo}

La persona interesa al mercado común en tanto factor de producción, en su faz laboral. No se regula, por ejemplo, el matrimonio, porque no incide en la producción de un mercado; lo que se regula no es si la persona es capaz o incapaz, sino si tiene \textbf{aptitud o capacidad legal para trabajar}.

\begin{nota}
El \textbf{Protocolo de Tucumán} regula específicamente el libre tránsito laboral en el Mercosur. A diferencia de la Unión Europea, que habla del libre tránsito de personas en sentido amplio, en el Mercosur el enfoque es exclusivamente laboral.
\end{nota}

Para que el libre tránsito laboral sea efectivo deben regularse, entre otras, las siguientes cuestiones: la norma de cada país no puede \textbf{discriminar} entre trabajadores según su nacionalidad de origen dentro del bloque; no basta con establecer la igualdad de trato, sino que es necesario definir qué se considera \textbf{salario} y qué tipo de actividad se considera \textbf{trabajo} --el protocolo se refiere al trabajo en relación de dependencia--; y es necesario \textbf{calificar}, en el sentido propio del Derecho Internacional, qué se entiende por trabajador y qué aspectos de esa persona pueden regularse dentro del espacio de integración. Para facilitar el libre tránsito puede establecerse además un documento que garantice la identidad en igualdad de condiciones --el Documento Mercosur--, que habilita el tránsito entre Estados sin necesidad de pasaporte, en tanto la persona es ciudadana mercosureña.

\section{Libre circulación de servicios}

Dentro de los servicios se encuentran los servicios profesionales, lo que plantea todas las cuestiones vinculadas a la \textbf{reválida de títulos}: el reconocimiento de un título obtenido en una universidad de otro Estado, y si el profesional puede o no ejercer.

\begin{importante}
Hasta el momento, en el Mercosur la reválida y el reconocimiento de títulos profesionales tienen efecto \textbf{únicamente académico}. Un profesional de otro Estado miembro puede venir a dar un posgrado y cobrar honorarios, pero no ejercer libremente la profesión. La razón: en medicina, por ejemplo, el cuerpo humano es el mismo en cualquier país, pero la cuestión es la \textbf{capacitación} recibida, que varía entre universidades y entre contextos --la carrera de arquitectura en un país con actividad sísmica enfoca contenidos distintos que en un país sin esa realidad.
\end{importante}

\begin{ejemplo}{El modelo de la Unión Europea para la reválida}
Cuando la Unión Europea abordó la libre circulación de servicios y la reválida de títulos, exigió que profesionales como abogados, arquitectos o ingenieros trabajaran durante determinada cantidad de años --aproximadamente seis-- bajo la firma y la protección de un profesional recibido en el Estado receptor, quien garantizaba haber controlado el trabajo. Transcurrido ese período, se reconocía el derecho a ejercer la profesión con efecto pleno, y no solo académico.
\end{ejemplo}

Para regular los servicios turísticos con libre tránsito de personas se necesita garantizar la igualdad de trato: el turismo de los propios nacionales solo traslada el dinero de un lugar a otro sin generar nuevos ingresos, mientras que para generar nuevos ingresos es necesario que el dinero provenga de otro país del bloque.

\begin{nota}
\textbf{Derechos del consumidor: un desafío pendiente.} Hasta el momento no se logró una normativa unificada de defensa del consumidor en el Mercosur. La idea era tomar la ley más desarrollada de cada Estado y elevar los estándares, pero Brasil --que cuenta con un Código de Defensa del Consumidor muy completo-- se niega a bajar sus estándares para admitir productos de Paraguay con otro nivel de exigencia, porque ello implicaría reducir los derechos de sus propios ciudadanos. Este caso ilustra, en la práctica, el límite que la armonización legislativa (Capítulo~\ref{cap:caracteristicas}) puede encontrar cuando los niveles de protección de partida son muy desiguales.
\end{nota}

\section{Libre circulación de capitales e inversiones}

El libre tránsito de capitales se relaciona con las inversiones, tanto de los socios como de terceros países. No podría crearse, por ejemplo, un impuesto que limitara la libre circulación de capitales entre socios del bloque --lo que sí puede exigirse es la declaración e imposición tributaria sobre ese capital--. Un ejemplo del tratamiento de los capitales dentro del bloque es el régimen uruguayo de sociedades: Uruguay permite constituir sociedades cuyo objeto no se desarrolla en su territorio (sociedades \emph{offshore}), especialmente reguladas y con beneficios propios.

\section{Cuadro Cornell}

\begin{cornell}
\cpregunta{¿Cuáles son los cuatro factores de un mercado común?}{Mercaderías, personas (en su faz laboral), servicios (incluidos profesionales y turísticos) y capitales.}
\cpregunta{¿Qué aspecto de la persona regula el mercado común?}{Su aptitud o capacidad legal para trabajar, no cuestiones como el matrimonio ni la capacidad o incapacidad en general (Protocolo de Tucumán).}
\cpregunta{¿Cuál es el estado actual de la reválida de títulos en el Mercosur?}{Tiene solo efecto académico; la Unión Europea, en cambio, exige años de práctica supervisada para reconocer el ejercicio pleno de la profesión.}
\cpregunta{¿Por qué no se logró una normativa unificada de consumo en el Mercosur?}{Porque Brasil, con un Código de Defensa del Consumidor muy avanzado, se niega a bajar sus estándares para admitir productos con menores exigencias de calidad.}
\cresumen{El mercado común aspira a la libre circulación de cuatro factores: mercaderías, personas en su faz laboral, servicios y capitales. En el Mercosur el libre tránsito de personas es exclusivamente laboral, la reválida de títulos tiene solo efecto académico, la libre circulación de capitales impide crear impuestos que la limiten, y la normativa unificada de consumo no se logró por la disparidad de estándares entre los Estados miembros.}
\end{cornell}

% ----------------------------------------------------------------
\chapter{Comunidad económica y unión económica}
\chaptermark{Comunidad y unión económica}
% ----------------------------------------------------------------

\section{Comunidad económica}

En la comunidad económica no se habla solamente del mercado y de los factores de producción, sino de la economía en general.

\begin{definicion}{Comunidad económica}
Añade al mercado común la integración económica plena, incluyendo moneda única con banco central supranacional, política financiera común (tasas de interés, reservas, emisión monetaria), y aspectos no solo comerciales sino también económicos y financieros de amplio alcance.
\end{definicion}

Una moneda única requiere que varios países usen la misma moneda, pero además exige definir quién le da el valor, cuántas monedas se emiten y quién decide. No se trata de reunir las reservas en un solo lugar: se exige a cada Estado que mantenga determinada reserva. Cada país conserva su propio banco central, que emite la moneda; pero un órgano supranacional determina cuánta moneda se emite, qué valor tendrá y qué tasa de interés puede cobrarse a nivel económico.

\begin{importante}
El \textbf{Banco Central Europeo} es un ejemplo de institución supranacional que establece los parámetros de emisión monetaria para cada Estado miembro, exige a cada Estado mantener determinado nivel de reservas, fija tasas de interés mínimas y máximas, y puede sancionar a un Estado que emita moneda de más, porque ello desvaloriza la moneda común a nivel general.
\end{importante}

En el debate académico se planteó alguna vez la propuesta de dolarizar el Mercosur. Se objeta que no es posible estabilizar, como moneda de un espacio de integración, una moneda cuyo valor no se maneja: en una moneda única, la moneda debe valer lo mismo para todos los Estados miembros, y si un Estado no puede manejar el valor de su propia moneda, menos aún puede pretender manejar el de una moneda ajena.

\section{Unión económica y la Unión Europea}

La denominación actual de este nivel es \textbf{unión económica}: ya no se trata de una economía en común, sino de una sola entidad.

\begin{nota}
En el Tratado de Lisboa, la Unión Europea regula aspectos que van más allá de los factores de producción: derecho al ambiente sano, protección a la ancianidad y derecho a morir dignamente. Ya no se considera a la persona solo desde una perspectiva productiva, sino desde la perspectiva de los derechos humanos.
\end{nota}

\begin{importante}
Existe una relación \textbf{directamente proporcional} entre el grado de profundidad del área de integración y su complejidad institucional: a mayor compromiso, mayor cantidad y complejidad de órganos, mayor complejidad del derecho aplicable y más funciones y facultades para esos órganos.
\end{importante}

Si todos los países debieran ponerse de acuerdo en cada decisión, no se lograría consenso; por eso algunas decisiones se toman por \textbf{mayoría}, mientras que otras, una vez establecido determinado marco, se resuelven por \textbf{unanimidad}, \textbf{consenso}, \textbf{mayorías especiales} o \textbf{doble mayoría} --cuestión que se retoma en detalle al estudiar los órganos de la Unión Europea en la Parte~V--. La Comunidad Andina también cuenta con órganos supranacionales: aunque no alcanza la complejidad de la Unión Europea, ya se encuentra dentro del ámbito de una comunidad económica.

\section{Cuadro comparativo general}

\begin{table}[H]
\centering
\small
\begin{tabularx}{\textwidth}{
    >{\raggedright\arraybackslash}p{0.17\textwidth}
    >{\raggedright\arraybackslash}p{0.13\textwidth}
    >{\raggedright\arraybackslash}p{0.07\textwidth}
    >{\raggedright\arraybackslash}X
    >{\raggedright\arraybackslash}X
}
\toprule
\textbf{Nivel} & \textbf{Arancel intrazona} & \textbf{AEC} & \textbf{Libre circulación de factores} & \textbf{Política económica común} \\
\midrule
Zona franca & No aplica & No & No & No \\
Unión fronteriza & Reducido & No & No & No \\
Zona de libre comercio & Cero (objetivo) & No & Mercaderías & No \\
Unión aduanera & Cero (objetivo) & Sí & Mercaderías & Arancel externo común \\
Mercado común & Cero & Sí & Mercaderías, personas, servicios, capitales & Parcial \\
Comunidad económica & Cero & Sí & Todos los factores & Moneda única, banco central \\
Unión económica & Cero & Sí & Todos los factores y derechos humanos & Política común plena \\
\bottomrule
\end{tabularx}
\caption{Niveles de integración económica.}
\end{table}

\section{Cuadro Cornell}

\begin{cornell}
\cpregunta{¿Qué diferencia al mercado común de la comunidad económica?}{La comunidad económica incorpora, además de los cuatro factores, una moneda única y una política financiera común, con instituciones supranacionales como el Banco Central Europeo.}
\cpregunta{¿Por qué no es viable dolarizar el Mercosur?}{Porque la moneda de un espacio de integración debe valer lo mismo para todos los miembros y su valor debe poder manejarse; si un Estado no maneja el valor de su propia moneda, no puede pretender manejar el de una moneda ajena.}
\cpregunta{¿Qué es la unión económica?}{Es la denominación actual del nivel de integración plena: una sola entidad, que incorpora también la perspectiva de los derechos humanos (Tratado de Lisboa).}
\cpregunta{¿Cuál es la relación entre integración y complejidad institucional?}{Es directamente proporcional: a mayor profundidad del área de integración, mayor cantidad y complejidad de órganos y de derecho aplicable.}
\cresumen{La comunidad económica agrega al mercado común una moneda única y una política financiera común. La unión económica, denominación actual del nivel de integración plena, incorpora además la perspectiva de los derechos humanos. A mayor profundidad de la integración, mayor complejidad institucional y jurídica, y las decisiones se adoptan según distintos sistemas de mayorías, consenso o unanimidad, según el tema.}
\end{cornell}

% ================================================================
\part{Antecedentes y desarrollo de la integración latinoamericana}
% ================================================================

% ----------------------------------------------------------------
\chapter{Etapas históricas de la integración latinoamericana}
\chaptermark{Etapas de la integración latinoamericana}
% ----------------------------------------------------------------

La integración en América Latina puede periodizarse en dos grandes etapas: una \textbf{anterior a la década de 1970} y otra \textbf{posterior a la década de 1970-80}. El criterio de corte responde a un quiebre político regional: durante la década de 1970 la mayoría de los países latinoamericanos atravesó gobiernos militares, lo que produjo una discontinuidad institucional profunda en toda la región.

\section{Integración antes de la década de 1970: primero el mercado interno}

En la primera etapa, las áreas de integración entendían que, al conformarse, lo primero que debía establecerse era el \textbf{mercado interno}: estabilizar primero las economías internas de cada Estado, «puertas adentro», y una vez consolidado ese mercado interno, avanzar hacia una política exterior común frente a terceros países. Durante este período predominó la creación de \textbf{zonas de libre comercio} (Capítulo~\ref{cap:zlc}) como primer escalón del proceso.

\begin{ejemplo}{La analogía de la escolaridad}
Esta lógica escalonada se compara con el sistema educativo formal: así como una persona debe completar el nivel primario antes de acceder al secundario, cada etapa no es un mero requisito formal, sino que supone la adquisición efectiva de una base que sostiene el paso al escalón siguiente. Del mismo modo, en la primera etapa de la integración latinoamericana se entendía que un mercado interno estabilizado era la base indispensable antes de coordinar una política exterior común.
\end{ejemplo}

\section{Integración después de la década de 1970-80: el retorno democrático}

Con el retorno de los países de la región a gobiernos democráticos, cambió también la lógica de los procesos de integración. En muchos casos, las economías internas se encontraban seriamente deterioradas tras los períodos de gobiernos de facto, lo que llevó a los Estados a encarar de manera simultánea la necesidad de fortalecer el mercado interno y la de conseguir inversión extranjera y una inserción internacional favorable: para hacer crecer la economía se necesitan recursos, porque no crece de la nada.

\begin{table}[H]
\centering
\renewcommand{\arraystretch}{1.3}
\begin{tabularx}{\textwidth}{>{\raggedright\arraybackslash}p{0.20\textwidth}>{\raggedright\arraybackslash}X>{\raggedright\arraybackslash}X}
\toprule
 & \textbf{Antes de la década de 1970} & \textbf{Después de la década de 1970-80} \\
\midrule
\textbf{Contexto} & Etapa previa al quiebre político de los gobiernos militares. & Retorno democrático; economías internas deterioradas tras los gobiernos de facto. \\
\textbf{Estrategia} & Primero estabilizar el mercado interno; luego, una política exterior común. & Encarar simultáneamente el fortalecimiento del mercado interno y la búsqueda de inversión extranjera. \\
\textbf{Modelo predominante} & Zonas de libre comercio como primer escalón. & Acuerdos bilaterales entre Argentina y Brasil como base del Mercosur, planteado directamente como mercado común. \\
\bottomrule
\end{tabularx}
\caption{Comparación de las dos etapas históricas de la integración latinoamericana.}
\end{table}

\section{Los antecedentes bilaterales entre Argentina y Brasil}

En este período se ubica el origen del proceso que luego derivaría en el Mercosur. El punto de partida fue un acercamiento bilateral entre \textbf{Argentina y Brasil}, los dos países de mayor extensión territorial y mayor población de América del Sur y, por lo tanto, los de mayor peso relativo como mercados, que comenzaron con acuerdos parciales y bilaterales. El antecedente directo fue la \textbf{Declaración de Foz de Iguazú}, suscripta el 30 de noviembre de 1985 entre los entonces presidentes de Argentina y Brasil, que dio inicio a un Programa de Integración y Cooperación Económica entre ambos países. A partir de ese acercamiento se sucedieron distintos acuerdos bilaterales de complementación económica, hasta desembocar en la firma del tratado multilateral que dio origen formal al Mercosur, desarrollado en el Capítulo~\ref{cap:asuncion}.

Un dato relevante para comprender la profundidad institucional del proceso es que el acuerdo finalmente alcanzado por los cuatro Estados fundadores no se limitó a establecer una zona de libre comercio --como había sido habitual en la primera etapa histórica--, sino que planteó directamente la constitución de un \textbf{mercado común}, un compromiso institucional considerablemente más profundo.

\section{Cuadro Cornell}

\begin{cornell}
\cpregunta{¿Cuáles son las dos grandes etapas de la integración latinoamericana?}{Una anterior a la década de 1970, centrada en estabilizar primero el mercado interno, y otra posterior al retorno democrático, en la que los Estados encararon simultáneamente el fortalecimiento interno y la apertura a la inversión extranjera.}
\cpregunta{¿Cuál fue el punto de partida del proceso que derivó en el Mercosur?}{Un acercamiento bilateral entre Argentina y Brasil, con antecedente directo en la Declaración de Foz de Iguazú (30 de noviembre de 1985).}
\cpregunta{¿En qué se diferenció el acuerdo de los cuatro Estados fundadores del modelo de la primera etapa?}{No se limitó a una zona de libre comercio, sino que planteó directamente la constitución de un mercado común.}
\cresumen{La integración latinoamericana se divide en una etapa anterior a la década de 1970, centrada en estabilizar primero el mercado interno, y una posterior al retorno democrático, en la que se encaró simultáneamente el fortalecimiento interno y la apertura a la inversión extranjera. El acercamiento bilateral entre Argentina y Brasil, iniciado en Foz de Iguazú, dio origen al proceso que culminó en el Mercosur, planteado directamente como mercado común.}
\end{cornell}

% ----------------------------------------------------------------
\chapter{El marco de la ALADI}
\chaptermark{El marco de la ALADI}
\label{cap:aladi}
% ----------------------------------------------------------------

\section{El Tratado de Montevideo de 1980 y la creación de la ALADI}

La \textbf{Asociación Latinoamericana de Integración (ALADI)} es un tratado del cono sur de 1980 que abarca prácticamente toda América Latina, orientado a liberar el comercio.

\begin{normacitada}{Tratado de Montevideo de 1980}
Instrumento fundacional de la ALADI, aún vigente. No debe confundirse con los tratados de Montevideo de Derecho Internacional Privado (1889 y 1940), que regulan materias distintas.
\end{normacitada}

La ALADI, si bien incluye el libre comercio, no aspira solamente a constituir una zona de libre comercio, sino a lograr una integración más profunda: un desarrollo social, económico y cultural más global. Su antecedente es la \textbf{ALALC} (Asociación Latinoamericana de Libre Comercio), del Tratado de Montevideo de 1960, orientada exclusivamente al libre comercio y que finalmente fracasó, dando paso a la ALADI.

\section{Fin versus medio en los tratados marco}

Cuando un tratado habla de «establecer en forma gradual y progresiva un mercado común latinoamericano», el mercado común no es el fin del tratado, sino el \textbf{medio} para alcanzar un fin más amplio: la integración latinoamericana progresiva --distinción ya presentada en el Capítulo~3 al tratar la relación entre objetivo y medio--. Lo que no debería cambiar entre distintos tratados es la finalidad; lo que puede variar es el medio elegido. El Tratado de Asunción, antecedente del Mercosur, retoma esta misma lógica: el fin sigue siendo la integración latinoamericana progresiva, mientras que el mercado común del Mercosur se constituye como un medio para lograr esa integración mayor, optimizando la parte productiva, las economías de escala y la complementación comercial. El Mercosur es, en rigor, una \textbf{subregión} dentro de la ALADI.

\section{La cláusula de la nación más favorecida}

\begin{definicion}{Cláusula de la nación más favorecida}
Dentro de un espacio de integración, un Estado miembro no puede otorgar beneficios específicos a otro Estado sin extenderlos al resto de los miembros, salvo que dichos beneficios estén amparados por un acuerdo de alcance parcial o un área de integración específicamente autorizada dentro del marco general.
\end{definicion}

\begin{ejemplo}{La analogía familiar de los regalos}
Si se le compran zapatos solamente a uno de los hijos, el resto reclamará: «para todos es igual». Del mismo modo, dentro de un área de integración, si un Estado le otorga beneficios específicos a otro Estado miembro, está obligado a extenderlos a todos los demás miembros, salvo que exista un acuerdo de alcance parcial que lo autorice a limitar ese beneficio.
\end{ejemplo}

\section{Acuerdos de alcance parcial y acuerdos bilaterales}

La ALADI crea, entre sus órganos, una secretaría general, una de cuyas funciones es permitir y controlar acuerdos bilaterales o multilaterales entre los Estados parte. La ALADI no permite hacer acuerdos de integración con Estados ajenos al bloque --esos deben regirse por el derecho internacional común--, pero sí permite a sus miembros celebrar entre ellos acuerdos bilaterales o multilaterales, de integración o de alcance parcial, siempre que no vayan en contra de sus principios y objetivos generales.

\begin{ejemplo}{La analogía del doble matrimonio}
Pertenecer a dos áreas de integración con el mismo objetivo se compara con casarse dos veces: ¿a los hijos de qué matrimonio se envía a la escuela privada si no alcanza el dinero para ambos? Según la doctrina, un Estado no puede pertenecer a dos áreas de integración con el mismo objetivo, salvo que una sea, por ejemplo, política y la otra social o cultural, es decir, de naturaleza distinta.
\end{ejemplo}

El origen histórico del Mercosur se enmarca en este esquema: con el retorno de la democracia, Argentina y Brasil firmaron numerosos acuerdos bilaterales de alcance parcial, limitados al ámbito comercial, que fueron el antecedente directo del Mercosur (Capítulo~9). Cuando se propuso transformar ese vínculo bilateral en un tratado de integración, Uruguay y Paraguay solicitaron sumarse, dando origen al Tratado de Asunción entre los cuatro países, siempre dentro del ámbito que permite la ALADI. La secretaría de la ALADI registra y controla los distintos acuerdos, verificando que no contraríen los objetivos generales del bloque, precisamente por la vigencia de la cláusula de la nación más favorecida.

\section{Cuadro Cornell}

\begin{cornell}
\cpregunta{¿Qué es la ALADI y cuál es su antecedente?}{La Asociación Latinoamericana de Integración, creada por el Tratado de Montevideo de 1980. Su antecedente es la ALALC, del Tratado de Montevideo de 1960, orientada exclusivamente al libre comercio.}
\cpregunta{¿Cuál es la diferencia entre el fin y el medio en un tratado de integración?}{El fin es el objetivo último (la integración latinoamericana progresiva); el medio es el instrumento institucional elegido (por ejemplo, constituir un mercado común). El medio puede variar; la finalidad, en principio, no debería cambiar.}
\cpregunta{¿En qué consiste la cláusula de la nación más favorecida?}{Un Estado miembro de un área de integración no puede otorgar beneficios específicos a otro sin extenderlos a los demás, salvo acuerdos de alcance parcial expresamente autorizados.}
\cpregunta{¿Qué controla la secretaría general de la ALADI?}{Registra y controla los acuerdos de alcance parcial y de integración que celebran los Estados miembros entre sí, verificando que no contradigan los objetivos generales del bloque.}
\cresumen{La ALADI --creada por el Tratado de Montevideo de 1980 como heredera de la ALALC de 1960-- constituye hoy el marco general de la integración latinoamericana, sujeto a la cláusula de la nación más favorecida y al control de su secretaría general. El Mercosur es una subregión dentro de ese marco mayor, y su origen se explica por los acuerdos bilaterales de alcance parcial entre Argentina y Brasil.}
\end{cornell}

% ----------------------------------------------------------------
\chapter{Integración frente a derecho internacional común}
\chaptermark{Integración frente a derecho internacional}
% ----------------------------------------------------------------

Distinguir un tratado de integración de un tratado internacional común no siempre es sencillo en la práctica, aunque en la teoría la diferencia parezca clara. Tres ejemplos permiten trabajar esta distinción.

\section{El caso del gas boliviano}

\begin{ejemplo}{Compraventa internacional de gas}
La Argentina, al no producir gas suficiente, le compraba gas a Bolivia. Ese tratado de compraventa es un contrato internacional de derecho privado, no un acuerdo enmarcado en un área de integración: Bolivia no tenía obligación de venderle gas a la Argentina al mismo precio al que se lo vendía a otro Estado, porque el vínculo no estaba alcanzado por la cláusula de la nación más favorecida propia de un área de integración.
\end{ejemplo}

La clave para distinguir este caso de un acuerdo de integración está en la finalidad económica de cada parte: la Argentina buscaba proveerse de gas, y Bolivia buscaba cobrar por su producto. No existía una finalidad de beneficio recíproco propia de un área de integración, sino un intercambio de intereses individuales típico de un contrato internacional de compraventa.

\section{El Convenio de Montreal y el transporte aéreo}

\begin{normacitada}{Convenio de Montreal}
Instrumento internacional en materia aeronáutica que regula las relaciones internacionales aéreas de las compañías de transporte. No constituye un tratado de integración, sino una regulación técnica del tráfico aéreo internacional.
\end{normacitada}

\begin{definicion}{Interdependencia}
Elemento distintivo de un área de integración económica: significa que la actitud económica que adopta un Estado repercute directamente en la economía de los demás Estados del bloque, porque comparten un mercado único. Sin ese efecto recíproco, no puede hablarse propiamente de integración, aunque exista una regulación internacional común sobre la materia.
\end{definicion}

En el caso del Convenio de Montreal, lo único que se regula es que los aviones no colisionen entre sí; no existe un objetivo de interdependencia económica plena entre los Estados firmantes. Distinto sería el caso, por ejemplo, de un acuerdo que estableciera libre circulación e igualdad de trato con fines turísticos entre dos Estados: allí sí podría constituirse una integración económica, aunque focalizada en un solo servicio, porque existiría un beneficio económico recíproco.

\section{El tratado de cooperación judicial del Mercosur}

Existe, dentro del Mercosur, un tratado orientado a la cooperación en materia civil, comercial, administrativa y laboral. Este instrumento no es un tratado marco --no es el que crea el área de integración--, sino una norma elaborada por los órganos del Mercosur ya existente, que regula un aspecto específico.

\begin{definicion}{Derecho derivado (primera aproximación)}
Conjunto de normas elaboradas por los órganos propios de un área de integración ya constituida, a diferencia del tratado marco, que es el instrumento fundacional. El tratado de cooperación judicial del Mercosur se califica como derecho derivado y no como un tratado fundacional. Esta distinción se retoma con mayor precisión en el Capítulo~\ref{cap:fuentes-mercosur}.
\end{definicion}

\section{Cuadro Cornell}

\begin{cornell}
\cpregunta{¿Por qué la compraventa de gas entre Argentina y Bolivia no es un tratado de integración?}{Porque es un contrato de derecho privado en el que cada parte persigue su propio interés, sin la interdependencia económica recíproca propia de un área de integración.}
\cpregunta{¿Por qué el Convenio de Montreal no constituye, por sí solo, un tratado de integración?}{Porque regula un aspecto técnico específico --el tráfico aéreo-- sin generar interdependencia económica recíproca.}
\cpregunta{¿Qué es la interdependencia en materia de integración económica?}{El fenómeno por el cual la actitud económica de un Estado repercute directamente en la economía de los demás Estados del bloque, porque comparten un mercado único.}
\cresumen{Distinguir un tratado de integración de un tratado internacional común exige identificar si existe o no interdependencia económica recíproca entre los Estados, más allá de que ambos tipos de instrumentos puedan tratar aspectos económicos, como muestran los casos del gas boliviano y el Convenio de Montreal.}
\end{cornell}

% ----------------------------------------------------------------
\chapter{El Mercosur y la Comunidad Andina dentro de la ALADI}
\chaptermark{Mercosur y Comunidad Andina}
% ----------------------------------------------------------------

\section{El Grupo Andino y su evolución hacia la Comunidad Andina}

Durante la década de 1960, dentro de la ALALC, los dos países con mayor peso económico en Latinoamérica eran Argentina y Brasil, que en los hechos imponían las reglas del juego. Los países andinos --Bolivia, Colombia, Ecuador, Perú y Venezuela-- se agruparon para poder defenderse de esa influencia permanente, sin salir del proyecto de integración latinoamericana, dando origen al \textbf{Grupo Andino}.

\begin{normacitada}{Acuerdo de Cartagena (1969)}
Tratado fundacional del Grupo Andino (hoy Comunidad Andina), integrado originalmente por Bolivia, Colombia, Ecuador, Perú y Venezuela. Su objetivo declarado: promover el desarrollo equilibrado y armónico de los países miembros mediante la integración y la cooperación económica y social.
\end{normacitada}

A diferencia de la ALALC, que fracasó y tuvo que transformarse en la ALADI, el Grupo Andino fue evolucionando hasta convertirse en la actual \textbf{Comunidad Andina}, que hoy cuenta con órganos supranacionales propios, entre ellos un tribunal de justicia supranacional --un desarrollo institucional cercano al que tuvo la Unión Europea--, y cuyos Estados miembros conocen y sostienen activamente su normativa.

\section{La interdependencia económica entre socios}

Dentro de un área de integración resulta habitual que existan mecanismos de ayuda entre los Estados miembros: si existe libre tránsito de personas y uno de los Estados atraviesa una alta tasa de desempleo, se produce migración hacia los Estados con mejores condiciones laborales, lo que a su vez puede generar desempleo en el Estado receptor. Por esta interdependencia, la crisis económica, la hiperinflación o el deterioro de un Estado socio termina repercutiendo, tarde o temprano, en los demás miembros del bloque.

\begin{ejemplo}{La crisis de deuda de Grecia}
La crisis de deuda de Grecia dentro de la Unión Europea obligó a los demás socios europeos a intervenir para asistirla financieramente, mediante el Banco Europeo y el Banco de Inversiones de la Unión Europea (desarrollados en el Capítulo~\ref{cap:organos-consultivos}), imponiéndole a la vez condiciones de control sobre el uso de esos fondos. Este tipo de mecanismos de asistencia financiera entre socios normalmente no está previsto en el tratado fundacional del bloque, sino que forma parte del derecho derivado de la integración.
\end{ejemplo}

\section{Cuadro Cornell}

\begin{cornell}
\cpregunta{¿Cuál es el origen histórico del Grupo Andino (hoy Comunidad Andina)?}{Surgió en 1969, cuando Bolivia, Colombia, Ecuador, Perú y Venezuela se agruparon para contrarrestar la influencia dominante de Argentina y Brasil dentro de la ALALC.}
\cpregunta{¿En qué se diferencia hoy la Comunidad Andina de otros bloques latinoamericanos?}{Cuenta con órganos supranacionales propios, incluido un tribunal de justicia supranacional, con desarrollo institucional comparable al de la Unión Europea.}
\cpregunta{¿Por qué la interdependencia económica exige mecanismos de auxilio financiero entre socios?}{Porque el deterioro económico de un socio repercute en el resto del bloque por la interdependencia del mercado único (ejemplo: la crisis de deuda de Grecia).}
\cresumen{Dentro del marco de la ALADI conviven subregiones como el Mercosur --definido como un paso hacia una integración latinoamericana mayor-- y bloques más antiguos y consolidados como la Comunidad Andina, con órganos supranacionales propios. La interdependencia económica entre socios explica tanto los riesgos compartidos como los mecanismos de auxilio financiero creados para asistir a un Estado en crisis.}
\end{cornell}

% ----------------------------------------------------------------
\chapter{Integración no formal: los casos de África y Asia}
\chaptermark{Integración no formal en África y Asia}
% ----------------------------------------------------------------

\section{Integración política y de defensa sin tratado marco}

Existen procesos de integración que se producen de hecho, sin un tratado marco formal, particularmente en materia de defensa y de política exterior.

\begin{definicion}{Integración no formal}
Proceso de integración que se manifiesta en la práctica --por ejemplo, mediante el voto conjunto y sistemático de un grupo de Estados en organismos internacionales, o mediante políticas de defensa compartidas-- sin que exista un tratado marco que lo formalice institucionalmente. Este fenómeno es frecuente en los bloques africano y asiático.
\end{definicion}

\section{Diversidad cultural y motivación geopolítica en Asia y África}

En Asia y África, una característica distintiva de los procesos de integración es la fuerte presencia de la \textbf{diversidad cultural}. A diferencia de la integración latinoamericana, donde existe cierta identificación entre los pueblos como «hermanos», en esas regiones cada nación conserva una identidad tan profunda que sus miembros no se sienten parte de un mismo todo. Esta situación se expresa con la fórmula «juntos, pero no revueltos»: se hace integración de tipo económico y se comparte un espacio común, pero no se pretende unificar la cultura ni las tradiciones de cada pueblo.

Otra característica de estas integraciones es su \textbf{motivación geopolítica}: responde también a una necesidad externa de presentar un frente común ante las demandas de potencias como Estados Unidos o Rusia, que frecuentemente condicionan la aceptación de socios comerciales a la instalación de bases militares en el territorio. Es más difícil hacer frente a esas presiones estando solo, de modo que la integración surge más por una necesidad defensiva externa que por el deseo genuino de ser socios.

\section{Fracasos de integración en África: causas estructurales}

La gran mayoría de los procesos de integración en África no logran consolidarse en la práctica, a pesar de figurar formalmente creados. Para explicar las causas de estos fracasos, debe considerarse que los Estados africanos son, en su mayoría, de formación reciente: hasta hace pocas décadas eran colonias europeas, y dependieron durante mucho tiempo, económica, legislativa y políticamente, de las decisiones de las potencias coloniales, además de reunir en un mismo territorio culturas internamente muy distintas entre sí.

\begin{ejemplo}{La analogía del padre colonial}
«Porque era colonia, papá no te dejaba pelear; pero una vez que se fue papá, no bajamos la guardia». Con esta imagen se explica que la salida del dominio colonial dejó al descubierto tensiones internas entre culturas distintas que antes eran contenidas por la potencia colonizadora.
\end{ejemplo}

En cuanto a las causas económicas, muchos Estados africanos carecen de riqueza propia consolidada y presentan una economía con mucha precariedad: aunque un país africano cría ganado, si sus estándares fitosanitarios no cumplen los parámetros exigidos por los mercados europeo o americano, sus productos no logran acceder a esos mercados, lo cual limita drásticamente sus posibilidades de desarrollo comercial. La ausencia de una base económica sólida dificulta seriamente la consolidación de la integración, más allá de la voluntad política declarada en los tratados, y esto pone en cuestión una explicación frecuente en ámbitos académicos, según la cual el fracaso de estos procesos se debe únicamente a la falta de órganos supranacionales: un bloque sin capacidad económica de base no puede sostener una integración real, tenga o no órganos supranacionales.

Como contraejemplo dentro del propio continente africano, se menciona el caso de la Comunidad Económica y Monetaria de África Central, presentada como el proceso más avanzado de la región, cuyo Banco Central, sin embargo, se encuentra bajo fuerte influencia de Francia, lo cual condiciona la autonomía real de las decisiones económicas de sus Estados miembros.

\section{Cuadro Cornell}

\begin{cornell}
\cpregunta{¿Qué caracteriza a la integración en Asia y África?}{Se realiza integración económica sin unificar las culturas («juntos, pero no revueltos»); su motivación es frecuentemente geopolítica: presentar un frente común ante potencias externas.}
\cpregunta{¿Cuáles son las causas estructurales del fracaso de gran parte de los procesos de integración africanos?}{La formación reciente de los Estados, su prolongada dependencia colonial, la coexistencia de culturas internamente distintas, y sobre todo la precariedad económica de base, más allá de que exista o no un órgano supranacional.}
\cresumen{Existen procesos de integración no formales, sin tratado marco, especialmente en materia de defensa y política exterior, frecuentes en Asia y en África. Sin embargo, ningún tratado, por más órganos supranacionales que prevea, puede sostener una integración real si los Estados involucrados carecen de una base económica sólida y de un compromiso recíproco genuino entre sus miembros, tal como evidencia el caso africano.}
\end{cornell}

% ================================================================
\part{El Mercosur: origen, fuentes y estructura institucional}
% ================================================================

% ----------------------------------------------------------------
\chapter{Antecedentes y creación del Mercosur}
\chaptermark{Antecedentes y creación}
\label{cap:asuncion}
% ----------------------------------------------------------------

\section{El Mercosur como subregión de la ALADI}

Como se adelantó en el Capítulo~\ref{cap:aladi}, el Mercosur no constituye un espacio de integración autónomo o independiente, sino que se inscribe dentro de un marco institucional más amplio: la ALADI. Esta permite a sus miembros celebrar \textbf{acuerdos de alcance parcial} que no comprometen a la totalidad de sus Estados, sino únicamente a los que los suscriben. El Mercosur es, precisamente, uno de esos espacios subregionales: no es un desprendimiento de la ALADI ni algo distinto de ella, sino un acuerdo que se da \textbf{dentro} de ella, sin contradecir sus objetivos generales.

\section{El acercamiento bilateral entre Argentina y Brasil}

El proceso que desembocó en la creación del Mercosur se remonta al restablecimiento de la democracia en la región durante la década de 1980. En ese contexto, Argentina --bajo la presidencia de Raúl Alfonsín-- y Brasil comenzaron un proceso de aproximación mediante negociaciones bilaterales desarrolladas dentro del ámbito de la ALADI, con antecedente directo en la Declaración de Foz de Iguazú de 1985 (Capítulo~9). Durante este período se celebraron numerosos acuerdos de alcance parcial que funcionaron como un ensayo de lo que vendría después: permitieron a ambos países explorar campos de cooperación, ajustar expectativas y generar confianza mutua.

Continuando ese proceso, ya consolidada la democracia en la región, Argentina, Brasil, Paraguay y Uruguay decidieron avanzar hacia un acuerdo de integración de mayor alcance y profundidad. Estos cuatro países son los \textbf{Estados fundadores del Mercosur}.

\begin{importante}
La decisión de formar el Mercosur fue, fundamentalmente, política: representó una apuesta por consolidar la integración regional como proyecto estratégico de largo plazo, más allá de las variaciones en las coyunturas económicas o en los gobiernos de turno.
\end{importante}

\section{El Tratado de Asunción}

\begin{definicion}{Tratado de Asunción}
Instrumento jurídico internacional suscripto en 1991 por Argentina, Brasil, Paraguay y Uruguay, que crea el Mercado Común del Sur (Mercosur) y constituye el derecho originario del espacio de integración.
\end{definicion}

El Tratado de Asunción cumple, dentro del Mercosur, una función equivalente a la que desempeña una Constitución dentro del ordenamiento jurídico interno de un Estado. Un dato relevante para comprender la profundidad institucional del proceso es que el acuerdo alcanzado por los cuatro Estados fundadores no se limitó a establecer una zona de libre comercio --como había sido habitual en la primera etapa histórica de la integración latinoamericana, descrita en el Capítulo~9--, sino que planteó directamente la constitución de un \textbf{mercado común} (Capítulo~\ref{cap:mercado-comun}).

Como se desarrolló al estudiar el mercado común en general, este nivel de integración aspira a la libre circulación de mercaderías, servicios y personas, y a la libre circulación de los factores de producción --capital y trabajo--, lo cual requiere, como condición previa, la adopción de políticas uniformes entre los Estados miembros. La realidad del bloque muestra, sin embargo, una considerable distancia entre ese objetivo declarado y lo efectivamente alcanzado en cada materia.

\section{El costo de la integración: la cesión de facultades soberanas}

Formar parte de un mercado común implica necesariamente una cesión de facultades: los Estados no pueden tomar decisiones unilaterales en materias que afecten el libre tránsito de personas, capitales, servicios o mercaderías dentro del bloque. Esta restricción de la soberanía decisoria constituye el núcleo de todo proceso de integración: los Estados aceptan límites a su capacidad de actuar autónomamente para lograr, a cambio, los beneficios que la integración ofrece. Las tensiones concretas que genera este costo en la práctica argentina se desarrollan en profundidad en el Capítulo~\ref{cap:tensiones}.

\section{Cuadro Cornell}

\begin{cornell}
\cpregunta{¿Por qué el Mercosur no es un espacio de integración autónomo respecto de la ALADI?}{Porque es uno de los espacios subregionales que la ALADI permite conformar mediante acuerdos de alcance parcial, sin contradecir sus objetivos generales.}
\cpregunta{¿Cuál fue el antecedente directo del Mercosur?}{El acercamiento bilateral entre Argentina y Brasil tras el retorno democrático de la década de 1980, formalizado en la Declaración de Foz de Iguazú (1985) y en sucesivos acuerdos de alcance parcial.}
\cpregunta{¿Qué crea el Tratado de Asunción (1991) y qué función cumple?}{Crea el Mercado Común del Sur entre Argentina, Brasil, Paraguay y Uruguay, y funciona como el derecho originario del bloque, equivalente a una Constitución.}
\cpregunta{¿Por qué se dice que la creación del Mercosur implicó un compromiso mayor que en la primera etapa histórica de la integración latinoamericana?}{Porque no se limitó a una zona de libre comercio, sino que planteó directamente la constitución de un mercado común.}
\cresumen{El Mercosur nació como un espacio subregional dentro de la ALADI, a partir del acercamiento bilateral entre Argentina y Brasil iniciado con el retorno democrático de la región. El Tratado de Asunción (1991), suscripto por los cuatro Estados fundadores, constituye su derecho originario y planteó directamente la meta de un mercado común, con el costo que ello implica: la cesión de facultades soberanas en las materias comprendidas por esa libre circulación.}
\end{cornell}

% ----------------------------------------------------------------
\chapter{Fuentes del derecho del Mercosur}
\chaptermark{Fuentes del derecho del Mercosur}
\label{cap:fuentes-mercosur}
% ----------------------------------------------------------------

\section{Derecho originario}

\begin{definicion}{Derecho originario}
Conjunto de normas de fuente convencional que conforman la estructura fundamental del espacio de integración. En el Mercosur está constituido por el Tratado de Asunción y por aquellos protocolos posteriores que, aunque sean instrumentos jurídicos distintos, se incorporan a él con su mismo rango y lo modifican estructuralmente.
\end{definicion}

Una confusión frecuente surge de la distinción entre tener \textbf{rango} constitucional y ser \textbf{parte} de la Constitución. El Protocolo de Ouro Preto, que modificó la estructura orgánica del bloque (Capítulo~\ref{cap:estructura-mercosur}), no simplemente \emph{tiene} rango equivalente al Tratado de Asunción: pasa a \emph{integrar}, junto con él, el propio derecho originario del sistema. No son dos normas de igual jerarquía coexistiendo; es una única fuente normativa fundamental en la que ambos instrumentos se han fusionado.

\begin{ejemplo}{Las enmiendas a la Constitución de los Estados Unidos}
El derecho originario no se agota en el tratado fundacional: cada vez que se incorpora un nuevo órgano mediante un instrumento posterior --el Protocolo de Ouro Preto o el Protocolo de Olivos, que crea el Tribunal Permanente de Revisión (Capítulo~\ref{cap:tpr})--, ese instrumento pasa a integrar también el derecho originario, del mismo modo en que todas las enmiendas a la Constitución de los Estados Unidos siguen siendo Constitución, aunque se hayan incorporado después.
\end{ejemplo}

El derecho originario del Mercosur está integrado, principalmente, por el Tratado de Asunción (1991), el Protocolo de Ouro Preto (1994), que estableció la estructura institucional definitiva del bloque, y el Protocolo de Olivos, que creó el Tribunal Permanente de Revisión y adicionó nuevos órganos a la estructura.

\section{Derecho derivado}

\begin{definicion}{Derecho derivado}
Conjunto de normas dictadas por los órganos que el tratado fundacional creó, en ejercicio de las funciones que ese tratado les atribuyó: resoluciones, decisiones, directivas y otros actos normativos emanados de la estructura institucional del bloque.
\end{definicion}

De manera análoga a como la Constitución Nacional crea los poderes Ejecutivo, Legislativo y Judicial, el Tratado de Asunción crea una estructura orgánica y dota a cada órgano de determinadas funciones y competencias. La actividad de esos órganos --cuando actúan conforme a las competencias que el tratado les asignó-- constituye derecho derivado.

\begin{importante}
El criterio fundamental para distinguir entre derecho originario y derecho derivado es la \textbf{cualidad en que actúa} la autoridad emisora, no la identidad de la persona física. Cuando los presidentes de los Estados miembros se reúnen y firman un instrumento actuando en representación de su Estado --como ocurrió con el Tratado de Asunción--, generan derecho originario. Cuando esos mismos presidentes, u otros funcionarios, actúan integrando un órgano creado por el tratado --por ejemplo, el Consejo del Mercado Común--, lo que producen es derecho derivado, aunque se trate de las mismas personas físicas.
\end{importante}

\begin{table}[H]
\centering
\renewcommand{\arraystretch}{1.3}
\begin{tabularx}{\textwidth}{>{\raggedright\arraybackslash}p{0.25\textwidth}>{\raggedright\arraybackslash}X>{\raggedright\arraybackslash}X}
\toprule
\textbf{Aspecto} & \textbf{Derecho originario} & \textbf{Derecho derivado} \\
\midrule
Equivalente interno & Constitución Nacional & Leyes, resoluciones \\
Instrumento & Tratado de Asunción y protocolos que lo modifican (Ouro Preto, Olivos) & Resoluciones, decisiones y directivas de los órganos del Mercosur \\
Quién lo dicta & Los Estados parte, en representación de su Estado & Los órganos creados por el tratado, dentro de su competencia \\
Función & Crea el sistema, la estructura orgánica y las reglas fundamentales & Desarrolla y ejecuta las funciones atribuidas por el derecho originario \\
Modificación & Requiere nueva convención de los Estados & Puede modificarse dentro de los límites de la competencia del órgano \\
\bottomrule
\end{tabularx}
\caption{Derecho originario y derecho derivado en el Mercosur.}
\end{table}

\begin{importante}
Determinar si una norma del bloque es derecho originario o derivado, y verificar que el órgano que la dictó era competente para hacerlo --principio de atribución de facultades, Capítulo~\ref{cap:atribucion}--, es la manera de establecer si esa norma resulta exigible. Una resolución que exceda las competencias de su órgano no obliga a los Estados miembros, aunque formalmente haya sido emitida.
\end{importante}

\section{Cuadro Cornell}

\begin{cornell}
\cpregunta{¿Qué instrumentos integran el derecho originario del Mercosur?}{El Tratado de Asunción (1991), el Protocolo de Ouro Preto (1994) y el Protocolo de Olivos, porque cada instrumento posterior que incorpora un nuevo órgano pasa a integrar también el derecho originario.}
\cpregunta{¿Cuál es el criterio para distinguir derecho originario de derecho derivado?}{La cualidad en que actúa quien dicta la norma: en representación directa del Estado (originario) o integrando un órgano creado por el tratado (derivado).}
\cpregunta{¿Qué consecuencia tiene que un órgano dicte una norma fuera de su competencia?}{La norma no obliga a los Estados miembros, aunque formalmente haya sido emitida: el principio de atribución de facultades opera como límite y como garantía.}
\cresumen{El derecho originario del Mercosur --Tratado de Asunción, Protocolo de Ouro Preto y Protocolo de Olivos-- funciona como su Constitución; el derecho derivado son las normas dictadas por los órganos creados por ese derecho originario, dentro de las competencias que este les atribuyó. La distinción depende de la función en la que actúa quien produce la norma, y su verificación es indispensable para establecer si esa norma resulta jurídicamente exigible.}
\end{cornell}

% ----------------------------------------------------------------
\chapter{Estructura institucional: de Asunción a Ouro Preto}
\chaptermark{Estructura institucional}
\label{cap:estructura-mercosur}
% ----------------------------------------------------------------

\section{La estructura provisoria del Tratado de Asunción}

Durante la vigencia inicial del Tratado de Asunción (1991), previa al Protocolo de Ouro Preto, únicamente tres órganos del Mercosur contaban con capacidad decisoria propia, mientras que el resto cumplía funciones de asesoramiento o consulta.

\begin{definicion}{Órgano intergubernamental de carácter decisorio}
Los tres órganos originarios del Mercosur con poder decisorio compartían tres características: (1) ser órganos \textbf{intergubernamentales}, es decir, representar directamente a los gobiernos y no a la ciudadanía; (2) decidir por \textbf{consenso}, con la presencia y el acuerdo de todos los Estados presentes en la reunión (desarrollado en el Capítulo~\ref{cap:consenso}); y (3) tener \textbf{poder decisorio}, a diferencia del resto de los órganos, de carácter consultivo.
\end{definicion}

Además de la Secretaría --que en esta etapa no era un órgano propiamente dicho, sino una secretaría administrativa-- existían otros organismos de consulta: el \textbf{Comité Económico y Social}, con participación de sindicatos, asociaciones de consumidores y representantes empresariales, de carácter estrictamente consultivo; y la \textbf{Comisión Parlamentaria Conjunta}, órgano de carácter informativo, que elaboraba informes y estadísticas comparativas sobre el estado de implementación de las normas del Mercosur en cada país, pero sin función legislativa alguna.

\begin{nota}
Pese a llamarse «parlamentaria», la Comisión Parlamentaria Conjunta no era un parlamento: no legislaba, solo informaba.
\end{nota}

Esta estructura institucional era provisoria: se había establecido un período de transición de apenas \textbf{tres años} desde 1991, durante el cual estos órganos gobernarían el proceso de formación del arancel externo común y de la zona de libre comercio, con vencimiento de mandatos en 1994.

\begin{importante}
La concepción original apuntaba a una estructura no meramente intergubernamental, sino \textbf{supranacional}, en el pensamiento de quienes trabajaron en el proyecto originario, aunque ese objetivo no llegó a alcanzarse plenamente.
\end{importante}

\section{El Protocolo de Ouro Preto y la estructura definitiva}

Ante el vencimiento de los mandatos en 1994, comenzó a trabajarse el Protocolo de Ouro Preto sobre la estructura institucional definitiva del Mercosur. En rigor, no se salió nunca del todo del período de transición, sino que se decidió tomar los mismos órganos existentes como definitivos y continuar trabajando a partir de ese esquema.

\subsection{La reforma constitucional argentina de 1994}

Este proceso se vincula con un hecho paralelo de gran relevancia ocurrido en la Argentina ese mismo año: la reforma de la Constitución Nacional, que incorporó una previsión especial para los procesos de integración.

\begin{normacitada}{Constitución Nacional argentina, artículo 75, inciso 24}
«Corresponde al Congreso: aprobar tratados de integración que deleguen competencias y jurisdicción a organizaciones supraestatales en condiciones de reciprocidad e igualdad, y que respeten el orden democrático y los derechos humanos. Las normas dictadas en su consecuencia tienen jerarquía superior a las leyes.»
\end{normacitada}

En virtud de esta norma, el Congreso argentino quedó facultado, con una mayoría especial, a delegar facultades legislativas propias de la Argentina en órganos de integración creados con carácter supranacional, otorgando además a esas normas jerarquía superior a la ley interna. La Argentina fue el único país de la región que, en ese momento, previó constitucionalmente la posibilidad de ceder competencias a órganos supranacionales, dejando abierta la vía para no requerir en el futuro una nueva reforma constitucional en caso de avanzarse hacia ese modelo.

\subsection{Continuidades y cambios respecto de la etapa de Asunción}

El Protocolo de Ouro Preto (1994) mantuvo la misma estructura de base que el Tratado de Asunción: el \textbf{Consejo del Mercado Común (CMC)}, el \textbf{Grupo Mercado Común (GMC)} y la \textbf{Comisión de Comercio del Mercosur (CCM)} continuaron siendo los tres órganos con función decisoria, conservando la intergubernamentalidad, la decisión por consenso y la representación directa de los gobiernos.

En la etapa de Asunción, el Consejo debía reunirse al menos una vez al año; con Ouro Preto pasó a reunirse \textbf{dos veces al año}. La presidencia del Consejo es \textbf{alterna}: rota entre los Estados parte por orden alfabético y por períodos de un año. Dado que el Mercosur posee personería jurídica propia y se presenta como bloque ante organismos internacionales, quien ejerce esa representación es el presidente en ejercicio, denominado presidente \emph{pro tempore}.

\begin{nota}
Los protocolos del Mercosur suelen llevar el nombre del lugar donde fueron firmados, coincidente con el país que ejercía la presidencia pro tempore en ese momento: el Protocolo de Ushuaia se firmó bajo presidencia argentina, el de Ouro Preto corresponde a Brasil, y el de Las Leñas se firmó en Mendoza, Argentina.
\end{nota}

Hasta el Protocolo de Ouro Preto, ninguno de los órganos del Mercosur tenía sede fija: las reuniones se organizaban en el país que ejercía la presidencia pro tempore. Recién con este protocolo, la \textbf{Secretaría} del Mercosur dejó de ser una simple oficina administrativa para adquirir estatus de \textbf{órgano}, con sede permanente en \textbf{Montevideo, Uruguay}. Con el correr de las reformas posteriores, la Secretaría dejó de ser puramente administrativa para adquirir también una función técnica: además de custodiar los documentos, puede advertir cuestiones vinculadas a la aplicación de las normas. Cada Estado debe ratificar internamente una norma del Mercosur mediante sus propios actos y leyes internas; la Secretaría recibe y reúne todas esas ratificaciones y, una vez completas, notifica que la norma entra en vigor simultáneamente para todos los Estados parte a partir de determinada fecha.

\section{Los tres órganos con capacidad decisoria}

\begin{definicion}{Consejo del Mercado Común (CMC)}
Órgano de máxima autoridad y nivel decisorio más alto del Mercosur, integrado por los presidentes de los Estados miembros, los ministros de Economía y los cancilleres. Define la política general y el rumbo estratégico del bloque: define orientaciones estratégicas, la relación política con otros espacios de integración, cuestiones fundamentales de política comercial y aranceles, y la aprobación de nuevos miembros o suspensiones.
\end{definicion}

La composición en este nivel máximo responde a la naturaleza de las decisiones que se toman: afectan simultáneamente dimensiones económicas (de allí la participación de los ministros de Economía) y de política exterior (de allí la de los cancilleres), y requieren, en última instancia, la aprobación política de cada jefe de Estado. El Consejo es el órgano que «maneja el timón» del Mercosur.

\begin{definicion}{Grupo Mercado Común (GMC)}
Órgano ejecutivo e implementador de las políticas definidas por el Consejo, integrado por funcionarios técnicos del Ministerio de Economía de cada Estado, organizados en \textbf{subgrupos de trabajo} temáticos --por ejemplo, el Subgrupo de Trabajo N.º~7, dedicado a la protección del consumidor--.
\end{definicion}

Si el Consejo decide que deben eliminarse ciertos aranceles entre los miembros, es el GMC quien diseña el cronograma, identifica los obstáculos prácticos, coordina con las autoridades nacionales y propone los mecanismos específicos de implementación.

\begin{definicion}{Comisión de Comercio del Mercosur (CCM)}
Órgano especializado responsable de la administración técnica de las cuestiones comerciales entre los Estados miembros --aranceles, barreras arancelarias y para-arancelarias--, organizado en \textbf{subcomités} por materias, como el subcomité en materia agropecuaria.
\end{definicion}

\begin{table}[H]
\centering
\renewcommand{\arraystretch}{1.3}
\begin{tabularx}{\textwidth}{>{\raggedright\arraybackslash}p{0.22\textwidth}>{\raggedright\arraybackslash}X>{\raggedright\arraybackslash}X}
\toprule
\textbf{Órgano} & \textbf{Composición} & \textbf{Función principal} \\
\midrule
Consejo del Mercado Común (CMC) & Presidentes, ministros de Economía y cancilleres & Define la política y el rumbo general del bloque \\
Grupo Mercado Común (GMC) & Funcionarios técnicos, organizados en subgrupos temáticos & Ejecuta e implementa las políticas fijadas por el Consejo \\
Comisión de Comercio del Mercosur (CCM) & Funcionarios especializados, organizados en subcomités temáticos & Regula y administra las cuestiones estrictamente comerciales \\
\bottomrule
\end{tabularx}
\caption{Estructura institucional decisoria del Mercosur.}
\end{table}

Un aspecto fundamental es que estos tres órganos son \textbf{intergubernamentales}: están integrados por representantes de los gobiernos de cada Estado miembro, y no por funcionarios independientes ni por representantes directos de la ciudadanía. Esta característica marca una diferencia fundamental respecto de otros modelos de integración, como la Comisión Europea (Capítulo~\ref{cap:organos-ue}).

\section{Cuadro Cornell}

\begin{cornell}
\cpregunta{¿Qué tres características compartían los órganos decisorios originarios del Mercosur?}{Ser intergubernamentales, decidir por consenso y tener poder decisorio (a diferencia de los órganos consultivos).}
\cpregunta{¿Qué modificó el Protocolo de Ouro Preto (1994) respecto de la etapa de Asunción?}{Estableció la estructura institucional definitiva, manteniendo los mismos tres órganos decisorios, y convirtió a la Secretaría en un órgano con sede fija en Montevideo.}
\cpregunta{¿Qué previó la reforma constitucional argentina de 1994 en materia de integración?}{El artículo 75, inciso 24, facultó al Congreso a aprobar tratados de integración que deleguen competencias en órganos supranacionales, con jerarquía superior a las leyes.}
\cpregunta{¿Cuáles son los tres órganos decisorios del Mercosur y sus funciones?}{El Consejo del Mercado Común define el rumbo político; el Grupo Mercado Común ejecuta esas políticas; la Comisión de Comercio administra las cuestiones comerciales.}
\cresumen{La estructura provisoria del Tratado de Asunción, con mandatos que vencían en 1994, fue reemplazada por la estructura definitiva del Protocolo de Ouro Preto, que mantuvo como órganos decisorios al Consejo del Mercado Común, el Grupo Mercado Común y la Comisión de Comercio --todos intergubernamentales y decisorios por consenso-- y dotó de estatus de órgano a la Secretaría, con sede en Montevideo.}
\end{cornell}

% ----------------------------------------------------------------
\chapter{Nuevos órganos: PARLASUR y Tribunal Permanente de Revisión}
\chaptermark{PARLASUR y Tribunal Permanente de Revisión}
\label{cap:tpr}
% ----------------------------------------------------------------

Una de las virtudes de la estructura definitiva adoptada en Ouro Preto fue precisamente no ser rígida: podía modificarse cada vez que resultara necesario, sin depender de plazos de caducidad como en la etapa anterior. De este modo pudieron crearse nuevos órganos --el Parlamento del Mercosur y el Tribunal Permanente de Revisión-- sin necesidad de reformar todo el esquema institucional.

\section{El Parlamento del Mercosur (PARLASUR)}

En la estructura original solo existía la Comisión Parlamentaria Conjunta, de carácter meramente informativo (Capítulo~\ref{cap:estructura-mercosur}). A través del \textbf{Tratado de Olivos} se dio un paso adicional con la creación del \textbf{Parlamento del Mercosur}, integrado por diputados que --a diferencia de los órganos tradicionales, todos intergubernamentales-- representan directamente a la sociedad, a la población del Mercosur, y no a los gobiernos.

\begin{importante}
El PARLASUR es el primer órgano del Mercosur que sale de la estructura intergubernamental tradicional. Sin embargo, en la práctica, los diputados llegan con una orientación política establecida por sus respectivos gobiernos y no se apartan totalmente de esa posición.
\end{importante}

La implementación del Parlamento fue gradual: los primeros diputados fueron directamente designados por cada uno de los Estados, y con el tiempo se le fueron incorporando funciones adicionales.

\begin{importante}
Los actos que emite el PARLASUR \textbf{no son vinculantes} para los Estados, a diferencia de las decisiones del Consejo, el Grupo o la Comisión de Comercio. La función de dictar normas obligatorias sigue en manos de esos tres órganos: si bien se modificó la denominación de los actos parlamentarios, en sustancia continúan siendo informes, dictámenes o proyectos de carácter no vinculante, sin que el PARLASUR constituya un órgano legislativo propiamente dicho.
\end{importante}

\section{De la negociación directa al arbitraje: el Protocolo de Brasilia}

Durante la vigencia del Tratado de Asunción no existía un órgano jurisdiccional dentro del Mercosur, sino un mecanismo de solución de controversias regulado por el \textbf{Protocolo de Brasilia}.

\begin{formula}{Mecanismo escalonado del Protocolo de Brasilia}
\begin{enumerate}
    \item \textbf{Negociaciones directas} entre los Estados en conflicto.
    \item Si fracasan, interviene el \textbf{Grupo Mercado Común} como mediador --pudiendo solicitar la colaboración técnica de la Comisión de Comercio--. Su dictamen, como el de todo mediador, no es vinculante.
    \item Si las partes no aceptan la propuesta del Grupo Mercado Común, se recurre a un \textbf{arbitraje}.
\end{enumerate}
\end{formula}

\begin{definicion}{Arbitraje ad hoc}
Arbitraje que se crea puntualmente para un caso concreto y se disuelve al terminarlo. Ante un nuevo conflicto entre los mismos Estados sobre el mismo tema, debe nombrarse otro arbitraje, que no está obligado a decidir en el mismo sentido que el anterior.
\end{definicion}

\begin{importante}
La falta de continuidad propia del arbitraje ad hoc impide lograr uniformidad de interpretación y aplicación de las normas del Mercosur, a diferencia de lo que ocurre con un tribunal permanente como el Tribunal de Justicia de la Unión Europea (Parte~VI).
\end{importante}

\section{El Tribunal Permanente de Revisión}

A través del Tratado de Olivos se creó el \textbf{Tribunal Permanente de Revisión (TPR)}, un órgano que no existía en la estructura originaria.

\begin{definicion}{Tribunal Permanente de Revisión}
La doctrina prefiere hablar de «tribunal arbitral» antes que de «tribunal jurisdiccional». El calificativo «permanente» no implica que el Tribunal sesione todo el tiempo, sino que sus integrantes están designados de antemano y se convocan cuando se presenta un caso concreto, a diferencia del Tribunal de Justicia de la Unión Europea, que sí trabaja de modo continuo interpretando normas.
\end{definicion}

La función original y principal del TPR consiste en \textbf{revisar el laudo arbitral} dictado en la instancia anterior, dando así una instancia de revisión judicial que antes no existía para los arbitrajes del Mercosur. Como todo tribunal revisor, debe atenerse a lo que consta en el expediente: no puede ampliar el objeto del litigio, sino interpretar y verificar si se cumplió o no el derecho conforme lo alegado. Adicionalmente, tiene otras funciones, entre ellas intervenir en medidas cautelares.

\begin{importante}
El Tribunal Permanente de Revisión no constituye una institución autónoma ni supralegal: conserva facultades bastante restringidas. La sola creación de un órgano no equivale a dotarlo de las funciones propias de un tribunal jurisdiccional pleno. Este sistema se compara en detalle con el de la Unión Europea en la Parte~VI.
\end{importante}

\section{Cuadro Cornell}

\begin{cornell}
\cpregunta{¿Qué diferencia al PARLASUR de los demás órganos del Mercosur?}{Es el primer órgano que representa directamente a la población y no a los gobiernos, aunque en la práctica sus diputados llegan con la orientación política de sus Estados de origen y sus actos no son vinculantes.}
\cpregunta{¿Cuáles son los tres pasos del mecanismo del Protocolo de Brasilia?}{Negociaciones directas, mediación (no vinculante) del Grupo Mercado Común y, en caso de fracaso, arbitraje ad hoc.}
\cpregunta{¿Qué problema genera el arbitraje ad hoc?}{Al disolverse tras cada caso, impide lograr uniformidad de interpretación y aplicación de las normas del bloque.}
\cpregunta{¿Cuál es la función principal del Tribunal Permanente de Revisión?}{Revisar el laudo arbitral dictado en la instancia anterior, sin poder ampliar el objeto del litigio, además de intervenir en medidas cautelares.}
\cresumen{El PARLASUR, creado por el Tratado de Olivos, es el primer órgano del Mercosur con representación directa de la población, aunque sus actos no son vinculantes. La solución de controversias evolucionó de la negociación directa y la mediación del Grupo Mercado Común, previstas en el Protocolo de Brasilia, hasta el arbitraje ad hoc y, finalmente, la creación del Tribunal Permanente de Revisión, que revisa laudos arbitrales pero conserva facultades restringidas.}
\end{cornell}

% ----------------------------------------------------------------
\chapter{El consenso como mecanismo de decisión}
\chaptermark{El consenso como mecanismo de decisión}
\label{cap:consenso}
% ----------------------------------------------------------------

\section{El consenso como regla de decisión}

\begin{definicion}{Consenso}
Mecanismo de adopción de normas que exige unanimidad: la norma se dicta únicamente si todos los órganos y miembros que participan en la decisión llegan a un acuerdo. Un solo voto en contra impide que la norma llegue a nacer.
\end{definicion}

Los tres órganos del Mercosur con capacidad decisoria --el Consejo, el Grupo Mercado Común y la Comisión de Comercio-- adoptan sus normas por consenso, es decir, por unanimidad. La consecuencia es decisiva: si un solo Estado miembro no está de acuerdo, la norma directamente no se dicta. No es que la norma exista pero no se aplique a ese Estado, sino que la norma \textbf{no llega a nacer}: no existe jurídicamente y no genera obligación para ningún Estado.

\begin{importante}
El consenso es un mecanismo binario: o todos acuerdan y la norma existe, o al menos uno se opone y la norma simplemente no nace.
\end{importante}

Un Estado que se opone a una norma no queda obligado por ella simplemente porque los demás la aprobaron: nadie se obliga sin su consentimiento, principio fundamental del derecho internacional convencional que el Mercosur mantiene en su mecanismo de consenso. Con este sistema, cada Estado tiene poder de veto sobre cualquier norma a la que se oponga.

\section{Comparación con el sistema de mayorías de la Unión Europea}

La Unión Europea, en contraste directo, adopta las decisiones según distintos tipos de mayorías: simples, absolutas, calificadas y, en ciertos casos, dobles mayorías (mayoría de votos y de representación poblacional). El sistema de mayorías implica que una norma puede aprobarse aun cuando algunos Estados se oponen: los disidentes no pueden impedirla por voto de veto, y deben acatarla igual.

Cuando la entonces Comunidad Económica Europea fue fundada en 1957, integrada por apenas seis Estados, sus fundadores optaron ya por un sistema de mayorías en lugar de unanimidad, reconociendo que la unanimidad habría resultado impracticable. Hoy, con veintisiete Estados miembros, una regla de unanimidad habría hecho prácticamente imposible dictar cualquier norma, porque siempre habría al menos un Estado en desacuerdo. Además del mecanismo de decisión, la estructura institucional de la Unión Europea es notablemente más compleja: cuenta con órganos que representan a la propia Unión como entidad supranacional (la Comisión Europea), órganos que representan a los Estados miembros (el Consejo de la Unión) y órganos que representan a la población (el Parlamento Europeo) --desarrollados en la Parte~V--, una división de poderes ausente en el esquema del Mercosur, donde todos los órganos decisorios son puramente intergubernamentales.

\begin{table}[H]
\centering
\renewcommand{\arraystretch}{1.3}
\begin{tabularx}{\textwidth}{>{\raggedright\arraybackslash}p{0.2\textwidth}>{\raggedright\arraybackslash}X>{\raggedright\arraybackslash}X}
\toprule
\textbf{Aspecto} & \textbf{Mercosur} & \textbf{Unión Europea} \\
\midrule
Mecanismo de decisión & Consenso (unanimidad) & Mayoría (simple, absoluta, calificada, doble) \\
Efecto de un voto en contra & La norma no se dicta & La norma puede aprobarse igual, sin ese Estado \\
Representación en los órganos & Solo representantes de los gobiernos (intergubernamental) & Representan a la Unión, a los Estados y a la población \\
Lógica dominante & Negociación internacional & Integración supranacional \\
Miembros originarios & 4 (Argentina, Brasil, Paraguay, Uruguay) & 6 (Francia, Alemania, Italia, Bélgica, Países Bajos, Luxemburgo) \\
\bottomrule
\end{tabularx}
\caption{Comparación entre los sistemas de decisión del Mercosur y la Unión Europea.}
\end{table}

\section{Significado político del consenso}

La exigencia de consenso revela algo profundo sobre la naturaleza política del Mercosur: los Estados miembros no están dispuestos a someterse a la voluntad de la mayoría, y cada uno se reserva el derecho de veto sobre cualquier decisión que considere que afecta sus intereses. Esta característica evidencia, además, la ausencia de una voluntad política real de avanzar hacia órganos supranacionales: si esa voluntad existiera, la estructura del bloque sería distinta, con órganos independientes con capacidad de decisión y con el reconocimiento de que ciertos asuntos quedan fuera del control exclusivo de cada Estado. Esta idea se retoma en el Capítulo~\ref{cap:tensiones} al analizar la brecha entre el discurso integracionista y la práctica concreta de los Estados de la región.

\section{Cuadro Cornell}

\begin{cornell}
\cpregunta{¿Qué efecto produce un voto en contra bajo el mecanismo de consenso?}{Impide que la norma nazca jurídicamente: no es que exista y no obligue al disidente, sino que directamente no llega a dictarse.}
\cpregunta{¿Por qué la Unión Europea adoptó desde su origen un sistema de mayorías y no de consenso?}{Porque incluso con solo seis Estados fundadores (1957) se reconoció que la unanimidad sería impracticable; hoy, con veintisiete miembros, sería directamente inviable.}
\cpregunta{¿Qué revela la persistencia del consenso sobre la voluntad política del Mercosur?}{Que los Estados no están dispuestos a someterse a la voluntad de la mayoría ni a ceder poder real a órganos supranacionales.}
\cresumen{El Mercosur decide por consenso: un solo voto en contra impide que la norma nazca. La Unión Europea, en cambio, decide por distintos tipos de mayorías desde su fundación, lo que le permite funcionar con veintisiete Estados miembros. La persistencia del consenso en el Mercosur revela la ausencia de una voluntad política real de avanzar hacia la supranacionalidad.}
\end{cornell}

% ----------------------------------------------------------------
\chapter{Socios plenos, Estados asociados y la cláusula democrática}
\chaptermark{Socios, asociados y cláusula democrática}
% ----------------------------------------------------------------

\section{Socios plenos y Estados asociados}

El Mercosur distingue entre dos categorías de Estados con grados distintos de vinculación y derechos.

\begin{definicion}{Estado socio pleno}
Estado que se ha incorporado plenamente al Mercosur, asumiendo la obligación de implementar la totalidad de la normativa del bloque --tanto originaria como derivada-- y gozando de plenos derechos de voto en los órganos decisorios (CMC, GMC, CCM).
\end{definicion}

Los cuatro Estados fundadores del Mercosur son Argentina, Brasil, Paraguay y Uruguay. Posteriormente, \textbf{Venezuela} se incorporó como socio pleno, aunque más adelante fue suspendida del bloque.

\begin{definicion}{Estado asociado}
Estado que mantiene una relación de integración parcial con el Mercosur: solo se obliga por la normativa a la que adhiere expresamente, no por su totalidad, y no tiene derecho a voto en los órganos decisorios, aunque puede participar de las reuniones.
\end{definicion}

Entre los Estados asociados se encuentran Chile, Colombia y Bolivia --este último con un trámite en curso de incorporación como socio pleno--. La condición de Estado asociado está ligada a la pertenencia de estos países a la ALADI: por formar parte de ese marco más amplio, se les permite vincularse parcialmente con el Mercosur sin asumir la totalidad de sus compromisos, generando una situación de integración gradual.

\section{El ingreso de Venezuela como socio pleno}

Venezuela se incorporó como socio pleno del Mercosur durante la presidencia de Hugo Chávez, en un contexto de cercanía política con el gobierno argentino de la época. Desde una perspectiva puramente económica, el ingreso de Venezuela no representó un aporte significativo para el bloque: su economía, aunque importante, no contribuía de manera determinante a fortalecer las economías del Mercosur ni ampliaba significativamente el mercado común. Sin embargo, desde una perspectiva política, cumplió una función estratégica: permitió conformar, frente a un Brasil cada vez más fuerte económicamente, un frente regional que compartiera una determinada orientación política y económica.

\section{La cláusula democrática (Protocolo de Ushuaia)}

Una de las premisas fundamentales del Mercosur es que sus Estados miembros deben operar dentro del Estado de Derecho y en un ámbito democrático. Esta premisa se formaliza en la llamada \textbf{cláusula democrática del Mercosur}, contenida en el Protocolo de Ushuaia.

\begin{definicion}{Cláusula democrática}
Establece que los Estados que se apartan de los principios fundamentales de democracia y Estado de Derecho pueden ser \textbf{suspendidos} del bloque.
\end{definicion}

\begin{importante}
Un Estado suspendido pierde el derecho a participar en la toma de decisiones (voto), no puede gozar de los beneficios correspondientes a los Estados parte, puede participar de las reuniones pero sin capacidad decisoria, y permanece en esa condición hasta que se resuelva su reincorporación. La suspensión afecta tanto los derechos como las obligaciones del Estado suspendido.
\end{importante}

\section{El caso de Paraguay: suspensión e ingreso de Venezuela}

En un momento crítico de la historia del Mercosur, el mismo día en que se aprobó el ingreso de Venezuela como socio pleno, se había decidido previamente suspender a Paraguay del bloque, cuyo gobierno de ese momento era cuestionado por no ajustarse a las premisas democráticas del Protocolo de Ushuaia.

\begin{ejemplo}{La «rosca política» en el ingreso de Venezuela}
Paraguay era, precisamente, el Estado que habría votado en contra del ingreso de Venezuela, como parte de su capacidad de veto dentro del sistema de consenso. Al quedar suspendido, Paraguay perdió su derecho a voto, lo que permitió que el ingreso de Venezuela se aprobara sin que todos los miembros plenos completaran el requisito de unanimidad. Cuando posteriormente se restableció el ingreso de Paraguay, ese país cuestionó la validez del ingreso de Venezuela, alegando que «los miembros permanentes no votaron todos» en esa sesión. Este episodio se caracteriza como un ejemplo de «rosca política»: el uso de mecanismos formales --la cláusula democrática, la suspensión-- para alterar el resultado de un proceso de toma de decisiones que, de otro modo, habría sido distinto.
\end{ejemplo}

En la actualidad, el Mercosur cuenta formalmente con cinco Estados socios plenos, incluida Venezuela. Sin embargo, Venezuela se encuentra a su vez suspendida del bloque, lo que genera una situación paradójica: es miembro pleno, pero está privada de los derechos y beneficios correspondientes a esa condición.

\section{Cuadro Cornell}

\begin{cornell}
\cpregunta{¿Qué diferencia a un socio pleno de un Estado asociado?}{El socio pleno está obligado por toda la normativa del bloque y tiene derecho a voto. El Estado asociado solo se obliga por la normativa a la que adhiere expresamente y no tiene derecho a voto.}
\cpregunta{¿Cuáles son los Estados socios plenos actuales del Mercosur?}{Argentina, Brasil, Paraguay, Uruguay y Venezuela --aunque esta última se encuentra suspendida--.}
\cpregunta{¿Qué es la cláusula democrática y qué efectos tiene?}{Prevista en el Protocolo de Ushuaia, permite suspender a un Estado que se aparte de los principios democráticos y del Estado de Derecho, privándolo del voto y de los beneficios del bloque.}
\cpregunta{¿Qué ocurrió con Paraguay y el ingreso de Venezuela?}{Paraguay fue suspendido por apartarse de premisas democráticas, precisamente cuando iba a votar en contra del ingreso de Venezuela; sin su voto, este se aprobó, lo que luego Paraguay cuestionó.}
\cresumen{El Mercosur distingue entre socios plenos --obligados por toda su normativa y con derecho a voto-- y Estados asociados, vinculados parcialmente en el marco de la ALADI. La cláusula democrática del Protocolo de Ushuaia permite suspender a un Estado que se aparte de la democracia y el Estado de Derecho, como ocurrió con Paraguay en el episodio que precedió al ingreso de Venezuela como socio pleno, hoy también suspendida.}
\end{cornell}

% ----------------------------------------------------------------
\chapter{Tensiones entre política interna y compromisos de integración}
\chaptermark{Política interna y compromisos de integración}
\label{cap:tensiones}
% ----------------------------------------------------------------

\section{El costo de la incoherencia: el voto «no positivo» de 2008}

Una de las dificultades recurrentes en el funcionamiento del Mercosur surge de la tensión entre las decisiones que adopta cada Estado en su ámbito interno y los compromisos que ha asumido frente a sus socios del bloque. Los responsables políticos suelen tomar decisiones atendiendo a la coyuntura inmediata, sin advertir que esas decisiones --aunque justificadas por circunstancias internas-- pueden resultar incompatibles con los compromisos externos previamente asumidos.

\begin{ejemplo}{El voto «no positivo» de 2008}
Durante un conflicto relativo a la regulación de retenciones a las exportaciones del sector agropecuario, el Senado de la Nación argentina se encontró en una situación de empate al votar un proyecto de ley. La presidencia del Senado debió ejercer su voto de desempate, que resultó negativo, rechazando la iniciativa. Lo significativo del episodio radica en su cronología: la misma semana en que el Senado argentino rechazaba esta iniciativa, el gobierno nacional estaba suscribiendo, en el ámbito del Mercosur, compromisos de libre circulación con sus socios regionales.
\end{ejemplo}

\begin{importante}
Cualquier límite que un Estado imponga a la libre circulación debería dirigirse únicamente a terceros países ajenos al bloque; nunca debe dirigirse contra sus propios socios, porque de lo contrario se estaría violando directamente el compromiso asumido en la integración.
\end{importante}

Esta clase de incoherencias es frecuente y genera tensiones permanentes: los gobiernos enfrentan presiones políticas internas que, en el corto plazo, generan incentivos para adoptar medidas proteccionistas, mientras que los beneficios de la integración tienden a manifestarse a largo plazo y a distribuirse de manera desigual entre los distintos sectores de la economía.

\section{El costo de las decisiones de corto plazo: privatizaciones y desindustrialización}

Durante ciertos períodos de la historia económica argentina reciente se tomaron decisiones de política económica con consecuencias profundas a largo plazo. Una de las más significativas fue la privatización de los ferrocarriles del Estado, que trasladó en la práctica el transporte de mercaderías hacia el transporte automotor, reorganizando toda la estructura de costos y eficiencia de la economía. De manera contemporánea, se produjo una apertura a productos importados que afectó gravemente a la industria nacional --en barrios como Parque Patricios, que históricamente albergaban numerosas fábricas, entre ellas una reconocida fábrica de chocolates, esas instalaciones fueron clausurando o reconvirtiéndose a otros usos--.

\begin{ejemplo}{La escasez de neumáticos durante la pandemia}
Durante la pandemia de COVID-19 la Argentina experimentó una escasez de neumáticos, porque la capacidad productiva nacional del sector había sido prácticamente eliminada durante los años en que se desmanteló la industria nacional. La escasez no era un problema abstracto: los camiones son las arterias de la economía, y sin neumáticos no pueden circular ni transportar productos desde el campo hacia los centros urbanos, generando un problema concreto de abastecimiento de alimentos y bienes esenciales.
\end{ejemplo}

Este episodio ilustra con crudeza cómo una decisión que parece razonable en el corto plazo --permitir el ingreso de productos importados más baratos-- genera consecuencias que, décadas después, resultan en crisis de abastecimiento.

\section{La necesidad de una política de Estado}

Cuando cada gobierno que asume revierte lo hecho por el anterior, ningún proyecto de largo plazo llega a concretarse: una industria requiere años para desarrollarse, la investigación requiere continuidad, y las políticas de integración requieren predictibilidad. Si un gobierno protege cierto sector industrial con aranceles, pero el siguiente reduce esos aranceles, la industria no tiene estabilidad suficiente para invertir en tecnología, capacitación o expansión.

\begin{ejemplo}{La guerra de Vietnam como ejemplo de continuidad política más allá del signo de gobierno}
Durante casi cuarenta años, sucesivos presidentes de los Estados Unidos, de distintos partidos e ideologías, mantuvieron el compromiso con la guerra de Vietnam, a pesar de su costo humano y económico. Ninguno se atrevió públicamente a reconocer el error ni a ordenar un retiro que implicara asumir personalmente esa responsabilidad, porque hacerlo habría tenido costos políticos enormes. Esta dificultad --la de los responsables políticos para asumir costos inmediatos en beneficio de un resultado de largo plazo, o para reconocer públicamente errores pasados-- es universal en política, y el Mercosur no es la excepción.
\end{ejemplo}

\section{La brecha entre el discurso integracionista y la práctica}

El discurso de integración latinoamericana es poderoso e invoca referencias históricas profundas: los ideales bolivarianos de unidad regional, las aspiraciones sanmartinianas de liberación y unidad, la narrativa de una región que debe hablar con una voz única frente al resto del mundo. Sin embargo, la práctica de los Estados es, frecuentemente, otra: marcada por la desconfianza mutua, por decisiones que contradicen ese discurso, y por la priorización de intereses nacionales inmediatos sobre los compromisos regionales.

Si realmente existiera una voluntad política de avanzar hacia una integración más profunda --comparable a la de la Unión Europea--, la estructura del Mercosur sería completamente distinta: habría órganos independientes, con funcionarios no delegados por los gobiernos, con capacidad de decisión vinculante sobre los Estados. La ausencia de esos órganos no responde a una imposibilidad técnica o económica, sino a una decisión política: ningún Estado miembro está dispuesto a comprometerse de manera permanente e irrevocable con sus socios regionales en la medida en que ello le quitaría control sobre sus propias políticas. La persistencia del consenso (Capítulo~\ref{cap:consenso}) como mecanismo de decisión es, en sí misma, un indicador de esa falta de voluntad: garantiza que ningún Estado será obligado a nada contra su voluntad, pero también que el bloque nunca podrá actuar con la rapidez y coherencia de una entidad verdaderamente integrada.

\begin{importante}
La pregunta fundamental no es técnica, sino política: ¿están realmente los gobiernos dispuestos a subordinar sus decisiones al proceso de integración? La Unión Europea logró avanzar porque sus Estados miembros aceptaron ceder poder a órganos supranacionales que no cambian con cada elección; el Mercosur no ha logrado eso, porque sus Estados no están dispuestos a ceder ese poder.
\end{importante}

\section{Cuadro Cornell}

\begin{cornell}
\cpregunta{¿Qué mostró el episodio del voto «no positivo» de 2008?}{Que un Estado puede restringir internamente el comercio mientras se compromete externamente, en el mismo momento, a la libre circulación dentro del bloque: una incoherencia que solo debería dirigirse a terceros países.}
\cpregunta{¿Qué consecuencia tuvo la desindustrialización sobre la capacidad productiva argentina?}{La pérdida, prácticamente irreversible, de sectores productivos completos, con efectos concretos como la escasez de neumáticos durante la pandemia.}
\cpregunta{¿Por qué es necesaria una política de Estado de largo plazo para sostener la integración?}{Porque la industria, la investigación y la integración requieren continuidad y predictibilidad; si cada gobierno revierte lo hecho por el anterior, ningún proyecto de largo plazo se concreta.}
\cpregunta{¿Qué indica la persistencia del consenso sobre la voluntad política real del Mercosur?}{Que los Estados no están dispuestos a ceder poder real a órganos supranacionales, a diferencia de lo ocurrido en la Unión Europea.}
\cresumen{El funcionamiento real del Mercosur revela una tensión permanente entre los compromisos asumidos frente a los socios y las decisiones de política interna, agravada por la falta de continuidad de las políticas de Estado entre gobiernos sucesivos. La brecha entre el discurso integracionista y la práctica --evidenciada en episodios como el voto «no positivo» o la desindustrialización-- muestra que la principal barrera para una integración más profunda no es institucional, sino la falta de voluntad política de los Estados para ceder poder real a órganos supranacionales.}
\end{cornell}

% ================================================================
\part{La Unión Europea: origen y estructura institucional}
% ================================================================

Frente al modelo predominantemente intergubernamental del Mercosur, la Unión Europea ofrece el caso de mayor profundidad institucional entre los procesos de integración estudiados. Esta parte recorre su origen histórico, su compleja arquitectura de pilares, sus órganos principales y su evolución hasta el Tratado de Lisboa, cerrando con los límites reales que la soberanía nacional sigue imponiendo a su proyecto integracionista.

% ----------------------------------------------------------------
\chapter{Orígenes de la integración europea: CECA, EURATOM y CEE}
\chaptermark{Orígenes de la Unión Europea}
\label{cap:origenes-ue}
% ----------------------------------------------------------------

\section{Un desarrollo distinto al del Mercosur}

Existe un contraste de partida entre ambos procesos de integración: mientras que el Mercosur nació ya con la idea de constituir un mercado común y no avanzó mucho más allá de ese objetivo desde entonces (Parte~IV), la Unión Europea tuvo un desarrollo completamente distinto, no enmarcado dentro de un espacio de integración más amplio como sí ocurre con el Mercosur dentro de la ALADI (Capítulo~\ref{cap:aladi}).

\section{Las causas fundacionales: reconstrucción económica y defensa común}

La integración europea surgió de dos grandes necesidades históricas tras la Segunda Guerra Mundial: la \textbf{reconstrucción económica} de Europa y la organización de una \textbf{defensa común}.

\begin{normacitada}{Tratado de París (1951)}
Constituyó la \textbf{Comunidad Europea del Carbón y del Acero (CECA)}, transfiriendo a órganos comunes, de manera exclusiva, la explotación, investigación, producción y comercialización del carbón y del acero. Se eligieron justamente esos dos recursos por ser centrales tanto para la industria como para el armamento.
\end{normacitada}

\begin{ejemplo}{La magnitud de la reconstrucción de posguerra}
Europa se encontraba devastada tras la guerra: ciudades como Londres y París habían sufrido bombardeos de gran magnitud, en una situación comparable a la de un desastre natural de proporciones extremas, donde la reconstrucción exige comenzar prácticamente desde cero. Por esa razón, la industria del acero y del carbón --fuente de energía predominante en ese momento-- resultaba fundamental, lo que llevó a estructurar la CECA alrededor de dos objetivos claros: la industria y el rearme.
\end{ejemplo}

En 1957, mediante los Tratados de Roma, se crearon dos comunidades adicionales: la \textbf{Comunidad Europea de la Energía Atómica (EURATOM)} y la \textbf{Comunidad Económica Europea (CEE)}. Esta última no se limitaba al libre tránsito de mercaderías: incluía también cuestiones monetarias y financieras, es decir, abarcaba mucho más que un simple mercado común.

A diferencia del Mercosur --donde lo económico y lo político quedaron separados en vías institucionales distintas--, la Unión Europea terminó creando una única estructura institucional que maneja todos sus asuntos, cualquiera sea el tema que se esté tratando en cada sesión. La CECA, como comunidad autónoma, desapareció con el tiempo: el carbón y el acero pasaron a ser una mercadería más dentro de la Comunidad Económica Europea, ya sin el carácter estratégico central que tenían en la posguerra, reemplazados por otras fuentes de energía y otros tipos de armamento.

\section{La secuencia de tratados hasta Lisboa}

El proceso de integración europea se fue construyendo mediante una secuencia de tratados sucesivos, cada uno de los cuales amplió o reformó el anterior:

\begin{enumerate}
    \item \textbf{Tratado de París (1951)}: Comunidad Europea del Carbón y el Acero (CECA).
    \item \textbf{Tratados de Roma (1957)}: Comunidad Económica Europea (CEE) y EURATOM.
    \item \textbf{Tratado de Maastricht (1992)}.
    \item \textbf{Tratado de Ámsterdam (1997)}.
    \item \textbf{Tratado de Lisboa (2007)}: el vigente, que incorpora las grandes reformas institucionales (Capítulo~\ref{cap:evolucion-ue}).
\end{enumerate}

\section{Cuadro Cornell}

\begin{cornell}
\cpregunta{¿Qué dos necesidades históricas dieron origen a la integración europea?}{La reconstrucción económica de posguerra y la organización de una defensa común, tras la devastación causada por la Segunda Guerra Mundial.}
\cpregunta{¿Qué creó el Tratado de París (1951) y por qué se eligieron esos dos recursos?}{Creó la CECA, transfiriendo a órganos comunes la explotación del carbón y del acero, recursos centrales tanto para la industria como para el armamento.}
\cpregunta{¿Qué se creó en 1957 y en qué se diferenció la CEE de un simple mercado común?}{Se crearon la EURATOM y la CEE; esta última incluía también cuestiones monetarias y financieras, no solo el libre tránsito de mercaderías.}
\cpregunta{¿En qué se diferencia la estructura institucional de la Unión Europea de la del Mercosur?}{La Unión Europea creó una única estructura institucional para todos sus asuntos; en el Mercosur, lo económico y lo político transitan por vías institucionales separadas.}
\cresumen{La integración europea nació de la necesidad de reconstrucción económica y defensa común de posguerra, con la CECA (Tratado de París, 1951) como primer paso, seguida por la EURATOM y la CEE (Tratados de Roma, 1957). A diferencia del Mercosur, la Unión Europea desarrolló una única estructura institucional para todos sus asuntos, y su proceso normativo se construyó mediante una secuencia de tratados sucesivos hasta el Tratado de Lisboa (2007).}
\end{cornell}

% ----------------------------------------------------------------
\chapter{Los pilares de la Unión Europea y los límites de la supranacionalidad}
\chaptermark{Pilares y límites de la supranacionalidad}
\label{cap:pilares}
% ----------------------------------------------------------------

\section{Una afirmación imprecisa: ¿la Unión Europea decide siempre por mayoría?}

Existe una idea muy extendida pero imprecisa: que la Unión Europea decide siempre por mayoría y que todos sus órganos son supranacionales. Esa afirmación solo es válida para una parte de sus competencias, lo que puede mostrarse mediante un esquema de tres columnas o «pilares» temáticos, correspondiente a la estructura institucional previa al Tratado de Lisboa mencionado en el Capítulo~\ref{cap:origenes-ue}, y que continúa siendo útil para explicar en qué materias los órganos actúan hoy con lógica supranacional y en cuáles no.

\begin{importante}
No todas las relaciones que se dan dentro de la Unión Europea presentan la misma lógica de decisión: afirmar que la Unión Europea decide siempre por mayoría supranacional sería impreciso.
\end{importante}

\section{Primer pilar: el ámbito económico, comercial y financiero}

En esta primera columna se ubican la ex-CECA, la EURATOM y la Comunidad Económica Europea (hoy Unión Europea): todo lo vinculado al libre tránsito de mercaderías, bienes, servicios y capitales, y a la moneda única.

\begin{table}[H]
\centering
\renewcommand{\arraystretch}{1.3}
\begin{tabularx}{\textwidth}{>{\raggedright\arraybackslash}p{0.28\textwidth}>{\raggedright\arraybackslash}X}
\toprule
\textbf{Pilar} & \textbf{Ámbito económico, comercial y financiero} \\
\midrule
Órganos que intervienen & Órganos comunes de la ex-CECA, EURATOM y la Comunidad Económica Europea \\
Forma de decisión & Mayoría (absoluta o doble mayoría, según el tema) \\
Carácter & Supranacional: la norma tiene prioridad frente al derecho interno \\
\bottomrule
\end{tabularx}
\caption{Primer pilar de la Unión Europea.}
\end{table}

Solo cuando los órganos sesionan tratando estos temas económicos, las facultades delegadas son verdaderamente exclusivas y la norma resultante tiene prioridad frente al derecho interno de cada Estado, decidiéndose por mayoría (absoluta o doble mayoría, según el tema).

\section{Segundo y tercer pilar: seguridad, justicia, política exterior y defensa}

Los otros dos ámbitos temáticos comprenden, por un lado, la seguridad interna y la cooperación judicial, y por otro, la política exterior y la defensa común.

\begin{ejemplo}{Europol y Eurojust}
\textbf{Europol} es la policía europea, con jurisdicción en todo el territorio de la Unión en materias como trata de personas o tráfico de drogas. Sus límites son claros: si un agente persigue a un traficante que cruza de un país a otro, no puede efectuar la detención por sí mismo del otro lado de la frontera; lo que sí puede hacer es dar aviso y trabajar en cooperación con la fuerza local. \textbf{Eurojust} es la agencia de cooperación judicial europea: no crea una justicia única, pero establece pautas de colaboración, como el reconocimiento de la competencia de un juez de un Estado en otro y la vigencia de resoluciones judiciales dentro de toda la Unión.
\end{ejemplo}

En este ámbito, las reglas de funcionamiento, las facultades y los presupuestos no se deciden por mayoría, sino que requieren acuerdo entre los Estados, precisamente porque se trata de materias vinculadas a la soberanía y al auxilio de la fuerza pública: ningún Estado cede fácilmente funciones de seguridad interna sin su consentimiento expreso.

\begin{ejemplo}{La posición de la Unión Europea frente a la guerra entre Rusia y Ucrania}
La posición de la Unión Europea frente a la invasión rusa a Ucrania ilustra el funcionamiento de este ámbito: las decisiones sobre qué política adoptar frente a un conflicto de esta naturaleza no se toman por mayoría, sino por consenso, debido a que están vinculadas a decisiones tan sensibles como la eventual participación en un conflicto bélico.
\end{ejemplo}

La Unión Europea trabajó, con dificultades, en la conformación de fuerzas armadas propias para una defensa común (no las fuerzas armadas de un país en particular, sino de la Unión), aunque durante mucho tiempo el aporte de soldados y de presupuesto fue voluntario, ya que este ámbito requiere \textbf{unanimidad} y no simple mayoría: ningún Estado puede ser obligado a delegar la defensa de su propia frontera.

\begin{nota}
La estructura de tres pilares fue formulada, en las fuentes de este libro, de dos maneras que no coinciden exactamente en el orden de sus columnas: una las agrupa en económico, seguridad-justicia y política exterior-defensa; otra distingue expresamente entre la Política Exterior y de Seguridad Común (PESC) como segundo pilar y la Cooperación en Justicia y Asuntos de Interior (JAI) como tercero. En ambos casos las materias comprendidas y su carácter --supranacional en lo económico, intergubernamental en lo demás-- coinciden plenamente; solo varía cuál de esos dos ámbitos se numera segundo y cuál tercero. \% TODO: verificar consistencia entre fuentes originales respecto del orden numérico de los pilares segundo y tercero.
\end{nota}

\begin{importante}
Cada vez que se habla de supranacionalidad y supralegalidad en la Unión Europea, normalmente se hace referencia al aspecto económico, pura y exclusivamente. En materia de seguridad interior o de defensa exterior, esa lógica ya no rige de la misma manera. La estructura de la Unión Europea es compleja porque no se limita a un aspecto económico, comercial o financiero: cuenta con una pluralidad de órganos, cada uno con su propia competencia, y en algunos casos se decide por mayoría y en otros por unanimidad o consenso. La misma estructura institucional --los mismos órganos-- actúa de manera diferente según el pilar en que se encuentre.
\end{importante}

\section{Cuadro Cornell}

\begin{cornell}
\cpregunta{¿En qué ámbito rige efectivamente la supranacionalidad en la Unión Europea?}{Únicamente en el pilar económico, comercial y financiero: allí las normas tienen prioridad sobre el derecho interno y se deciden por mayoría.}
\cpregunta{¿Qué son Europol y Eurojust, y qué límites tienen?}{Europol es la policía europea, sin facultad de detener por sí misma fuera de su jurisdicción de origen; Eurojust es la agencia de cooperación judicial. Ambos operan en un ámbito donde se decide por acuerdo entre Estados, no por mayoría.}
\cpregunta{¿Cómo se decide la política exterior y de defensa común?}{Por consenso o unanimidad, porque afecta directamente la soberanía y la seguridad de cada Estado, como ilustra la posición frente a la guerra entre Rusia y Ucrania.}
\cresumen{La Unión Europea no decide siempre por mayoría supranacional: solo en el pilar económico, comercial y financiero sus normas tienen prioridad sobre el derecho interno; en seguridad, justicia, política exterior y defensa rige la lógica del consenso o la unanimidad, porque los Estados no ceden fácilmente su soberanía en esas materias.}
\end{cornell}

% ----------------------------------------------------------------
\chapter{Órganos principales de la Unión Europea}
\chaptermark{Órganos principales}
\label{cap:organos-ue}
% ----------------------------------------------------------------

\section{Advertencia terminológica: los distintos «Consejos»}

Existen tres instituciones con denominaciones similares, lo que puede generar confusión en los materiales de consulta.

\begin{importante}
\begin{enumerate}
    \item \textbf{Consejo de Europa}: organización internacional independiente. \textbf{No forma parte de la Unión Europea}. Tiene sede en Estrasburgo y se ocupa de derechos humanos.
    \item \textbf{Consejo Europeo}: órgano de la Unión Europea, integrado por los jefes de Estado y de gobierno de los Estados miembros. Traza el rumbo político general.
    \item \textbf{Consejo de la Unión Europea} (o Consejo de Ministros): órgano de la Unión Europea, integrado por representantes ministeriales de cada Estado miembro. Colegisla con el Parlamento.
\end{enumerate}
\end{importante}

\section{El Consejo Europeo}

\begin{definicion}{Consejo Europeo}
Órgano político de más alto nivel de la Unión Europea, integrado por los jefes de Estado y de gobierno de todos los Estados miembros (presidentes en las repúblicas y primeros ministros en las monarquías constitucionales).
\end{definicion}

\begin{nota}
En la Unión Europea, los integrantes se denominan \textbf{Estados miembros}; en el Mercosur, \textbf{Estados parte}. Técnicamente, cuando se habla de «democracia» en este contexto, en rigor debería decirse «sistema republicano», porque hay monarquías constitucionales en la Unión Europea; una monarquía no es una democracia republicana, aunque tenga división de poderes. A las reuniones del Consejo Europeo no concurre el monarca, sino el primer ministro.
\end{nota}

La función del Consejo Europeo es trazar el \textbf{rumbo político general} de la Unión: posiciones internacionales, grandes lineamientos geopolíticos y acuerdos estratégicos. No emite normas comunitarias en sentido técnico; su accionar es más propio del derecho internacional público (relación entre Estados soberanos) que del derecho comunitario.

\begin{ejemplo}{La cuestión Malvinas}
Corresponde al Consejo Europeo decidir, por ejemplo, si la Unión Europea se alinea con Argentina en la cuestión Malvinas o se mantiene al margen, o qué postura toma cuando el tema se trata en la ONU. No habrá una norma de la Unión Europea que diga «Malvinas es argentina»; lo que se define, desde lo político, es la postura que adoptará la Unión.
\end{ejemplo}

\section{El Consejo de la Unión Europea (Consejo de Ministros)}

\begin{definicion}{Consejo de la Unión Europea (Consejo de Ministros)}
Órgano que representa a los Estados miembros en la función legislativa y ejecutiva de la Unión, integrado por un representante de rango ministerial de cada Estado, que varía según el tema en agenda: si se trata un asunto de economía, concurren los ministros de economía; si es agricultura, los de agricultura, y así sucesivamente.
\end{definicion}

El Consejo Europeo marca el rumbo político y el Consejo de Ministros decide qué se hace con lo que se marcó, sin que exista dependencia de uno respecto del otro. El Consejo de Ministros, al representar a los Estados, \textbf{defiende los intereses nacionales}: sus miembros votan siguiendo las instrucciones de sus gobiernos, por mayoría, pero cada uno según la postura de su propio Estado.

\section{El Parlamento Europeo}

\begin{definicion}{Parlamento Europeo}
Órgano de representación ciudadana de la Unión Europea. Sus miembros son elegidos directamente por los ciudadanos de cada Estado miembro, mediante sufragio universal.
\end{definicion}

\begin{nota}
A modo exclusivamente didáctico, puede compararse esta estructura con el sistema argentino: el Consejo de Ministros equivaldría al Senado (representa a los Estados), el Parlamento Europeo a la Cámara de Diputados (representa a la ciudadanía) y la Comisión Europea al Poder Ejecutivo (ejecuta las normas). Se trata de una analogía útil para comprender la lógica de representación, no de una equivalencia literal.
\end{nota}

\begin{definicion}{Proporcionalidad inversa}
El Parlamento tiene un número fijo de bancas: ningún país puede tener más de 96 ni menos de 6 representantes. Los escaños se asignan de forma proporcional a la población, pero de manera inversa: los países más pequeños tienen más representantes per cápita que los países grandes.
\end{definicion}

Para compensar esta situación, cada diputado de un país populoso representa a más ciudadanos cuando vota --puede estar votando «por siete», mientras que el de un país pequeño vota «por uno» o «por dos»--, porque en la sala no entran más de 96 bancas por país. Este sistema evita que tres o cuatro países con grandes poblaciones dominen todas las decisiones.

Los diputados del Parlamento Europeo \textbf{no representan a su país de origen}, sino al ciudadano europeo en general. En consecuencia, se organizan por orientación política y no por nacionalidad: se sientan agrupados por bloque ideológico y no por delegaciones nacionales.

\begin{ejemplo}{Diputados italianos de distinto signo político}
Puede haber tres diputados italianos de derecha y cinco de izquierda. En el recinto se sentarán con sus bloques ideológicos respectivos, no juntos por ser italianos. Cuando votan, lo hacen en virtud de la postura político-partidaria de su ciudadanía, y no del interés del Estado, lo que contrasta con el Consejo de Ministros, donde los representantes defienden el interés nacional.
\end{ejemplo}

\section{La Comisión Europea}

\begin{definicion}{Comisión Europea}
Órgano ejecutivo de la Unión Europea. Está integrada por comisarios, uno por cada Estado miembro, que no representan a sus Estados sino a la Unión Europea en su conjunto. La función de los comisarios equivale a la del Poder Ejecutivo en los sistemas nacionales.
\end{definicion}

La Comisión ejecuta las normas aprobadas por el Consejo y el Parlamento, dicta los actos reglamentarios necesarios para ponerlas en funcionamiento (análogo al decreto reglamentario presidencial) y gestiona el presupuesto de la Unión Europea.

\begin{ejemplo}{La reglamentación de detalle}
Si la norma dispuso que habrá libre tránsito de mercadería de frutas y verduras, es el comisario quien establece la normativa reglamentaria de detalle: qué aranceles sí y qué aranceles no.
\end{ejemplo}

El \textbf{Parlamento controla a la Comisión}: puede interpelar a los comisarios e incluso destituirlos.

\begin{table}[H]
\centering
\renewcommand{\arraystretch}{1.3}
\begin{tabularx}{\textwidth}{>{\raggedright\arraybackslash}p{0.22\textwidth}>{\raggedright\arraybackslash}X>{\raggedright\arraybackslash}X>{\raggedright\arraybackslash}X}
\toprule
\textbf{Órgano} & \textbf{Integración} & \textbf{Representa a} & \textbf{Función principal} \\
\midrule
Consejo Europeo & Jefes de Estado y de gobierno & Los Estados (nivel político) & Dirección política macro \\
Consejo de Ministros & Ministros sectoriales & Los Estados (función legislativa) & Colegislación con el Parlamento \\
Parlamento Europeo & Diputados de sufragio universal & Los ciudadanos europeos & Colegislación; control de la Comisión \\
Comisión Europea & Comisarios (uno por Estado) & La Unión Europea & Ejecución de normas y presupuesto \\
\bottomrule
\end{tabularx}
\caption{Órganos principales de la Unión Europea.}
\end{table}

\section{Cuadro Cornell}

\begin{cornell}
\cpregunta{¿En qué se diferencian el Consejo de Europa, el Consejo Europeo y el Consejo de la Unión Europea?}{El Consejo de Europa no forma parte de la Unión y se ocupa de derechos humanos; el Consejo Europeo fija el rumbo político (jefes de Estado y de gobierno); el Consejo de la Unión (Consejo de Ministros) colegisla con el Parlamento (representantes ministeriales).}
\cpregunta{¿Qué es la proporcionalidad inversa del Parlamento Europeo?}{Los escaños se asignan según la población pero de manera inversa: los países pequeños tienen más representación per cápita, y un diputado de un país grande vota representando a más ciudadanos.}
\cpregunta{¿Por qué los diputados europeos se agrupan por ideología y no por país?}{Porque representan al ciudadano europeo en general, no a su Estado de origen, a diferencia del Consejo de Ministros, cuyos miembros defienden el interés nacional.}
\cpregunta{¿Qué función cumple la Comisión Europea y quién la controla?}{Ejecuta las normas, dicta la reglamentación de detalle y gestiona el presupuesto; el Parlamento la controla y puede interpelar o destituir a sus comisarios.}
\cresumen{La Unión Europea distingue con precisión terminológica entre el Consejo de Europa (ajeno a la Unión), el Consejo Europeo (rumbo político) y el Consejo de la Unión o de Ministros (colegislación estatal). El Parlamento Europeo representa a la ciudadanía mediante un sistema de proporcionalidad inversa y agrupación ideológica, y controla a la Comisión Europea, órgano ejecutivo que aplica las normas y administra el presupuesto.}
\end{cornell}

% ----------------------------------------------------------------
\chapter{Evolución institucional: de la Constitución fracasada al Tratado de Lisboa}
\chaptermark{Evolución institucional: Tratado de Lisboa}
\label{cap:evolucion-ue}
% ----------------------------------------------------------------

\section{Del Consejo dominante al Parlamento colegislador}

La evolución del rol del Parlamento Europeo dentro de la estructura institucional es significativa y puede rastrearse en las propias normas comunitarias. En los orígenes, el Consejo era el que tomaba las decisiones y, en algunos casos, se requería que también interviniera el Parlamento; en las normativas antiguas se encuentran reglamentos dictados solo por el Consejo, mientras que en fechas posteriores comienzan a aparecer reglamentos del Consejo \textbf{y} del Parlamento, porque después del Tratado de Lisboa la intervención del Parlamento es obligatoria. Este cambio fue consecuencia directa del intento frustrado de sancionar una Constitución Europea.

\section{El proyecto de Constitución Europea y su fracaso}

En algún momento la Unión Europea quiso dotarse de una Constitución propia --tal como se adelantó en el Capítulo~\ref{cap:atribucion} al describir los límites del proceso evolutivo de la integración--. La Constitución fue redactada y estaba prácticamente lista para ser aprobada. Sin embargo, en aquellos países que habían acordado someterla a referéndum popular, la respuesta ciudadana fue negativa.

\begin{importante}
La ciudadanía percibía la Constitución Europea como una amenaza a la soberanía nacional y a las identidades culturales. Existió un «cortocircuito» entre el mando y la gente: los dirigentes creían que la Unión Europea era un gran proyecto, mientras que la población lo percibía como una invasión a su soberanía, a su condición de nacional, y consideraba que iba a perder identidad.
\end{importante}

\begin{ejemplo}{El independentismo catalán frente a la ciudadanía europea}
En España, sectores independentistas --como el catalán-- rechazaban la idea de ser ciudadanos europeos, porque ya cuestionaban ser ciudadanos españoles. El rechazo a un nivel de pertenencia mayor se apoyaba en un rechazo previo al nivel de pertenencia inmediatamente anterior.
\end{ejemplo}

\section{La campaña de información y el Tratado de Lisboa}

Ante el rechazo, la Unión Europea cambió de estrategia: lanzó una campaña masiva de difusión e información hacia la ciudadanía y, en lugar de una Constitución, incorporó todas las reformas deseadas dentro del \textbf{Tratado de Lisboa} (2007). Todo lo que se pretendía incluir en la Constitución se incluyó en Lisboa: se cambió la forma, pero la reforma que se quería hacer quedó hecha igual. El resultado fue un tratado votado por los Estados, ubicado por encima de la soberanía del Estado, y el derecho que se pretendía garantizar mediante una Constitución terminó garantizado en el tratado.

\section{Mecanismos de participación ciudadana}

Entre los mecanismos de participación ciudadana que se incorporaron o reforzaron con estas reformas se destacan:

\begin{enumerate}
    \item \textbf{Iniciativa popular}: los ciudadanos europeos pueden impulsar un tema en la agenda legislativa reuniendo un mínimo de firmas (un millón de ciudadanos de al menos siete Estados miembros).
    \item \textbf{Defensor del Pueblo (Ombudsman)}: órgano independiente que actúa autónomamente y protege a los ciudadanos frente a la mala administración de las instituciones europeas.
    \item \textbf{Parlamento colegislador}: el Parlamento deja de ser consultivo y pasa a colegislar con el Consejo de Ministros en igualdad de condiciones.
    \item \textbf{Control ciudadano vía Parlamento}: el Parlamento controla a la Comisión, que ejecuta el presupuesto; los ciudadanos controlan a la Unión Europea a través de sus representantes electos.
    \item \textbf{Transparencia presupuestaria}: publicación detallada de cuánto cuesta el funcionamiento de la Unión Europea por ciudadano.
\end{enumerate}

La transparencia presupuestaria permite tomar conciencia de cuánto cuesta la Unión --por ejemplo, cien euros al año por ciudadano-- y contrastarlo con los derechos de los que se goza, como el derecho a estudiar en la universidad de cualquier lugar de la Unión: analizado de este modo, si se lo sabe aprovechar, el costo resulta reducido.

\begin{ejemplo}{La experiencia cotidiana de la libre circulación en Europa}
En Europa se cruza la calle y se pasa de un país a otro. No es lo mismo que en Argentina, donde pueden transcurrir muchas horas de viaje sin haber salido siquiera de la provincia. En Europa, en una hora se cambia de país, y es habitual pasar el fin de semana en el país vecino, sin fronteras y sin trámites.
\end{ejemplo}

\section{Cuadro Cornell}

\begin{cornell}
\cpregunta{¿Por qué fracasó el proyecto de Constitución Europea?}{Porque la ciudadanía la percibió como una amenaza a la soberanía nacional y a la identidad cultural; en los países que la sometieron a referéndum, la respuesta fue negativa.}
\cpregunta{¿Qué hizo la Unión Europea tras el fracaso de la Constitución?}{Lanzó una campaña de información e incorporó las mismas reformas mediante el Tratado de Lisboa (2007), bajo la forma de un tratado votado por los Estados en lugar de una Constitución.}
\cpregunta{¿Qué cambió respecto del rol del Parlamento tras Lisboa?}{Pasó de tener una intervención limitada a ser colegislador obligatorio junto con el Consejo de Ministros.}
\cpregunta{¿Qué mecanismos de participación ciudadana se incorporaron o reforzaron?}{La iniciativa popular, el Defensor del Pueblo, el Parlamento colegislador, el control ciudadano sobre la Comisión y la transparencia presupuestaria.}
\cresumen{El fracaso del proyecto de Constitución Europea, rechazado por percibirse como una amenaza a la soberanía y la identidad nacional, llevó a la Unión Europea a incorporar las mismas reformas mediante el Tratado de Lisboa (2007), que convirtió al Parlamento en colegislador pleno y reforzó mecanismos de participación ciudadana como la iniciativa popular, el Defensor del Pueblo y la transparencia presupuestaria.}
\end{cornell}

% ----------------------------------------------------------------
\chapter{El caso Brexit}
\chaptermark{El caso Brexit}
% ----------------------------------------------------------------

\section{El estatus especial del Reino Unido}

El caso Brexit se presenta como ejemplo paradigmático de las consecuencias de no informar adecuadamente a la ciudadanía sobre los alcances de la integración. El Reino Unido no fue socio fundador: ingresó después, más por una cuestión representativa y política que por una necesidad económica, y su participación estuvo históricamente teñida de condiciones especiales.

\begin{itemize}
    \item \textbf{No adhirió a la moneda única} (el euro): mantuvo la libra esterlina.
    \item \textbf{Cláusula opt-out inicial}: las normas aprobadas por mayoría \emph{no} se aplicaban automáticamente en el Reino Unido; el gobierno británico debía aceptarlas expresamente.
    \item \textbf{Cláusula opt-out invertida} (posteriormente): las normas \emph{sí} se aplicaban, salvo que el gobierno británico se opusiera específicamente mediante una presentación formal.
\end{itemize}

Era un socio con un privilegio que no tenían Portugal, Italia ni ningún otro Estado: un «socio con coronita».

\section{El Brexit como maniobra política}

El Reino Unido no tenía tanta intención de irse de la Unión Europea: no fue una cuestión de que firmemente se tomara la decisión, sino una maniobra de partido político. El Brexit fue propuesto como un referéndum por el entonces primer ministro, como una jugada política para consolidar apoyo interno, sin prever que la ciudadanía, poco informada sobre las consecuencias reales, votaría afirmativamente.

\begin{ejemplo}{La mecánica del referéndum y sus consecuencias imprevistas}
El primer ministro propuso un referéndum sobre la permanencia en la Unión Europea como estrategia política. La ciudadanía votó a favor del Brexit, muchos sin comprender las consecuencias prácticas; una vez explicadas esas implicaciones, muchos votantes reconocieron: «no lo sabía; si lo hubiera sabido, votaba que no». El proceso de negociación se extendió por más de dos años porque no había preparación alguna para una salida, y el día efectivo de la salida, centenares de camiones quedaron varados en la aduana: la libre circulación se cortó abruptamente y los productores no tenían la documentación exigida por la Unión Europea.
\end{ejemplo}

\section{El common law inglés como complicación adicional}

La complejidad del Brexit se agravó por el sistema jurídico del Reino Unido, de \textbf{creación pretoriana}: no es la norma escrita la que se aplica en primer término, sino los precedentes judiciales, que adquieren principios fundamentales trasladables a otras situaciones idénticas. Cuando el Reino Unido salió de la Unión Europea, toda la normativa comunitaria que había sido incorporada a su sistema dejó de tener vigencia. Al ser un sistema de \emph{common law} basado en precedentes, se planteó qué ley nacional quedaba para reemplazar esa normativa comunitaria aplicada durante décadas: si se quita la norma comunitaria y no hay una ley que la reemplace --porque toda la legislación y la economía se habían acomodado a ese esquema comunitario--, el vacío legal resulta enorme, como efectivamente ocurrió.

\section{Reflexión sobre legislación y consecuencias no previstas}

El caso Brexit permite una reflexión más general sobre la importancia de pensar las consecuencias de la legislación antes de aprobarla: a veces se incorpora algo en una norma sin tomar en consideración las consecuencias a las que ello va a llevar.

\begin{ejemplo}{Embriones en la fecundación asistida}
La legislación argentina permitió la fecundación in vitro y la criopreservación de embriones. Años después surgió el debate sobre qué hacer con los embriones no utilizados: si la ley establece que la vida comienza en la concepción, ¿destruir un embrión es un delito? Esta consecuencia no se había legislado al momento de regular la fecundación asistida.
\end{ejemplo}

\begin{ejemplo}{Matrimonio igualitario y adopción}
Al aprobarse el matrimonio entre personas del mismo sexo en la Argentina, no se previó de forma integral la cuestión de la adopción y el derecho a formar familia. Las discusiones que siguieron demostraron que la ley había regulado el acto sin considerar todas sus consecuencias.
\end{ejemplo}

\begin{importante}
La discusión sobre las consecuencias de una norma debe darse antes de sancionarla: cuando se legisla, no es para conformar el momento, sino para organizar un proyecto de vida.
\end{importante}

\section{Cuadro Cornell}

\begin{cornell}
\cpregunta{¿Qué estatus especial tenía el Reino Unido dentro de la Unión Europea?}{No fue socio fundador, mantuvo la libra esterlina en lugar del euro, y gozó de cláusulas de opt-out (inicial e invertida) respecto de las normas aprobadas por mayoría.}
\cpregunta{¿Por qué el Brexit resultó tan complejo de implementar?}{Porque fue una maniobra política mal calculada, sin preparación para una eventual salida ni información suficiente para la ciudadanía, lo que generó una negociación de más de dos años y un colapso logístico en la aduana.}
\cpregunta{¿Qué complicación adicional generó el sistema de common law inglés?}{Al basarse en precedentes judiciales, la salida de la normativa comunitaria dejó un vacío legal enorme, porque legislación y economía se habían acomodado durante décadas a ese esquema.}
\cpregunta{¿Qué enseñanza general deja el Brexit sobre la técnica legislativa?}{Que es necesario prever las consecuencias de una norma antes de sancionarla, como ilustran también los casos de los embriones en fecundación asistida y la adopción tras el matrimonio igualitario.}
\cresumen{El Reino Unido tuvo un estatus privilegiado dentro de la Unión Europea, y el Brexit surgió como una maniobra política sin preparación ni información suficiente para la ciudadanía. Sus consecuencias fueron prácticas --colapso aduanero, negociación de más de dos años-- y jurídicas --un vacío legal profundo, agravado por el sistema de common law--, y el caso ilustra la importancia de prever las consecuencias de toda decisión normativa antes de adoptarla.}
\end{cornell}

% ----------------------------------------------------------------
\chapter{Órganos consultivos y financieros; tipos de normas}
\chaptermark{Órganos consultivos, financieros y normas}
\label{cap:organos-consultivos}
% ----------------------------------------------------------------

\section{Órganos consultivos}

Además de los órganos decisorios (Capítulo~\ref{cap:organos-ue}), la Unión Europea cuenta con órganos consultivos que asesoran en materias específicas: el \textbf{Comité Económico y Social}, que representa a los diferentes sectores económicos y sociales europeos --empleadores, trabajadores, sociedad civil-- y emite dictámenes sobre propuestas legislativas; y el \textbf{Comité de las Regiones}, que representa a las regiones y entidades locales y se reúne para tratar cuestiones que afectan especialmente a determinadas áreas geográficas.

\begin{ejemplo}{El Comité de las Regiones del Mediterráneo}
En una reunión del Comité de las Regiones del Mediterráneo se analizan los problemas específicos que tienen esos países por compartir el Mediterráneo --como la crisis migratoria abordada en el Capítulo~\ref{cap:tension-discurso}--: problemas que a Alemania no le preocupan del mismo modo. Son cuestiones que se tratan mejor en ese formato regional, porque de otro modo, al pasar al pleno general, los Estados que no comparten esa situación no les prestan la misma atención.
\end{ejemplo}

\section{Órganos financieros}

\begin{definicion}{Banco Central Europeo (BCE)}
Establece la política monetaria de la zona euro: determina cuánto puede emitir cada Estado miembro, fija el tipo de cambio del euro, establece el tipo de interés bancario y regula toda la política económico-financiera de la zona euro.
\end{definicion}

\begin{ejemplo}{El euro y la emisión monetaria}
El euro que se usa en España tiene la cara del rey español, pero cuánto puede emitir España lo decide el Banco Central Europeo. Ningún país puede «prender la maquinita» por sí solo, porque desestabiliza la moneda en todo el territorio --idea ya adelantada al tratar la comunidad económica en el Capítulo~7--.
\end{ejemplo}

\begin{definicion}{Banco Europeo de Inversiones (BEI)}
Presta dinero a los Estados miembros de la Unión Europea para proyectos de desarrollo (infraestructura, educación, etc.), siempre que favorezcan el objetivo general de la integración europea.
\end{definicion}

En lugar de acudir al Fondo Monetario Internacional, los Estados recurren al BEI; pero no se presta para cualquier finalidad, sino solo para favorecer a la integración europea. Incluso cuando un país en dificultades pide un préstamo, el BEI evalúa si el proyecto fortalece la estabilidad regional, porque un Estado económicamente inestable desestabiliza al conjunto --el mismo mecanismo de auxilio financiero recíproco descrito, para el caso griego, en el Capítulo~12--.

\section{Tipos de normas de la Unión Europea}

\begin{formula}{El tipo de norma depende del efecto}
En la Unión Europea, el tipo de norma depende del \textbf{tipo de efecto} que produce. El proceso de elaboración siempre es el mismo (Consejo de Ministros y Parlamento), pero puede dar como resultado distintos tipos de normas --reglamentos, decisiones, resoluciones, directivas--, con distintos efectos jurídicos.
\end{formula}

\begin{table}[H]
\centering
\renewcommand{\arraystretch}{1.3}
\begin{tabularx}{\textwidth}{>{\raggedright\arraybackslash}p{0.2\textwidth}>{\raggedright\arraybackslash}X}
\toprule
\textbf{Norma} & \textbf{Efecto y alcance} \\
\midrule
Reglamento & De aplicación directa y obligatoria en todos los Estados miembros. No requiere transposición al derecho interno. Es el más «potente» en cuanto a alcance. \\
Directiva & Obliga a los Estados miembros en cuanto al resultado a alcanzar, pero les deja libertad en cuanto a la forma y los medios. Requiere ser «transpuesta» al derecho nacional. \\
Decisión & Obligatoria para sus destinatarios específicos (un Estado, una empresa o una persona). No tiene alcance general. \\
Resolución & Instrumento de naturaleza política o declarativa. No necesariamente vinculante. \\
\bottomrule
\end{tabularx}
\caption{Tipos de normas de la Unión Europea.}
\end{table}

\begin{importante}
En el Mercosur, cada órgano emite un tipo específico de norma. En la Unión Europea, el mismo proceso de elaboración puede dar lugar a distintos tipos de normas según el efecto jurídico buscado.
\end{importante}

\section{Cuadro Cornell}

\begin{cornell}
\cpregunta{¿Qué diferencia al Comité Económico y Social del Comité de las Regiones?}{El primero representa a sectores económicos y sociales (empleadores, trabajadores, sociedad civil); el segundo, a regiones y entidades locales con problemáticas geográficas específicas.}
\cpregunta{¿Qué funciones cumplen el Banco Central Europeo y el Banco Europeo de Inversiones?}{El BCE regula la política monetaria de la zona euro (emisión, tipo de cambio, tasas); el BEI presta a los Estados miembros para proyectos que favorezcan la integración europea.}
\cpregunta{¿De qué depende el tipo de norma que dicta la Unión Europea?}{Del efecto jurídico buscado, no del órgano que la dicta: el mismo proceso (Consejo y Parlamento) puede producir reglamentos, directivas, decisiones o resoluciones.}
\cresumen{Junto a sus órganos decisorios, la Unión Europea cuenta con órganos consultivos --el Comité Económico y Social y el Comité de las Regiones-- y órganos financieros --el Banco Central Europeo y el Banco Europeo de Inversiones--. Sus normas (reglamento, directiva, decisión, resolución) se distinguen por su efecto jurídico y no por el órgano que las dicta, a diferencia del Mercosur, donde cada órgano emite un tipo específico de norma.}
\end{cornell}

% ----------------------------------------------------------------
\chapter{La tensión entre el discurso y la práctica en la Unión Europea}
\chaptermark{Tensión entre discurso y práctica}
\label{cap:tension-discurso}
% ----------------------------------------------------------------

Así como en el Mercosur existe una brecha entre el discurso integracionista y la práctica de los Estados (Capítulo~\ref{cap:tensiones}), en la Unión Europea también pueden identificarse casos en los que el discurso institucional entra en tensión con decisiones concretas de los Estados miembros.

\section{La crisis migratoria mediterránea}

Existen comisiones regionales dentro de la estructura europea que atienden problemáticas específicas de determinadas zonas, como la \textbf{Comisión Regional del Mediterráneo}, que agrupa a los países con costa sobre ese mar, históricamente una vía de comercio y traslado de personas.

\begin{ejemplo}{La crisis de los cayucos y la discusión sobre migración}
España e Italia, por su condición de países peninsulares con costa mediterránea, fueron durante años los primeros receptores de embarcaciones que cruzaban de manera irregular. Se advierte una contradicción entre el discurso institucional --que declara que los migrantes son bienvenidos-- y las dificultades reales que genera su recepción, en tanto se trata de personas que necesitan salud, vivienda y educación. Resulta preferible prometer menos y cumplirlo, evaluando luego la posibilidad de avanzar más, antes que declarar principios que después no pueden sostenerse en la práctica.
\end{ejemplo}

\section{La libertad religiosa en España}

\begin{ejemplo}{Los símbolos religiosos en el espacio público}
Como ejemplo de tensión cultural dentro de la Unión, puede citarse el debate en España sobre manifestaciones religiosas en la vía pública: por un lado, se cuestionaba a ciudadanos católicos por manifestar públicamente su fe cerca de una mezquita, mientras se defendía el derecho de otras comunidades religiosas a manifestar la propia. La Convención de Derechos Humanos reconoce el derecho de toda persona a manifestar sus ideas religiosas tanto en público como en privado. Este tipo de situaciones generó además cuestionamientos sobre el uso de símbolos religiosos en oficinas públicas, en relación con la libertad de expresión y de conciencia.
\end{ejemplo}

\section{La expulsión de comunidades gitanas de Francia}

\begin{ejemplo}{Expulsión y libre tránsito en la Unión Europea}
Ciertas disposiciones francesas dispusieron la expulsión de miembros de la comunidad gitana, lo que plantea un conflicto de derechos: si una persona de origen gitano y de nacionalidad, por ejemplo, polaca, es ciudadana de la Unión Europea, tiene en principio derecho a circular y permanecer en Francia. La justificación oficial de esas expulsiones fue calificada de discriminatoria, en la medida en que la pertenencia étnica no puede constituir por sí sola prueba de la comisión de un delito. Este caso muestra cuán compleja resulta en la práctica la aplicación del principio de libre tránsito dentro de la Unión.
\end{ejemplo}

Si bien para los países del Mercosur el libre tránsito de personas resulta una experiencia relativamente natural (Capítulo~\ref{cap:mercado-comun}), para la Unión Europea acordar las pautas de libre tránsito y los derechos que corresponden a un ciudadano de un Estado miembro en otro Estado miembro fue y sigue siendo un proceso complejo, que depende del tiempo de permanencia y de requisitos específicos según cada caso.

\begin{importante}
El principio de libre tránsito resulta central para la subsistencia del proyecto europeo: si dicho principio se viera afectado de manera sustancial, la propia continuidad de la Unión Europea quedaría comprometida.
\end{importante}

Corresponde diferenciar entre cuestiones que forman parte de la agenda internacional en la que la Unión Europea participa como sujeto de derecho internacional --junto con otros actores, como el Mercosur o los Estados Unidos-- y cuestiones que constituyen política interna propia de la Unión, como ciertas políticas ambientales o de protección a la ancianidad incorporadas en sus propios tratados (Capítulo~10). Estos ejemplos ilustran una misma idea: las decisiones de política exterior y de política interna deben analizarse de manera relacionada, sin asumir compromisos declarativos que luego no puedan sostenerse en la práctica.

\section{Cuadro Cornell}

\begin{cornell}
\cpregunta{¿Qué tensión revela la crisis migratoria mediterránea?}{La distancia entre el discurso institucional que declara bienvenidos a los migrantes y las dificultades reales de recepción que enfrentan España e Italia como primeros países receptores.}
\cpregunta{¿Qué ilustra el debate español sobre símbolos religiosos?}{Una tensión cultural interna en torno a la libertad de manifestar la propia fe en el espacio público, en relación con el derecho reconocido por la Convención de Derechos Humanos.}
\cpregunta{¿Por qué la expulsión de comunidades gitanas de Francia generó controversia?}{Porque un ciudadano de la Unión, aunque pertenezca a una minoría étnica, tiene en principio derecho a circular y permanecer en cualquier Estado miembro; la pertenencia étnica no puede por sí sola justificar una expulsión.}
\cresumen{Al igual que en el Mercosur, en la Unión Europea existe una brecha entre el discurso institucional y la práctica de los Estados, visible en la crisis migratoria mediterránea, el debate sobre libertad religiosa en España y la expulsión de comunidades gitanas de Francia. Estos casos muestran que el libre tránsito y los derechos fundamentales, aunque centrales para la subsistencia del proyecto europeo, entran con frecuencia en tensión con decisiones políticas internas condicionadas por factores sociales, económicos y culturales concretos.}
\end{cornell}

% ================================================================
\part{El Tribunal de Justicia de la Unión Europea y la solución de controversias}
% ================================================================

Tanto el Mercosur como la Unión Europea crearon mecanismos propios para resolver los conflictos que surgen entre sus Estados miembros en torno a la interpretación y aplicación de su derecho. El Capítulo~\ref{cap:tpr} presentó ya la evolución general del sistema del Mercosur, desde el Protocolo de Brasilia hasta el Tribunal Permanente de Revisión. Esta parte profundiza ese mecanismo, presenta en detalle el sistema europeo --estructurado alrededor de un tribunal único e intérprete exclusivo del derecho comunitario-- y compara ambos sistemas punto por punto, para concluir con la pregunta que atraviesa toda la materia: en qué medida cada bloque garantiza efectivamente la aplicación uniforme del derecho que crea.

% ----------------------------------------------------------------
\chapter{El Tribunal de Justicia de la Unión Europea}
\chaptermark{El Tribunal de Justicia de la UE}
\label{cap:tjue}
% ----------------------------------------------------------------

\section{La función interpretativa exclusiva}

Cierta doctrina española reconoce al Tribunal de Justicia de la Unión Europea una relevancia en el ámbito legislativo, más allá de su carácter jurisdiccional, a pesar de que formalmente no tiene facultades para legislar. La razón está en su función esencial: el único que interpreta, en forma exclusiva, el derecho comunitario es el Tribunal de la Unión Europea; no hay otra posibilidad. Desde el origen del proceso de integración europea, los Estados entendieron que de nada servía dictar una norma única si cada uno la interpretaba a su manera; por eso se concentró en un solo órgano la función de fijar el sentido del derecho comunitario, de modo que ese criterio se mantenga uniforme en el tiempo y en todos los Estados.

Al ejercer esa función, el Tribunal no se limita a aplicar la letra de la norma: cuando la interpreta, en los hechos la moldea, completando sentidos que la norma dejaba abiertos. Por eso la doctrina le atribuye un cierto carácter «legiferante» --ya adelantado en el Capítulo~\ref{cap:tpr}--, aunque no sea su función formal: sus criterios de interpretación terminan operando, en la práctica, como una fuente más del derecho de la Unión.

\begin{importante}
El Tribunal de Justicia de la Unión Europea es el \textbf{intérprete exclusivo} del derecho comunitario, pero no el único que lo aplica. Los jueces nacionales de cada Estado miembro aplican el derecho comunitario en los casos concretos que se les presentan; el Tribunal interviene cuando existe duda sobre cómo debe interpretarse una norma o cuando se cuestiona su cumplimiento.
\end{importante}

\begin{ejemplo}{Conflicto entre una empresa francesa y una italiana}
Un conflicto entre una empresa francesa y una empresa italiana que se rige por derecho comunitario se resuelve ante el juez nacional competente, no directamente ante el Tribunal de la Unión Europea. Solo si ese juez tiene dudas sobre cómo interpretar la norma aplicable recurre al Tribunal.
\end{ejemplo}

\section{Vías de acceso al Tribunal}

\begin{table}[H]
\centering
\renewcommand{\arraystretch}{1.25}
\begin{tabularx}{\textwidth}{>{\raggedright\arraybackslash}p{0.27\textwidth}>{\raggedright\arraybackslash}X}
\toprule
\textbf{Vía} & \textbf{Descripción} \\
\midrule
\textbf{Cuestión prejudicial} & El juez nacional que tiene dudas sobre la interpretación de una norma comunitaria, o sobre si su propia norma interna la contradice, eleva la consulta al Tribunal. Una vez resuelta, esa interpretación es obligatoria para ese tribunal y para todos los demás. \\
\textbf{Recurso directo} & Permite a los particulares, a los Estados y a los propios órganos de la Unión acudir directamente al Tribunal para que resuelva sobre la compatibilidad entre normas o exija a un Estado que aplique el derecho comunitario que debe prevalecer sobre su derecho interno. \\
\textbf{Recurso por incumplimiento del Estado} & Se utiliza cuando un Estado no está aplicando una norma que se encuentra en vigor. \\
\bottomrule
\end{tabularx}
\caption{Vías de acceso al Tribunal de Justicia de la Unión Europea.}
\end{table}

Cuando el Tribunal indica al juez cómo debe interpretarse la norma, el juez está obligado a aplicarla así: no es una simple consulta. Esa vinculatoriedad se extiende, además: una vez fijado el criterio en un caso, todos los demás tribunales quedan obligados a interpretar la norma de esa misma manera, lo que constituye la clave de la \textbf{función unificadora} del Tribunal. Los particulares, a diferencia de lo que ocurre en el Mercosur (Capítulo~\ref{cap:mercosur-ue}), pueden acceder directamente al recurso ante el Tribunal.

\section{Composición y formaciones de juzgamiento}

El Tribunal está integrado por un juez por cada Estado miembro. No todos los jueces intervienen en todos los casos: según la relevancia y complejidad del asunto, se determina si lo resuelve una sala pequeña, una gran sala, o si el asunto se eleva al pleno.

\begin{table}[H]
\centering
\renewcommand{\arraystretch}{1.25}
\begin{tabularx}{\textwidth}{>{\raggedright\arraybackslash}p{0.27\textwidth}>{\raggedright\arraybackslash}X}
\toprule
\textbf{Formación} & \textbf{Cuándo interviene} \\
\midrule
Sala reducida (tres jueces) & Cuestiones reiteradas, con abundantes precedentes, donde no hay debate real sobre el criterio a aplicar. \\
Gran sala (aproximadamente quince jueces) & Cuestiones más complejas o novedosas. \\
Pleno & Interviene en muy pocas ocasiones: cuando el asunto reviste una relevancia institucional de primer orden, o cuando se plantea directamente cambiar un criterio de interpretación ya asentado. \\
\bottomrule
\end{tabularx}
\caption{Formaciones de juzgamiento del Tribunal de Justicia de la Unión Europea.}
\end{table}

La clasificación entre sala pequeña, gran sala y pleno no responde a categorías rígidas predeterminadas para cada tipo de asunto, sino que se decide caso por caso según la complejidad y la existencia --o no-- de precedentes consolidados. El procedimiento interno es el mismo sin importar cuántos magistrados intervienen, y las decisiones se toman por mayoría.

\begin{definicion}{Abogado general}
Quien instruye el expediente: recibe y provee las pruebas, arma y analiza el caso, y eleva a los jueces una propuesta de resolución fundada. Es un auxiliar del Tribunal.
\end{definicion}

El límite de esta función es que el abogado general \textbf{no vota}: da su opinión, pero después son los jueces quienes se sientan a discutir y deciden. Para ser juez o abogado general se deben reunir las condiciones exigidas en el propio Estado de origen para ocupar el más alto grado judicial --la capacidad para ser juez de la corte suprema de ese país--, por lo que el nivel de exigencia depende del sistema jurídico de procedencia de cada magistrado.

No hay sorteo para la designación de jueces y abogados generales: cada Estado elige a su representante según sus propios criterios internos, y ese magistrado ocupa el cargo de forma permanente hasta la renovación correspondiente. Lo que sí se sortea, caso por caso, es qué jueces y qué abogado general intervienen en cada expediente concreto. La renovación de los miembros del Tribunal se realiza \textbf{por mitades}, nunca en bloque, precisamente para preservar la continuidad de los criterios interpretativos ya asentados.

\section{Métodos de interpretación normativa}

Los jueces aplican dos criterios complementarios:

\begin{itemize}
    \item \textbf{Criterio finalista}: la interpretación debe orientarse a cumplir la finalidad que tuvo la norma cuando fue dictada. Si una lectura posible conduce a un resultado que contradice ese propósito, esa interpretación no es válida.
    \item \textbf{Criterio sistémico}: la interpretación de una norma debe ser coherente con el resto del sistema normativo del que forma parte; no puede leerse de manera aislada, en contradicción con otras disposiciones del mismo cuerpo de normas.
\end{itemize}

Ambos métodos se vinculan con la técnica general de interpretación de la ley: buscar la intención del legislador y los fundamentos que dieron origen a la norma --idea ya presentada, para los tratados de integración en general, en el Capítulo~3--, de forma que la interpretación posibilite el cumplimiento de esa finalidad. La necesidad de interpretar surge precisamente porque la letra de la norma no siempre es clara.

\section{El Tribunal General y los tribunales especializados}

La evolución institucional del sistema judicial europeo parte de un único tribunal. Con el tiempo, en la época del Tratado de Maastricht, los Estados advirtieron que ese único órgano no daba abasto para atender todos los recursos directos que le llegaban de los particulares, y crearon el \textbf{Tribunal General} --originalmente Tribunal de Primera Instancia--, que hoy recibe recursos de cualquier legitimado (particulares, Estados, órganos), mientras que el Tribunal de Justicia concentra las consultas de los jueces nacionales, de los demás órganos y de los Estados.

No existe entre ambos una relación de jerarquía: son dos tribunales de igual rango, cada uno con su propia competencia, salvo que el Tribunal de Justicia interviene específicamente cuando el conflicto se plantea entre dos Estados. El número de jueces de cada tribunal es independiente, aunque coincidente (uno por Estado), por lo que cada Estado termina designando dos magistrados: uno para cada tribunal.

\begin{ejemplo}{El Tribunal de Asuntos Administrativos}
Durante un tiempo funcionó el Tribunal de Asuntos Administrativos, creado para resolver conflictos entre los órganos de la Unión Europea. Finalmente fue disuelto, posiblemente porque el costo institucional de mantenerlo no se justificaba, ya que el Tribunal General podía resolver esas mismas cuestiones sin necesidad de sostener una estructura adicional con representación de todos los Estados miembros.
\end{ejemplo}

Cuando se habla del «Tribunal de la Unión Europea», en realidad se está hablando de algo más amplio que una corte de unos pocos jueces: es todo un sistema de solución de conflictos.

\section{Control de legalidad}

El Tribunal ejerce además un control de legalidad entre los propios órganos de la Unión: puede intervenir cuando un órgano sostiene que una norma dictada por otro órgano resulta incompatible con otra norma superior.

\begin{definicion}{Control de legalidad}
Aplicado al derecho de la integración, implica verificar: (1) que la norma haya sido emitida por el órgano competente; (2) que se haya seguido el procedimiento establecido para dictar ese tipo de norma; y (3) que no vulnere ninguna norma superior del ordenamiento.
\end{definicion}

\section{Cuadro Cornell}

\begin{cornell}
\cpregunta{¿Qué función cumple, en exclusiva, el Tribunal de Justicia de la Unión Europea?}{Es el único intérprete del derecho comunitario. No es el único que lo aplica: la aplicación en casos concretos corresponde a los jueces nacionales, que recurren al Tribunal por vía de cuestión prejudicial cuando tienen dudas sobre la interpretación.}
\cpregunta{¿Cuáles son las tres vías de acceso al Tribunal?}{La cuestión prejudicial (consulta de un juez nacional, con efecto obligatorio y general), el recurso directo (particulares, Estados y órganos) y el recurso por incumplimiento del Estado.}
\cpregunta{¿Qué es el abogado general y qué no puede hacer?}{Instruye el expediente y propone una resolución a los jueces, pero no vota ni participa de la deliberación.}
\cpregunta{¿Qué diferencia hay entre el Tribunal de Justicia y el Tribunal General?}{No hay jerarquía entre ambos. El Tribunal General recibe recursos de cualquier legitimado; el Tribunal de Justicia concentra las cuestiones prejudiciales, las consultas de otros órganos y los conflictos entre Estados.}
\cresumen{El Tribunal de Justicia de la Unión Europea es el intérprete exclusivo del derecho comunitario, con tres vías de acceso (cuestión prejudicial, recurso directo, recurso por incumplimiento), salas de distinto tamaño según la complejidad del asunto, abogados generales que instruyen sin votar, dos criterios de interpretación (finalista y sistémico), un Tribunal General de igual rango para los recursos directos, y una función de control de legalidad entre los propios órganos de la Unión.}
\end{cornell}

% ----------------------------------------------------------------
\chapter{Profundización del sistema de solución de controversias del Mercosur}
\chaptermark{Solución de controversias del Mercosur (profundización)}
\label{cap:solucion-mercosur}
% ----------------------------------------------------------------

El Capítulo~\ref{cap:tpr} presentó la evolución general del sistema de solución de controversias del Mercosur, desde el Protocolo de Brasilia hasta el Tribunal Permanente de Revisión. Este capítulo profundiza ese recorrido con los elementos necesarios para compararlo, punto por punto, con el sistema europeo recién descrito.

\section{Autocomposición y heterocomposición}

\begin{table}[H]
\centering
\renewcommand{\arraystretch}{1.3}
\begin{tabularx}{\textwidth}{>{\raggedright\arraybackslash}p{0.2\textwidth}>{\raggedright\arraybackslash}X>{\raggedright\arraybackslash}X}
\toprule
\textbf{Concepto} & \textbf{Definición} & \textbf{Ejemplos} \\
\midrule
Autocomposición & Son las propias partes las que resuelven el conflicto, llegando a un acuerdo sobre cómo solucionarlo. & Negociación directa, conciliación, mediación \\
Heterocomposición & Interviene un tercero imparcial e independiente de las partes, que resuelve el conflicto. & Arbitraje, tribunales judiciales o permanentes \\
\bottomrule
\end{tabularx}
\caption{Autocomposición y heterocomposición.}
\end{table}

Dentro del Mercosur, el sistema combina ambas lógicas: primero instancias de autocomposición (negociaciones directas y, en cierta medida, la intervención conciliadora del Grupo Mercado Común) y, si estas fracasan, instancias de heterocomposición (arbitraje y, luego, el Tribunal Permanente de Revisión).

\section{El Protocolo de Brasilia y los ajustes posteriores}

El Tratado de Asunción no reguló ningún procedimiento de solución de controversias; ese vacío fue cubierto por el \textbf{Protocolo de Brasilia} (1991), que previó tres etapas: negociaciones directas (plazo de quince días, prorrogable), intervención \textbf{obligatoria} del Grupo Mercado Común como mediador, y arbitraje ante un Tribunal Arbitral Ad Hoc, integrado por un mínimo de tres árbitros --uno por cada Estado en conflicto y un tercero neutral que preside--, surgidos de nóminas de expertos que cada Estado presenta al momento de crearse el bloque.

El \textbf{Protocolo de Ouro Preto} mantuvo en lo sustancial este esquema, con una modificación relevante: quitó el carácter obligatorio de la intervención del Grupo Mercado Común, que pasó a ser facultativa. Con el \textbf{Protocolo de Olivos} (firmado el 18 de febrero de 2002, vigente desde el 1.º de enero de 2004), que reemplazó a Brasilia y creó el Tribunal Permanente de Revisión, todo el procedimiento pasó a ser alternativo: las partes pueden optar por omitir etapas intermedias y avanzar directamente hacia el arbitraje.

\begin{importante}
El arbitraje es una forma de justicia privada, no estatal ni nacional. Al no ser un tribunal permanente sino ad hoc en su primera instancia, los laudos son obligatorios entre las partes de cada caso, pero no generan jurisprudencia vinculante para casos futuros: distintos tribunales, integrados por distintos árbitros, pueden interpretar la misma norma de manera diferente.
\end{importante}

\section{Las tres funciones del Tribunal Permanente de Revisión}

Además de revisar el laudo arbitral --su función original, presentada en el Capítulo~\ref{cap:tpr}--, el Tribunal Permanente de Revisión cumple otras dos funciones:

\begin{itemize}
    \item \textbf{Instancia única \emph{per saltum}}: cuando, concluidas las negociaciones directas sin acuerdo, las partes deciden ir directamente al Tribunal sin pasar por el arbitraje ad hoc, generalmente por la urgencia o la importancia del asunto. En este supuesto, el Tribunal actúa con la totalidad de las facultades de un tribunal de instancia única: puede producir prueba, dictar medidas y sentenciar.
    \item \textbf{Medidas provisionales o cautelares}: puede dictarlas en cualquier momento del procedimiento, incluso durante la negociación directa, cuando exista riesgo de un daño grave o irreparable.
\end{itemize}

\begin{ejemplo}{Medida cautelar ante un daño irreparable}
Ante la importación de bananas o de algún producto perecedero, si ninguna de las partes toma medidas para la conservación del material, es posible acudir directamente al Tribunal Permanente de Revisión para que adopte una medida urgente de protección.
\end{ejemplo}

\section{El objeto del mecanismo y la posibilidad de acudir a otros foros}

El mecanismo de solución de controversias del Mercosur --al igual que el de la Unión Europea-- apunta a resolver conflictos relativos a la aplicación, interpretación y vigencia del derecho de la integración, y no a resolver el conflicto particular entre dos partes privadas. Cuando un Estado interpreta y aplica la norma de un modo distinto al fijado por el órgano de solución de controversias, ese Estado debe corregir su normativa interna.

A diferencia del Protocolo de Brasilia, que establecía el mecanismo del Mercosur como la única vía a seguir, el Protocolo de Olivos permite que, de común acuerdo, los Estados sometan la controversia a otro mecanismo distinto, si estuvieran vinculados a él por otro convenio.

\begin{ejemplo}{El caso de «las papeleras»}
El conflicto conocido como «el caso de las papeleras», entre Argentina y Uruguay, vinculado a la instalación de plantas de celulosa sobre el río Uruguay, no se resolvió mediante el sistema de solución de controversias del Mercosur, sino ante otro tribunal internacional. Esta circunstancia se señala como una incoherencia respecto de los objetivos del bloque: el conflicto se resolvió atendiendo más a consideraciones comerciales que a la norma o a su ubicación dentro del sistema del Mercosur.
\end{ejemplo}

\begin{importante}
Permitir litigar fuera del sistema del Mercosur le resta autoridad y sentido al propio Tribunal Permanente de Revisión, dado que el objetivo declarado del mecanismo es justamente uniformar la interpretación y aplicación del derecho de la integración. El sistema es \textbf{opcional}: rige siempre que los Estados decidan impulsar la acción del tribunal, pero no existe sometimiento obligatorio.
\end{importante}

\section{Ejecución de los laudos y medidas compensatorias}

El sistema del Mercosur carece de supranacionalidad: las decisiones del Tribunal Permanente de Revisión y los laudos del Tribunal Ad Hoc no son ejecutables por la fuerza frente a un Estado que decide no cumplirlos. Si el Estado condenado no acepta cumplir, el Estado perjudicado queda habilitado para tomar \textbf{medidas compensatorias}, en principio dentro del mismo rubro en el que se produjo el perjuicio, y solo excepcionalmente sobre otro sector. El Estado perjudicado no puede fijar por sí mismo, en forma definitiva, el alcance de la medida: debe someter su propuesta a la aprobación del mismo órgano que dictó el laudo o la sentencia.

\section{Legitimación para accionar: Estados y particulares}

Únicamente los Estados tienen legitimación directa para accionar ante el sistema. Una empresa perjudicada no puede presentarse directamente: debe formular su reclamo ante la \textbf{Sección Nacional del Grupo Mercado Común} de su propio Estado, aportando toda la documentación respaldatoria. La Sección Nacional evalúa el caso y, si considera que existe un incumplimiento, puede decidir --de manera discrecional, con consideraciones también políticas y económicas-- asumirlo como un conflicto interestatal.

\begin{ejemplo}{Contraste entre una gran empresa y un monotributista}
Resulta poco probable que un Estado decida asumir un conflicto diplomático con otro Estado del bloque por un reclamo de pequeña escala --como el de un monotributista al que se le cobró un impuesto que entiende improcedente--, mientras que sí es más probable que lo haga frente al reclamo de una empresa de mayor envergadura.
\end{ejemplo}

\section{Cuadro Cornell}

\begin{cornell}
\cpregunta{¿Qué cambió el Protocolo de Ouro Preto respecto de Brasilia en materia de solución de controversias?}{Quitó el carácter obligatorio de la intervención del Grupo Mercado Común, que pasó a ser facultativa.}
\cpregunta{¿Cuáles son las tres funciones del Tribunal Permanente de Revisión?}{Revisar el laudo arbitral (apelación), actuar como instancia única \emph{per saltum}, y dictar medidas provisionales o cautelares.}
\cpregunta{¿Por qué se considera incoherente que los Estados puedan acudir a otros foros?}{Porque el Protocolo de Olivos lo permite de común acuerdo (caso de las papeleras), lo que debilita el objetivo declarado de uniformar la interpretación del derecho del bloque.}
\cpregunta{¿Cómo se ejecuta un laudo incumplido en el Mercosur?}{No es ejecutable por la fuerza; el Estado perjudicado puede aplicar medidas compensatorias, aprobadas por el mismo órgano que dictó la decisión incumplida.}
\cpregunta{¿Cómo accede un particular al sistema de solución de controversias del Mercosur?}{No accede directamente: debe presentar su reclamo ante la Sección Nacional del Grupo Mercado Común de su Estado, que decide de forma discrecional si lo asume como conflicto interestatal.}
\cresumen{El sistema de solución de controversias del Mercosur evolucionó del esquema obligatorio de tres etapas del Protocolo de Brasilia, a la flexibilización de Ouro Preto, hasta el Tribunal Permanente de Revisión creado por el Protocolo de Olivos, con funciones de apelación, instancia única \emph{per saltum} y medidas cautelares. El sistema carece de supranacionalidad --los laudos se hacen cumplir mediante medidas compensatorias, no por ejecución forzosa--, es opcional frente a otros foros, y los particulares no tienen acceso directo, sino solo a través de la Sección Nacional de su Estado.}
\end{cornell}

% ----------------------------------------------------------------
\chapter{Mercosur y Unión Europea: dos sistemas en espejo}
\chaptermark{Mercosur y Unión Europea en espejo}
\label{cap:mercosur-ue}
% ----------------------------------------------------------------

\section{La opinión consultiva ante el Tribunal Permanente de Revisión}

El punto de comparación más importante entre ambos sistemas es quién puede acceder al tribunal respectivo y con qué efecto. En el Mercosur no existe recurso directo para los particulares: solo la \textbf{Corte Suprema}, máximo tribunal de cada Estado parte, está habilitada para elevar una consulta al Tribunal Permanente de Revisión; los tribunales inferiores no pueden hacerlo directamente y deben canalizar la cuestión a través de la Corte. En la Unión Europea, en cambio, cualquier tribunal puede plantear la cuestión prejudicial.

\begin{importante}
\textbf{Consulta no vinculante.} A diferencia de la cuestión prejudicial europea, la opinión consultiva del Tribunal Permanente de Revisión no obliga a la Corte que la solicitó: es una consulta que, después, la Corte puede o no aplicar. La Corte puede apoyarse en esa opinión, pero también puede apartarse de ella si considera que existen motivos para interpretar la norma del Mercosur de otra manera dentro de su propio ordenamiento.
\end{importante}

Esta característica se conecta con un rasgo estructural del sistema argentino: al regir entre nosotros el \textbf{control de constitucionalidad difuso}, cualquier tribunal puede interpretar por sí mismo una norma del Mercosur, sin que esa interpretación sea competencia exclusiva de los órganos del bloque. Este diseño revela un escaso compromiso institucional: se crea una estructura común, pero cada Estado conserva, en los hechos, la última palabra.

\begin{ejemplo}{La consulta que se demora a propósito}
Si un Estado remite la consulta a su Corte Suprema sabiendo que la respuesta puede demorar mucho tiempo, en los hechos logra postergar indefinidamente la aplicación del criterio del Tribunal Permanente de Revisión y, mientras tanto, sigue aplicando su propio criterio. Con esta imagen se ilustra la diferencia entre un sistema que garantiza de raíz la aplicación uniforme del derecho --la Unión Europea, con un único tribunal y varias vías de acceso-- y uno que no la garantiza del mismo modo --el Mercosur--.
\end{ejemplo}

\section{Cuadro comparativo entre ambos sistemas}

Ambos mecanismos de solución de controversias fueron creados con el mismo objetivo declarado: garantizar la armonización legislativa y la aplicación uniforme del derecho que cada bloque crea.

\begin{table}[H]
\centering
\renewcommand{\arraystretch}{1.25}
\begin{tabularx}{\textwidth}{>{\raggedright\arraybackslash}p{0.19\textwidth}>{\raggedright\arraybackslash}X>{\raggedright\arraybackslash}X}
\toprule
\textbf{Criterio} & \textbf{Mercosur} & \textbf{Unión Europea} \\
\midrule
Naturaleza de los órganos & El arbitraje inicial es ad hoc, se conforma para cada controversia y luego se disuelve; el Tribunal Permanente de Revisión, en cambio, es un órgano permanente. & Órganos permanentes, de naturaleza judicial, con jueces que rotan por mitades. \\
Acceso al tribunal & Solo la Corte Suprema de cada Estado (los tribunales inferiores no acceden directamente); los particulares no tienen recurso directo. & Acceso amplio: jueces nacionales (cuestión prejudicial), Estados, órganos y particulares (recurso directo). \\
Efecto de la opinión consultiva & No es vinculante para la Corte que la solicitó. & La interpretación fijada por el Tribunal es obligatoria para el tribunal consultante y para todos los demás. \\
Cumplimiento del laudo arbitral & El laudo arbitral sí es vinculante; su incumplimiento habilita medidas compensatorias. & Las resoluciones del Tribunal son siempre obligatorias. \\
Tipo de reclamo posible ante el bloque & Recursos limitados: no existe la vía directa amplia que sí existe en la Unión Europea. & Múltiples vías: cuestión prejudicial, recurso directo, recurso por incumplimiento del Estado. \\
\bottomrule
\end{tabularx}
\caption{Mecanismos de solución de controversias: Mercosur y Unión Europea.}
\end{table}

La conclusión, anticipada desde la Parte~IV de este libro, es que el Mercosur se está comportando, en las relaciones entre sus Estados, como si fueran relaciones internacionales, en lugar de ser un espacio de integración plena.

\section{Cuadro Cornell}

\begin{cornell}
\cpregunta{¿Quién puede elevar una consulta al Tribunal Permanente de Revisión y a quién obliga la respuesta?}{Solo la Corte Suprema de cada Estado parte, y la respuesta --opinión consultiva-- no es vinculante para ella.}
\cpregunta{¿Qué distingue a la cuestión prejudicial europea de la opinión consultiva del Mercosur?}{La cuestión prejudicial obliga al tribunal consultante y a todos los demás; la opinión consultiva no obliga a la Corte que la solicitó, que puede apartarse de ella.}
\cpregunta{¿Qué revela el diseño institucional del Mercosur sobre el compromiso de sus Estados?}{Que, pese a crear una estructura común, cada Estado conserva en los hechos la última palabra, a diferencia de la Unión Europea, donde las resoluciones del Tribunal son siempre obligatorias.}
\cpregunta{¿Qué conclusión general se extrae de la comparación entre ambos sistemas?}{Que el Mercosur funciona, en la práctica, más cerca de una relación internacional clásica que de un verdadero espacio de integración con aplicación uniforme de su derecho.}
\cresumen{La comparación entre ambos sistemas de solución de controversias confirma una diferencia estructural que atraviesa toda la materia: la Unión Europea garantiza la aplicación uniforme de su derecho mediante un tribunal permanente, de acceso amplio y de interpretación obligatoria; el Mercosur, en cambio, limita el acceso a la Corte Suprema de cada Estado y no vincula a esa Corte con la opinión del Tribunal Permanente de Revisión, lo que explica por qué se sostiene que, en la práctica, el Mercosur se comporta más como una relación internacional clásica que como un espacio de integración plenamente consolidado.}
\end{cornell}

% ================================================================
\part{Inmunidad del Estado y de los órganos de integración}
% ================================================================

Todo lo desarrollado hasta aquí presupone una pregunta previa: ¿quién puede ser llevado a juicio, ante qué tribunal, y con qué posibilidad real de hacer cumplir lo decidido? Esta parte aborda esa pregunta en dos planos: la inmunidad del Estado frente a la jurisdicción de otro Estado, y la inmunidad de los órganos y funcionarios de un área de integración frente a la jurisdicción del Estado que los aloja.

% ----------------------------------------------------------------
\chapter{Inmunidad del Estado: de la inmunidad absoluta a la relativa}
\chaptermark{Inmunidad del Estado}
% ----------------------------------------------------------------

\section{El principio \emph{par in parem non habet imperium}}

\begin{definicion}{Par in parem non habet imperium}
«Entre iguales no hay autoridad.» Principio formulado por Bartolo de Sassoferrato en el siglo XI, en el marco de la primera escuela estatutaria italiana. Significa que si dos Estados son soberanos e iguales entre sí, ninguno puede quedar sometido a la jurisdicción del otro.
\end{definicion}

Este principio se prolongó y reforzó durante los siglos XV y XVI con el surgimiento de las monarquías absolutas, período en el cual la figura del rey se identificaba con el propio Estado: el soberano no podía ser cuestionado ni sometido a ningún tribunal, fuera este nacional o extranjero, porque cualquier acto emanado del monarca era, por definición, un acto del Estado.

\section{De la inmunidad absoluta a la inmunidad relativa}

Con el transcurso del tiempo, la concepción absoluta de la inmunidad fue cediendo ante nuevas realidades jurídicas y económicas. El primer quiebre conceptual se produjo cuando se aceptó que un Estado podía quedar sometido a las decisiones de una autoridad extranjera si él mismo había iniciado la acción ante esa jurisdicción --sumisión voluntaria: quien elige un foro para demandar no puede desconocer luego lo que ese foro resuelva--. El paso siguiente fue más profundo: el reconocimiento de que el Estado no actúa siempre en ejercicio de su soberanía, sino que en ocasiones opera como un agente comercial, con un carácter idéntico al de un particular. Esta dualidad de roles es la base de la \textbf{inmunidad relativa}.

\begin{table}[H]
\centering
\renewcommand{\arraystretch}{1.3}
\begin{tabularx}{\textwidth}{>{\raggedright\arraybackslash}p{0.22\textwidth}>{\raggedright\arraybackslash}X>{\raggedright\arraybackslash}X}
\toprule
\textbf{Aspecto} & \textbf{Inmunidad absoluta} & \textbf{Inmunidad relativa} \\
\midrule
Período histórico & Hasta fines del siglo XIX / principios del XX & Desde mediados del siglo XX hasta la actualidad \\
Fundamento & \emph{Par in parem non habet imperium}. Soberanía total & El Estado puede actuar con doble carácter: soberano o comercial \\
Distinción de actos & No se realizaba distinción alguna: todos los actos del Estado gozaban de inmunidad & Se distingue entre actos de imperio (soberanos) y actos de gestión (comerciales) \\
Consecuencia & El Estado nunca podía ser demandado ante tribunales extranjeros & El Estado puede ser demandado por sus actos de gestión; los actos de imperio mantienen la inmunidad \\
Postura argentina & Argentina fue más lenta en adoptar la distinción & Reconocida jurisprudencialmente a partir del fallo \emph{Manauta} (1994), Capítulo~\ref{cap:manauta} \\
\bottomrule
\end{tabularx}
\caption{Inmunidad absoluta e inmunidad relativa.}
\end{table}

\section{Actos de imperio y actos de gestión}

\begin{definicion}{Acto de imperio}
Acto directamente vinculado a la soberanía del Estado, que solo puede realizar el Estado en cuanto tal, pues involucra el ejercicio de sus funciones esenciales. Goza de inmunidad de jurisdicción.
\end{definicion}

\begin{definicion}{Acto de gestión}
Acto que el Estado realiza sin comprometer su soberanía ni su existencia, equivalente a los que podría realizar un particular --contratos comerciales, contratos laborales, emisión de deuda--. No goza de inmunidad.
\end{definicion}

\begin{formula}{Test de distinción}
La pregunta clave es: ¿podría un particular realizar este acto? Si la respuesta es negativa, se trata de un acto de imperio; si es afirmativa, se trata de un acto de gestión.
\end{formula}

\begin{ejemplo}{Discurso presidencial sobre recursos naturales en territorio en disputa}
Un discurso presidencial que anuncia que se sancionará a toda empresa que explote recursos naturales en territorio que el Estado considera propio, sin su autorización --idea ya presentada en el Capítulo~\ref{cap:tpr} respecto de la cuestión Malvinas--, constituye un \textbf{acto de imperio}: está directamente vinculado al ejercicio de la soberanía territorial.
\end{ejemplo}

\begin{ejemplo}{La compra de gas natural a Bolivia}
Cuando el Estado argentino le compraba gas natural al Estado boliviano para suplir el déficit de provisión interna --caso ya presentado en el Capítulo~13 al distinguir integración de derecho internacional común--, ese acto era un \textbf{acto de gestión}: aunque intervenían dos Estados soberanos, la operación consistía en la compraventa de una mercadería, actividad que también puede realizar un particular. Que ambas partes sean Estados no transforma por sí solo al acto en uno de carácter soberano.
\end{ejemplo}

\begin{ejemplo}{Aerolíneas Argentinas, YPF y la deuda pública}
La explotación de Aerolíneas Argentinas o de YPF, siendo empresas con participación estatal, constituye un \textbf{acto de gestión}, porque operan como personas de derecho privado; en cambio, si el Estado adquiriera aviones destinados específicamente a la defensa nacional y ese contrato no se cumpliera, la operación estaría intrínsecamente vinculada a la defensa de la soberanía y sería un \textbf{acto de imperio}. La emisión de deuda pública mediante bonos en mercados internacionales es también un acto de gestión: si el Estado no cumple con el pago, puede ser demandado por los acreedores.
\end{ejemplo}

\begin{ejemplo}{El Banco Central como garantía de un empréstito internacional}
En un caso real, un presidente de la Nación puso como garantía de un préstamo internacional los bienes del Banco Central de la República Argentina. Al ofrecer ese bien como garantía en un acto comercial, el Estado renunció a la inmunidad sobre ese bien específico: aunque los bienes vinculados al ejercicio de la soberanía son normalmente inembargables, la voluntad expresa del Estado de ofrecerlos como garantía levanta esa protección --el mismo mecanismo que opera cuando una persona levanta el bien de familia para darlo en garantía de un crédito bancario--.
\end{ejemplo}

\begin{importante}
Los actos de imperio continúan gozando de inmunidad en el plano internacional. Ningún organismo internacional puede imponerle a un Estado soberano una decisión sobre materias que hacen al núcleo de su soberanía, salvo que ese Estado se haya sometido voluntariamente.
\end{importante}

\section{Cuadro Cornell}

\begin{cornell}
\cpregunta{¿Qué significa el principio par in parem non habet imperium?}{Que entre Estados soberanos e iguales ninguno puede quedar sometido a la jurisdicción del otro; fundamento histórico de la inmunidad absoluta.}
\cpregunta{¿Qué distingue a la inmunidad absoluta de la relativa?}{La absoluta no distinguía entre tipos de actos: todo acto del Estado era inmune. La relativa distingue actos de imperio (inmunes) de actos de gestión (no inmunes), reconociendo que el Estado puede actuar como un agente comercial.}
\cpregunta{¿Cuál es el test para distinguir un acto de imperio de un acto de gestión?}{Preguntarse si un particular podría realizar ese mismo acto: si no puede, es un acto de imperio; si puede, es un acto de gestión.}
\cpregunta{¿Cómo puede un Estado renunciar a la inmunidad sobre un bien en particular?}{Ofreciéndolo expresamente como garantía en un acto comercial, como ocurrió con bienes del Banco Central puestos como garantía de un préstamo internacional.}
\cresumen{La inmunidad del Estado evolucionó del principio par in parem non habet imperium --que sostenía una inmunidad absoluta-- hacia una inmunidad relativa, que distingue entre actos de imperio (inmunes, exclusivos de la soberanía) y actos de gestión (no inmunes, equivalentes a los de un particular). Un Estado puede además renunciar expresamente a la inmunidad sobre un bien determinado al ofrecerlo como garantía en un acto comercial.}
\end{cornell}

% ----------------------------------------------------------------
\chapter{Mecanismos de solución de conflictos entre Estados}
\chaptermark{Solución de conflictos entre Estados}
% ----------------------------------------------------------------

Cuando un Estado no cumple con una norma del derecho de la integración --por ejemplo, si un acuerdo regional exime de aranceles a ciertos productos y un Estado miembro los cobra de todos modos--, se plantea la cuestión de cómo exigirle a un Estado soberano que acate esa norma. Como ya se presentó en el Capítulo~\ref{cap:solucion-mercosur}, estos mecanismos se agrupan en dos grandes familias: la autocomposición, en la que las propias partes resuelven el conflicto, y la heterocomposición, en la que interviene un tercero.

\begin{importante}
En el ámbito internacional, la distinción entre autocomposición y heterocomposición adquiere una particularidad relevante: dado que los Estados no están obligados a someterse a ningún mecanismo externo de resolución, cuando un Estado acepta voluntariamente la intervención de un mediador o árbitro, el compromiso moral y político de cumplir con el resultado es tan fuerte que la mediación funciona en la práctica con efectos similares a los de la heterocomposición vinculante. Un Estado que desobedeciera la propuesta de un mediador que él mismo eligió enfrentaría consecuencias diplomáticas de envergadura.
\end{importante}

\begin{definicion}{Mediación}
Mecanismo de heterocomposición en el que un tercero neutral propone una solución al conflicto. Las partes no están jurídicamente obligadas a aceptarla, aunque en el plano internacional el compromiso político suele ser determinante para su cumplimiento.
\end{definicion}

\begin{definicion}{Conciliación}
Mecanismo de heterocomposición en el que un tercero interviene para acercar a las partes y facilitar un acuerdo, sin imponer una solución.
\end{definicion}

A diferencia de estas dos, el \textbf{arbitraje} (ya desarrollado para el Mercosur en los Capítulos~\ref{cap:tpr} y~\ref{cap:solucion-mercosur}) dicta un laudo susceptible de cumplimiento forzoso, con fuerza vinculante para las partes.

\section{Canal de Beagle: la mediación papal}

\begin{ejemplo}{La mediación del Papa Juan Pablo II}
El conflicto entre Argentina y Chile por la soberanía sobre islas del Canal de Beagle involucraba el acceso al Océano Atlántico: la posesión de esas islas otorgaba al Estado titular una salida al Atlántico, evitando la circunnavegación por el Pacífico, el océano más extenso del mundo, con consecuencias comerciales y militares de primer orden. Argentina y Chile acordaron voluntariamente designar como mediador al Papa Juan Pablo II, por la confianza que ambos Estados depositaban en la imparcialidad de esa figura. El resultado fue una propuesta de mediación que asignó determinados territorios e islas y otorgó a Chile acceso al canal; ese acuerdo constituye el fundamento de la situación territorial vigente en la zona.
\end{ejemplo}

\begin{nota}
La propuesta papal no era un laudo susceptible de ejecución forzosa, sino una propuesta que los Estados podían aceptar o rechazar. Dado que ambos países la aceptaron voluntariamente, el cumplimiento fue pleno: en materia de actos soberanos, la aceptación o rechazo de una propuesta de mediación constituye en sí misma un ejercicio de la soberanía del Estado.
\end{nota}

\section{Límites cordilleranos: Perito Moreno y el arbitraje británico}

El conflicto limítrofe entre Argentina y Chile en la Cordillera de los Andes se dirimió mediante arbitraje, con Gran Bretaña como árbitro, ante la ambigüedad de los dos criterios internacionales para trazar fronteras en zonas montañosas: el criterio de los altos picos (una línea imaginaria entre los picos más altos de cada bloque montañoso) y el criterio del curso de las aguas (según hacia qué lado drena el agua del deshielo glaciar).

\begin{ejemplo}{La estrategia de Perito Moreno}
Francisco Pascasio Moreno (\emph{Perito Moreno}), naturalista y director del Museo de Ciencias Naturales de La Plata, fue designado para representar los intereses argentinos. Recorrió exhaustivamente la zona en disputa, elaboró cartografía detallada e identificó los sectores en los que, aplicando el criterio del drenaje de aguas, el árbitro inglés probablemente fallaría a favor de Chile. En esas zonas estratégicas, el gobierno argentino instaló destacamentos policiales o militares, una oficina de correos y una escuela, ejerciendo soberanía efectiva sobre el territorio: el ejercicio efectivo y continuado de soberanía constituye, en el derecho internacional, uno de los fundamentos más sólidos para reclamar la pertenencia de un territorio. Gracias a esta estrategia, Argentina conservó sectores patagónicos que de otro modo habría perdido; como reconocimiento, el Estado le donó tierras a Perito Moreno, y él las donó al propio Estado para que se constituyera un parque nacional.
\end{ejemplo}

\section{Cuadro Cornell}

\begin{cornell}
\cpregunta{¿Por qué la mediación internacional tiende a cumplirse aunque no sea jurídicamente vinculante?}{Porque al ser aceptada voluntariamente por el Estado, su incumplimiento generaría consecuencias diplomáticas de gran envergadura, lo que la acerca en la práctica a la heterocomposición vinculante.}
\cpregunta{¿Qué resolvió la mediación papal en el conflicto del Canal de Beagle?}{Asignó territorios e islas entre Argentina y Chile y otorgó a Chile acceso al canal, en una propuesta aceptada voluntariamente por ambos países.}
\cpregunta{¿Qué estrategia siguió Perito Moreno en el arbitraje limítrofe con Chile?}{Identificó las zonas donde Argentina perdería según el criterio del drenaje de aguas y promovió allí el ejercicio efectivo de soberanía (destacamentos, correo, escuela), fundamento sólido para reclamar el territorio ante el árbitro británico.}
\cresumen{Los mecanismos de solución de conflictos entre Estados --autocomposición y heterocomposición-- adquieren en el plano internacional una fuerza práctica que trasciende su carácter jurídicamente vinculante, como ilustran la mediación papal en el Canal de Beagle y el arbitraje británico en los límites cordilleranos con Chile, donde el ejercicio efectivo de soberanía, promovido estratégicamente por Perito Moreno, resultó decisivo.}
\end{cornell}

% ----------------------------------------------------------------
\chapter{Inmunidad de ejecución y el caso Manauta}
\chaptermark{Inmunidad de ejecución y caso Manauta}
\label{cap:manauta}
% ----------------------------------------------------------------

\section{Inmunidad de jurisdicción e inmunidad de ejecución}

La caída de la inmunidad de jurisdicción --es decir, la posibilidad de llevar a un Estado a juicio-- no implica necesariamente la caída de la inmunidad de ejecución. Son dos inmunidades conceptualmente distintas.

\begin{definicion}{Inmunidad de jurisdicción}
Prerrogativa del Estado extranjero de no ser sometido a los tribunales de otro Estado. Cuando se pierde, el Estado puede ser llevado a juicio.
\end{definicion}

\begin{definicion}{Inmunidad de ejecución}
Prerrogativa del Estado extranjero de que sus bienes no sean embargados ni ejecutados forzosamente, incluso si fue condenado en juicio.
\end{definicion}

\begin{importante}
Distintos países adoptan posturas diferentes: algunos consideran que la caída de la inmunidad de jurisdicción arrastra necesariamente la de ejecución; otros --como la Argentina, según la Corte Suprema-- los tratan como institutos independientes. Una sentencia condenatoria que no puede ejecutarse se convierte, en los hechos, en una norma de cumplimiento voluntario.
\end{importante}

No todos los bienes del Estado son susceptibles de embargo: los \textbf{bienes no ejecutables} son los vinculados al ejercicio de la soberanía (edificios de gobierno, bienes militares, reservas del Banco Central), que mantienen su protección aunque el Estado haya perdido la inmunidad de jurisdicción; los \textbf{bienes ejecutables} son los vinculados a actos de gestión, siempre que el Estado haya renunciado expresamente a la inmunidad sobre ellos.

\begin{ejemplo}{El caso de la Fragata Libertad}
La fragata ARA Libertad de la Armada Argentina fue retenida en Ghana por un acreedor de bonos de deuda soberana argentina, tras una escala por desperfectos técnicos que requirió reparaciones en un astillero local. Quien prestó el servicio invocó su derecho de retención sobre el bien hasta cobrar el trabajo realizado, pero el fundamento de fondo era más complejo: Argentina había renunciado previamente a la inmunidad de ese bien en un contrato vinculado a la deuda soberana. El derecho de retención fue el instrumento formal; la causa profunda fue esa renuncia previa.
\end{ejemplo}

\section{El caso Manauta c/ Embajada de la Federación Rusa}

\begin{definicion}{Caso Manauta}
Fallo de la Corte Suprema de Justicia de la Nación argentina que constituye el hito jurisprudencial que consolidó en Argentina el tránsito de la inmunidad absoluta a la inmunidad relativa de los Estados extranjeros.
\end{definicion}

La Embajada de la Unión de Repúblicas Socialistas Soviéticas (URSS) en Argentina tenía empleados contratados localmente en su sección de prensa. Cuando se produjo el despido de ese sector, los trabajadores reclamaron salarios caídos, indemnización y aportes previsionales no ingresados. La demanda se interpuso contra la \textbf{Federación Rusa}, y no contra la URSS: los empleados habían sido contratados por la URSS, que al momento del juicio ya no existía como Estado, disuelta tras la caída del Muro de Berlín. La primera defensa de Rusia fue la falta de legitimación pasiva --los empleados habían trabajado para la URSS, no para Rusia--, y el fallo desarrolla un análisis político-jurídico para justificar por qué Rusia sí era continuadora de la URSS a esos efectos.

\begin{normacitada}{Ley 1285 (1958)}
Norma argentina vigente al momento del caso, que establecía que los Estados extranjeros gozaban de inmunidad de jurisdicción, salvo que aceptaran voluntariamente someterse a los tribunales argentinos. El silencio ante la notificación de una demanda no podía interpretarse como aceptación; esta debía ser expresa.
\end{normacitada}

Bajo este régimen, la embajada rusa fue notificada reiteradamente y guardó silencio; los tribunales de primera instancia y la Cámara de Apelaciones declararon la incompetencia de los jueces argentinos. El caso llegó a la Corte Suprema, que resolvió adoptar la \textbf{inmunidad relativa}: en el marco de las relaciones internacionales contemporáneas, los Estados no gozan de inmunidad absoluta, sino que es necesario distinguir entre actos de imperio y actos de gestión. El contrato laboral entre la embajada y sus empleados fue calificado como un \textbf{acto de gestión}, por lo que la Federación Rusa podía ser sometida a la jurisdicción argentina.

A partir del fallo Manauta se sancionó la \textbf{Ley 24.488}, que codificó el criterio de inmunidad relativa adoptado por la Corte Suprema, estableciendo legalmente la distinción entre actos de imperio y actos de gestión en materia de inmunidad de jurisdicción de los Estados extranjeros en la Argentina.

\section{Cuadro Cornell}

\begin{cornell}
\cpregunta{¿Por qué la caída de la inmunidad de jurisdicción no siempre implica la caída de la inmunidad de ejecución?}{Porque son institutos jurídicamente independientes en sistemas como el argentino: un Estado puede ser llevado a juicio y, sin embargo, sus bienes soberanos seguir protegidos de la ejecución forzosa.}
\cpregunta{¿Qué bienes del Estado son, en principio, no ejecutables?}{Los vinculados al ejercicio de la soberanía: edificios de gobierno, bienes militares, reservas del Banco Central, salvo renuncia expresa a esa inmunidad.}
\cpregunta{¿Qué resolvió la Corte Suprema en el caso Manauta?}{Adoptó la inmunidad relativa, calificando el contrato laboral entre la embajada rusa y sus empleados como un acto de gestión, no amparado por la inmunidad.}
\cpregunta{¿Qué norma codificó el criterio de Manauta?}{La Ley 24.488, que estableció legalmente la distinción entre actos de imperio y actos de gestión en materia de inmunidad de jurisdicción.}
\cresumen{La inmunidad de jurisdicción y la inmunidad de ejecución son institutos distintos: en Argentina, la Corte Suprema los trata de manera independiente. El caso Manauta consolidó jurisprudencialmente la inmunidad relativa del Estado extranjero, al calificar un contrato laboral de una embajada como acto de gestión no amparado por la inmunidad, criterio luego codificado en la Ley 24.488.}
\end{cornell}

% ----------------------------------------------------------------
\chapter{Inmunidad de los órganos de integración}
\chaptermark{Inmunidad de los órganos de integración}
% ----------------------------------------------------------------

\section{Acuerdos de sede e inmunidad de jurisdicción}

En un primer momento, los órganos del Mercosur no tenían una sede fija: se reunían en el país que ejercía la presidencia \emph{pro tempore} en cada período. Con el tiempo se crearon órganos con sede física permanente, como la Secretaría --en Montevideo-- y el Tribunal Permanente de Revisión --en Asunción-- (Capítulos~\ref{cap:estructura-mercosur} y~\ref{cap:tpr}).

\begin{definicion}{Acuerdo de sede e inmunidad de jurisdicción}
Cuando un órgano del Mercosur se instala con sede fija en un Estado parte, ese Estado firma con el bloque un acuerdo de sede que garantiza la inmunidad de jurisdicción del organismo. En consecuencia, las acciones que realice ese organismo y los miembros que intervengan en su nombre no pueden ser sometidos a la jurisdicción nacional del Estado sede.
\end{definicion}

\begin{importante}
La inmunidad de jurisdicción de un órgano de integración conlleva necesariamente la inmunidad de ejecución sobre sus bienes: no es posible afectarlos de ninguna manera mientras subsista el acuerdo de sede.
\end{importante}

\begin{ejemplo}{Reclamos laborales contra la Secretaría del Mercosur}
El personal administrativo --no diplomático-- de la Secretaría del Mercosur suele ser incorporado por quienes ocupan los cargos políticos, y cuando esos funcionarios dejan sus puestos, ese personal queda en una situación de incertidumbre laboral. En los primeros cambios de autoridades, ese personal comenzó a presentar reclamos porque no se les respetó la estabilidad ni se los indemnizó por el despido. Esos reclamos se dirigieron contra la Secretaría ante la justicia uruguaya (Estado sede), y el Tribunal Superior de Uruguay declaró que el Estado uruguayo no tenía competencia para someter a la Secretaría del Mercosur a la justicia local. A diferencia del caso Manauta (Capítulo~\ref{cap:manauta}), dentro del Mercosur un reclamo de este tipo no cuenta con una vía judicial equivalente, porque los particulares no tienen acceso directo al mecanismo de solución de controversias del bloque.
\end{ejemplo}

\section{La inmunidad en la Unión Europea: un régimen más preciso}

El Tratado de Lisboa remite, en esta materia, al Tratado de Funcionamiento de la Unión Europea (TFUE), que regula estas cuestiones de un modo más técnico y preciso que el esquema del Mercosur.

\begin{nota}
El TFUE distingue los actos de imperio de los actos de gestión de las instituciones europeas de un modo análogo al desarrollado para la inmunidad del Estado (Capítulo~35). % TODO: verificar la numeración exacta del artículo del TFUE y del Protocolo n.º 7 sobre los Privilegios y las Inmunidades de la Unión Europea antes de citarla en un trabajo académico.
\end{nota}

Esa norma establece que los jueces nacionales no pueden someter a juicio ni a las instituciones de la Unión Europea ni a sus funcionarios por los actos, opiniones, declaraciones y votos emitidos en el ejercicio de sus funciones y para el cumplimiento de los objetivos de la Unión. Cuando se realiza un \textbf{acto de imperio} --un acto cumplido en ejercicio de la función y para el objetivo para el que el órgano fue creado--, no es posible someter a su autor a un tribunal nacional; en cambio, las cuestiones privadas de los funcionarios sí pueden ser sometidas a jurisdicción.

\begin{ejemplo}{Levantamiento de inmunidad de un eurodiputado}
Si un eurodiputado comete un delito ajeno al ejercicio de su función, se solicita a la Unión Europea que le retire la inmunidad, precisamente porque no se trata de un acto vinculado al ejercicio de sus funciones. Levantada la inmunidad, el eurodiputado puede ser sometido a la jurisdicción nacional.
\end{ejemplo}

\begin{importante}
Existe una diferencia estructural muy importante entre el Mercosur y la Unión Europea: en la Unión Europea existe un órgano jurisdiccional permanente y único --el Tribunal de Justicia de la Unión Europea (Capítulo~\ref{cap:tjue})-- con competencia exclusiva para evaluar los actos institucionales, al que los particulares pueden acceder de manera directa mediante recurso directo. El Mercosur no ofrece una vía equivalente.
\end{importante}

\section{Cuadro Cornell}

\begin{cornell}
\cpregunta{¿Qué es un acuerdo de sede y qué efecto produce?}{El acuerdo que un Estado firma con el bloque cuando un órgano de integración se instala en su territorio, que garantiza la inmunidad de jurisdicción --y, por extensión, la de ejecución-- de ese organismo frente a la justicia local.}
\cpregunta{¿Qué mostró el caso de los reclamos laborales contra la Secretaría del Mercosur?}{Que el Tribunal Superior de Uruguay se declaró incompetente para juzgar a la Secretaría, porque los particulares no tienen, dentro del Mercosur, una vía judicial directa equivalente a la del caso Manauta.}
\cpregunta{¿Cómo regula el TFUE la inmunidad de los funcionarios europeos?}{Distinguiendo actos de imperio (protegidos por la inmunidad) de actos privados o personales (no protegidos), del mismo modo que se distinguen los actos de imperio y de gestión del Estado.}
\cpregunta{¿Cuál es la diferencia estructural clave entre el Mercosur y la Unión Europea en esta materia?}{La Unión Europea cuenta con un tribunal permanente y único, con acceso directo de los particulares; el Mercosur no ofrece una vía judicial equivalente para ellos.}
\cresumen{Los órganos de integración con sede fija gozan de inmunidad de jurisdicción y de ejecución en virtud de los acuerdos de sede que celebran con el Estado que los aloja, como ilustra el caso de los reclamos laborales contra la Secretaría del Mercosur en Uruguay. La Unión Europea regula esta materia de un modo más preciso, distinguiendo actos de imperio de actos privados de sus funcionarios, y ofrece a los particulares un acceso directo a su Tribunal que el Mercosur no ofrece.}
\end{cornell}

% ================================================================
\part{Fuentes, jerarquía y relación con el derecho interno}
% ================================================================

% ----------------------------------------------------------------
\chapter{Derecho de la integración y derecho comunitario}
\chaptermark{Derecho de la integración y derecho comunitario}
% ----------------------------------------------------------------

\section{Una relación de género a especie}

Esta distinción conceptual suele generar confusión, porque en el lenguaje corriente ambas expresiones se usan como sinónimos.

\begin{definicion}{Derecho de la integración}
Es todo el derecho que rige en cualquier tipo de área de integración: una zona de libre comercio, una unión aduanera, un mercado común o una comunidad. Es el \textbf{género}.
\end{definicion}

\begin{definicion}{Derecho comunitario}
Es una \textbf{especie} dentro del género derecho de la integración. Solo existe a partir de que el proceso de integración alcanza el estadio de comunidad (Capítulo~10).
\end{definicion}

Por eso es incorrecto hablar de «derecho comunitario del Mercosur»: el Mercosur es un mercado común, no una comunidad, de modo que el derecho que produce es derecho de la integración, pero no derecho comunitario en sentido técnico. El derecho comunitario del Mercosur no es lo mismo que el de la Unión Europea, no porque sean distintos entre sí, sino porque el Mercosur directamente no tiene derecho comunitario: tiene derecho de la integración, porque no reúne las características de una comunidad.

Recordando la escala de niveles de integración económica desarrollada en la Parte~II --integración fronteriza, zona de libre comercio, unión aduanera, mercado común, comunidad económica y unión--, cada una de esas etapas produce un derecho con características propias, y a partir del estadio de comunidad esas características cambian de manera cualitativa, no solo por la mayor complejidad institucional.

\section{Derecho originario y derecho derivado: una síntesis}

La distinción entre derecho originario y derecho derivado --desarrollada en detalle para el Mercosur en el Capítulo~\ref{cap:fuentes-mercosur}-- se aplica, con los mismos criterios, a cualquier área de integración: el derecho originario equivale, salvando las distancias, a la Constitución dentro de un Estado, y está integrado por el tratado marco o tratado constitutivo, junto con los tratados posteriores que crean nuevos órganos o modifican la estructura institucional; el derecho derivado son las normas que emiten los órganos creados por ese tratado, en ejercicio de las funciones que este les asignó.

Esta distinción se ubica también respecto de las ramas clásicas del derecho: mientras se está creando derecho originario, los Estados actúan en el plano del derecho internacional público propiamente dicho, porque están limitando su propia soberanía o acordando con otros Estados soberanos en un pie de igualdad. Recién cuando se pasa al derecho derivado, propio de cada área específica --comercial, laboral, de seguridad--, el derecho de la integración empieza a volverse más concreto y especializado.

\section{Cuadro Cornell}

\begin{cornell}
\cpregunta{¿Qué relación hay entre derecho de la integración y derecho comunitario?}{Es una relación de género a especie: el derecho de la integración rige en cualquier área de integración; el derecho comunitario es una especie que solo existe desde el estadio de comunidad.}
\cpregunta{¿Por qué el Mercosur no tiene, en sentido técnico, «derecho comunitario»?}{Porque el Mercosur es un mercado común, no una comunidad; produce derecho de la integración, pero no reúne las características que exige el derecho comunitario.}
\cpregunta{¿En qué plano del derecho se ubica la creación del derecho originario?}{En el del derecho internacional público, porque los Estados actúan limitando su propia soberanía o acordando entre sí en pie de igualdad.}
\cresumen{El derecho de la integración es el género, y el derecho comunitario una especie que exige que el proceso alcance el estadio de comunidad; por eso el Mercosur --mercado común-- no tiene, en sentido técnico, derecho comunitario. La distinción entre derecho originario (equivalente a la Constitución del área) y derecho derivado (normas de los órganos que ese tratado creó) atraviesa cualquier bloque de integración, no solo el Mercosur.}
\end{cornell}

% ----------------------------------------------------------------
\chapter{Derecho internacional y derecho interno: monismo y dualismo}
\chaptermark{Monismo y dualismo}
% ----------------------------------------------------------------

Este tramo final retoma una pregunta de base: ¿qué relación existe entre los tratados y el derecho interno? Antes de la reforma constitucional argentina de 1994 (Capítulo~\ref{cap:estructura-mercosur}), no existía en el texto de la Constitución una norma que estableciera con claridad qué jerarquía ocupaban los tratados internacionales dentro del ordenamiento jurídico. Frente a esa pregunta se plantean las dos teorías clásicas sobre la relación entre el derecho internacional y el derecho interno: el \textbf{monismo} y el \textbf{dualismo}.

\begin{importante}
Hablar de «relación», de «conflicto» o de «primacía» entre un tratado y una ley solo tiene sentido si se parte de la idea de que ambos pertenecen al mismo sistema jurídico. Esa es la premisa del monismo. El dualismo, en cambio, niega esa premisa desde el inicio.
\end{importante}

\section{La teoría dualista}

\begin{definicion}{Teoría dualista}
Sostiene que el derecho internacional y el derecho interno se mueven en dos ámbitos totalmente distintos, sin punto de contacto: nunca puede haber conflicto entre un tratado y una ley, porque cada uno se aplica en un ámbito distinto.
\end{definicion}

Para sostener esta separación, el dualismo identifica tres diferencias estructurales entre ambos órdenes:

\begin{enumerate}
    \item \textbf{Relación entre el órgano creador y el destinatario de la norma}: en el derecho internacional, quien dicta la norma --los Estados soberanos, en igualdad de condiciones-- es el mismo sujeto que se somete a ella. En el derecho interno, el órgano legislativo dicta la norma y el Estado, a través de sus habitantes, es quien queda sometido a ella: son sujetos distintos.
    \item \textbf{Ámbito territorial de aplicación}: las normas internacionales están destinadas a regular relaciones entre Estados; las normas internas regulan relaciones dentro de las fronteras de un Estado, entre ese Estado y sus ciudadanos.
    \item \textbf{Procedimiento de creación}: en el orden interno interviene el Poder Legislativo, que sanciona la norma, y el Poder Ejecutivo, que la promulga; si el Ejecutivo veta la norma, esta no llega a existir. En el orden internacional el procedimiento es distinto, como se explica en la sección siguiente.
\end{enumerate}

\begin{ejemplo}{La casa propia y la del vecino}
Ambos sistemas no admiten comparación entre sí, porque no hay una jerarquía común entre ellos: es la casa propia y la casa del vecino, y en cada una manda quien la habita. Si los dos sistemas son autónomos, nunca puede existir en rigor un conflicto entre un tratado y una ley: sencillamente rigen en órbitas distintas.
\end{ejemplo}

\section{La teoría monista y el procedimiento complejo}

\begin{definicion}{Teoría monista}
Sostiene que el sistema de elaboración de las normas internacionales es un sistema complejo, cuya vigencia se construye en varios pasos sucesivos. Parte de la premisa de que el tratado y la ley pertenecen al mismo sistema jurídico.
\end{definicion}

\begin{importante}
\textbf{Procedimiento complejo de elaboración de un tratado (sistema argentino)}
\begin{enumerate}
    \item \textbf{Firma del tratado por el Poder Ejecutivo}: en ese momento el Estado asume el compromiso internacional, pero el tratado todavía no existe como norma dentro del ordenamiento jurídico argentino. % TODO: verificar consistencia entre fuentes originales respecto del alcance exacto de la responsabilidad internacional derivada de la sola firma.
    \item \textbf{Ratificación por el Congreso y depósito del instrumento de ratificación}: recién cuando intervienen los dos órganos que participan en la elaboración de la norma (Ejecutivo y Legislativo), el tratado queda vigente tanto en el ámbito internacional como en el ámbito interno, a la vez.
\end{enumerate}
\end{importante}

Hasta que no se completa la ratificación, la norma directamente no existe para el ordenamiento argentino, porque no cumplimentó todo el procedimiento de elaboración. Hay una diferencia de orden respecto del procedimiento interno: en la elaboración de una ley interna primero interviene el Legislativo y después el Ejecutivo promulga; en el sistema de tratados es al revés, primero interviene el Ejecutivo --con la firma-- y después el Congreso ratifica. En ambos casos la norma solo entra en vigor cuando intervinieron los dos órganos, y a partir de ese momento rige en los dos ámbitos, el internacional y el interno, simultáneamente.

\begin{importante}
\textbf{«Ratificar» y no «incorporar».} Dentro del sistema monista argentino no corresponde decir que «el tratado se incorpora» al derecho interno. Esa expresión pertenece al vocabulario dualista, porque presupone dos sistemas separados entre los cuales hay que trasladar la norma de uno a otro. En el sistema monista argentino, el Congreso \textbf{ratifica} el tratado, no lo incorpora, porque el tratado y la ley conviven desde el inicio en el mismo sistema jurídico.
\end{importante}

\begin{nota}
La doctrina reconoce, dentro del monismo, distintas variantes según la jerarquía que se asigna a la relación entre el derecho internacional y el derecho interno: monismo con igualdad jerárquica entre ambos órdenes, monismo con supremacía del derecho internacional y monismo con supremacía del derecho interno. Cuál de estas variantes rige en cada Estado --y cómo se aplica, en particular, al derecho de la integración-- excede el desarrollo de este libro, pero constituye la pregunta natural en la que continúa este último tramo de la materia.
\end{nota}

\section{Cuadro Cornell}

\begin{cornell}
\cpregunta{¿Qué sostiene la teoría dualista sobre la relación entre tratado y ley?}{Que nunca puede haber un conflicto real entre ambos, porque pertenecen a sistemas jurídicos distintos: distinto órgano creador y destinatario, distinto ámbito territorial de aplicación y distinto procedimiento de creación.}
\cpregunta{¿Por qué solo tiene sentido hablar de conflicto o primacía entre un tratado y una ley si se parte del monismo?}{Porque esas nociones presuponen que ambos pertenecen al mismo sistema jurídico; el monismo parte de esa premisa, y el dualismo la niega desde el inicio.}
\cpregunta{¿Cuál es el procedimiento complejo de elaboración de un tratado en el sistema monista argentino?}{Firma del Poder Ejecutivo (asume el compromiso internacional) y ratificación del Congreso con depósito del instrumento de ratificación; recién con ambos pasos el tratado entra en vigor a la vez en el ámbito internacional y en el interno.}
\cpregunta{¿Por qué es incorrecto decir que Argentina «incorpora» un tratado?}{Porque «incorporar» presupone dos sistemas jurídicos separados entre los cuales hay que trasladar la norma (vocabulario dualista); en el sistema monista argentino, el Congreso ratifica el tratado, que convive con la ley en el mismo sistema desde el inicio.}
\cresumen{Para el dualismo, el derecho internacional y el interno son órdenes independientes, por lo que no puede haber conflicto entre tratado y ley. Para el monismo --el sistema argentino--, ambos integran un mismo sistema jurídico, y la vigencia del tratado se construye mediante un procedimiento complejo: firma del Ejecutivo y ratificación del Congreso. En ese marco se ratifica un tratado, y no se lo incorpora.}
\end{cornell}

% ================================================================
\part{Síntesis final y preparación del examen}
% ================================================================

% ----------------------------------------------------------------
\chapter{Método de estudio y modalidad de examen}
\chaptermark{Método de estudio y examen}
% ----------------------------------------------------------------

\section{Modalidad de evaluación}

La modalidad de evaluación de la cátedra combina dos formatos: preguntas cerradas (verdadero/falso y opción múltiple) y preguntas a desarrollar.

\begin{table}[H]
\centering
\small
\renewcommand{\arraystretch}{1.3}
\begin{tabularx}{\textwidth}{>{\raggedright\arraybackslash}p{0.3\textwidth}>{\raggedright\arraybackslash}X}
\toprule
\textbf{Formato} & \textbf{Efecto sobre la evaluación} \\
\midrule
Preguntas a desarrollar (habitualmente dos) & Centradas en el régimen de solución de conflictos de la Unión Europea o del Mercosur, o en el derecho sustantivo de uno u otro bloque. Ninguna es eliminatoria por sí sola: en conjunto representan aproximadamente un punto y medio sobre diez. \\
Entre 10 y 15 preguntas de verdadero/falso y opción múltiple & Distribuidas sobre todos los temas del programa. Cada pregunta vale poco, por lo que ningún tema puntual compromete por sí solo la nota; lo que sí la compromete es no dominar un bloque completo de contenidos. \\
\bottomrule
\end{tabularx}
\end{table}

La desventaja de este sistema, tanto para los estudiantes como para la cátedra, es que se pregunta \textbf{absolutamente todos los temas} dados en la materia, sin excepción.

\begin{importante}
Las respuestas incorrectas \textbf{no restan puntaje}, pero, como contrapartida, el umbral para aprobar es más exigente que en un examen tradicional: se requiere el \textbf{60\,\%} del parcial correctamente respondido para llegar al cuatro (nota mínima de aprobación). La razón: si se aplicara la conversión directa sobre cien puntos, un 40\,\% equivaldría a un cuatro, lo que significaría que se sabe menos de lo que se ignora. Superado el umbral inicial, la nota sube rápidamente (60\,\% equivale a 4; 66\,\% a 5; 74\,\% a 6).
\end{importante}

Habrá dos instancias parciales a lo largo de la cursada, cada una con su propio recuperatorio en la misma modalidad, y en el recuperatorio también se puede promocionar la materia. Las notas de ambas instancias se promedian entre sí para la promoción: si en el parcial se obtuvo un dos y en el recuperatorio un ocho, la nota final es un ocho. El primer parcial evalúa fundamentalmente la primera parte de la materia --relaciones de integración en el ámbito público, entre Estados--; el segundo evalúa la segunda parte --relaciones de integración que afectan a los particulares--.

\section{Qué priorizar y qué no memorizar}

No se exige memorizar la cantidad de miembros de cada órgano, pero sí comprender \textbf{cómo se toman las decisiones} en cada uno --por consenso o por mayoría-- y \textbf{a quién representa} cada órgano.

\begin{ejemplo}{Trampa conceptual en una pregunta de verdadero o falso}
Afirmación: «el Consejo representa a la Unión Europea». La respuesta es \textbf{falsa}, porque quien representa a la Unión Europea es la Comisión Europea. El Consejo Europeo está integrado por los presidentes y mandatarios de los Estados, el Consejo de Ministros representa a los Estados y el Parlamento representa al pueblo (Capítulo~\ref{cap:organos-ue}).
\end{ejemplo}

\section{Estrategias de lectura comprensiva}

La pregunta del examen debe tomarse \textbf{literalmente}, sin completarla mentalmente con lo que se cree que falta. La distinción entre una enumeración cerrada y una enumeración abierta es, en este sentido, una cuestión de lectura comprensiva: conviene tomarse dos minutos para leer bien antes de responder.

\begin{ejemplo}{Enumeración cerrada y enumeración abierta}
Si se afirma «los únicos órganos con capacidad decisoria son estos» y falta uno, la afirmación es falsa, porque se dijo «únicos». Si en cambio se afirma «este, este y este tienen capacidad decisoria», aunque falte nombrar otro, la afirmación es verdadera, porque no se dijo que fueran los únicos.
\end{ejemplo}

\begin{importante}
\textbf{Criterio de corrección de las preguntas a desarrollar.} No se premian la extensión ni las frases genéricas: repetir la pregunta con otras palabras no suma nada. Lo que suma es identificar con precisión cada característica, enumerada punto por punto, y explicar qué significa cada una. También se penaliza copiar, en una pregunta a desarrollar, fragmentos de enunciados que aparecieron en las preguntas de verdadero o falso, sin razonar la respuesta propia.
\end{importante}

\section{Material de estudio y metodología de la cursada}

La cátedra no cuenta con un manual único obligatorio: los libros disponibles en el mercado suelen estar escritos desde la experiencia de la Unión Europea, que no es equiparable a la de otros bloques regionales como el Mercosur --un libro que describe «la integración» puede estar describiendo, en rigor, solo a la Unión Europea--. El material se pone a disposición a través del campus virtual, y los tratados internacionales relevantes se indican oportunamente, porque es necesario buscarlos y analizarlos.

Se espera una dinámica de participación y de cuestionamiento, con debate respetuoso sobre temas de actualidad política y económica vinculados a la integración regional, sin ofender a nadie.

\section{Cuadro Cornell}

\begin{cornell}
\cpregunta{¿Qué umbral de respuestas correctas se requiere para aprobar el parcial?}{El 60\,\%, equivalente a la nota mínima de aprobación (cuatro); a partir de allí, la nota sube rápido.}
\cpregunta{¿Qué se prioriza comprender sobre los órganos de cada bloque, en lugar de memorizar?}{Cómo se toman las decisiones en cada órgano (consenso o mayoría) y a quién representa cada uno.}
\cpregunta{¿Por qué es clave distinguir una enumeración cerrada de una abierta al leer una pregunta de examen?}{Porque una enumeración cerrada («los únicos») es falsa si falta un elemento; una enumeración abierta puede ser verdadera aunque no mencione todos los elementos.}
\cpregunta{¿Qué se valora en una pregunta a desarrollar?}{Identificar con precisión cada característica, punto por punto, y explicar su significado; no la extensión ni las frases genéricas.}
\cresumen{El examen combina preguntas cerradas sobre todo el programa con un número reducido de preguntas a desarrollar, bajo un umbral de aprobación exigente (60\,\%) que compensa la ausencia de penalización por error. Conviene priorizar la comprensión del funcionamiento de cada órgano --cómo decide y a quién representa-- por sobre la memorización de datos numéricos, leer cada enunciado literalmente y, en las preguntas a desarrollar, responder con precisión y en forma estructurada.}
\end{cornell}

% ----------------------------------------------------------------
\chapter{Síntesis integradora general}
\chaptermark{Síntesis integradora}
% ----------------------------------------------------------------

Al cabo de este recorrido, conviene reunir en un solo lugar las líneas que atraviesan todo el libro y que, en definitiva, son las que un examen final evalúa de manera integrada.

\section{El eje conceptual: qué es integrar}

El Derecho de la Integración es una rama autónoma, distinta del Derecho Internacional Público y del Derecho Internacional Privado, que regula tanto relaciones entre Estados como relaciones que afectan a los particulares, pero de una naturaleza propia (Parte~I). La integración se distingue de la simple cooperación por el compromiso recíproco y el objetivo común, y todo tratado de integración debe leerse distinguiendo con cuidado entre lo que reconoce, lo que garantiza y lo que promueve, y entre el objetivo perseguido y el medio elegido para alcanzarlo.

\section{Dos escalas paralelas}

A lo largo del libro conviven dos escalas de profundización que conviene no confundir: la escala de \textbf{niveles de integración económica} --zona franca, unión fronteriza, zona de libre comercio, unión aduanera, mercado común, comunidad económica y unión económica (Parte~II)-- y la escala de \textbf{modelos institucionales} --órganos intergubernamentales, que deciden por consenso y representan a los gobiernos, frente a órganos supranacionales, que deciden por mayoría y son independientes de los Estados que los designaron (Capítulo~\ref{cap:atribucion})--. El Mercosur y la Unión Europea permiten observar ambas escalas en funcionamiento simultáneo, pero con resultados muy distintos.

\section{Dos casos, un mismo espejo}

El Mercosur --nacido de un acercamiento bilateral entre Argentina y Brasil dentro del marco de la ALADI (Partes~III y~IV)-- planteó desde el Tratado de Asunción la meta de un mercado común, pero conservó una estructura predominantemente intergubernamental: sus órganos decisorios deciden por consenso, sus innovaciones más recientes --el PARLASUR y el Tribunal Permanente de Revisión-- conservan facultades limitadas, no vinculantes o restringidas a la revisión de lo ya decidido, y persiste una tensión permanente entre los compromisos asumidos frente a los socios y las decisiones de política interna de cada Estado.

La Unión Europea (Parte~V), en cambio, avanzó hacia una supranacionalidad real, pero acotada casi exclusivamente al pilar económico, comercial y financiero: en seguridad, justicia, política exterior y defensa vuelve a comportarse, en los hechos, de un modo cercano al del Mercosur, decidiendo por consenso o unanimidad. Los casos de tensión estudiados --la crisis migratoria mediterránea, la libertad religiosa en España, la expulsión de comunidades gitanas de Francia, el propio Brexit-- muestran que ningún bloque, por más avanzado que sea su diseño institucional, está exento de la distancia entre lo que declara y lo que efectivamente hace.

\begin{importante}
La lección metodológica que atraviesa toda la materia es que el grado real de cesión de soberanía de un Estado a un organismo internacional nunca debe presumirse a partir del nombre de una institución o del texto declarativo de un tratado, sino que debe verificarse concretamente observando quién decide, cómo decide y qué efecto tiene esa decisión sobre el derecho interno.
\end{importante}

\section{La solución de controversias como prueba de fuego}

Esa misma lección se verifica con particular nitidez en el sistema de solución de controversias de cada bloque (Parte~VI). La Unión Europea garantiza la aplicación uniforme de su derecho mediante un Tribunal de Justicia único, con múltiples vías de acceso --incluida la de los particulares-- y una interpretación siempre obligatoria. El Mercosur, en cambio, limita el acceso a la Corte Suprema de cada Estado, no vincula a esa Corte con la opinión del Tribunal Permanente de Revisión, y permite a los Estados acudir a otros foros, lo que le resta autoridad a su propio sistema. De allí la conclusión que se repite a lo largo del libro: el Mercosur se comporta, en la práctica, más como una relación internacional clásica que como un espacio de integración plenamente consolidado.

\section{Inmunidad y fuentes: el cierre del sistema}

La inmunidad del Estado (Parte~VII) evolucionó de la inmunidad absoluta --fundada en el principio \emph{par in parem non habet imperium}-- a la inmunidad relativa, que distingue actos de imperio de actos de gestión, consolidada en la Argentina a partir del caso Manauta. Esa misma lógica de distinción entre lo soberano y lo gestional se proyecta sobre la inmunidad de los órganos de integración frente al Estado que los aloja. Finalmente, comprender la jerarquía de fuentes --derecho originario y derivado, derecho de la integración y derecho comunitario-- y la relación entre el derecho internacional y el derecho interno --monismo y dualismo (Parte~VIII)-- es la base indispensable para saber, frente a un caso concreto, qué norma corresponde aplicar, con qué jerarquía y con qué grado de exigibilidad.

\section{El hilo que conecta todo}

Las piezas estudiadas en este libro --órganos y su naturaleza intergubernamental o supranacional, fuentes normativas y su jerarquía, sistemas jurisdiccionales y su alcance, inmunidad y su relación con la soberanía-- conforman un mismo mapa conceptual. La pregunta que subyace a todas ellas es siempre la misma: ¿hasta qué punto un Estado está dispuesto a que otro decida sobre lo que antes decidía él solo? La respuesta nunca es igual en todos los planos ni en todos los bloques, y precisamente esa variación --entre lo económico y lo político, entre el discurso y la práctica, entre lo declarado y lo efectivamente cedido-- es lo que este libro intentó mostrar en cada una de sus partes.

% ================================================================
\chapter*{Bibliografía y normativa citada}
\addcontentsline{toc}{chapter}{Bibliografía y normativa citada}
\markboth{Bibliografía y normativa citada}{Bibliografía y normativa citada}
% ================================================================

Esta lista reúne, sin agregar referencias externas a las fuentes originales, los tratados, protocolos, normas y precedentes jurisprudenciales efectivamente citados a lo largo del libro.

\subsection*{Tratados y protocolos de la integración latinoamericana}
\begin{itemize}
    \item Tratado de Montevideo (1960) --- ALALC (Asociación Latinoamericana de Libre Comercio).
    \item Tratado de Montevideo (1980) --- ALADI (Asociación Latinoamericana de Integración).
    \item Acuerdo de Cartagena (1969) --- Grupo Andino, hoy Comunidad Andina.
    \item Declaración de Foz de Iguazú (30 de noviembre de 1985) --- Argentina y Brasil.
    \item Tratado de Asunción (1991) --- creación del Mercosur.
    \item Protocolo de Brasilia (1991) --- sistema de solución de controversias del Mercosur.
    \item Protocolo de Ouro Preto (1994) --- estructura institucional definitiva del Mercosur.
    \item Protocolo de Ushuaia --- cláusula democrática del Mercosur.
    \item Protocolo de Olivos (2002) --- Tribunal Permanente de Revisión.
    \item Protocolo de Tucumán --- libre tránsito laboral en el Mercosur.
\end{itemize}

\subsection*{Tratados de la Unión Europea}
\begin{itemize}
    \item Tratado de París (1951) --- Comunidad Europea del Carbón y del Acero (CECA).
    \item Tratados de Roma (1957) --- Comunidad Económica Europea (CEE) y EURATOM.
    \item Tratado de Maastricht (1992).
    \item Tratado de Ámsterdam (1997).
    \item Tratado de Lisboa (2007) --- vigente, incorpora las reformas institucionales.
    \item Tratado de Funcionamiento de la Unión Europea (TFUE) --- inmunidad de instituciones y funcionarios.
    \item Protocolo n.\textordmasculine{}~7 sobre los Privilegios y las Inmunidades de la Unión Europea. \% TODO: verificar numeración exacta de los artículos citados.
\end{itemize}

\subsection*{Normativa argentina}
\begin{itemize}
    \item Constitución Nacional, artículo 2 --- sostenimiento del culto católico.
    \item Constitución Nacional, artículo 14 --- libertad de culto.
    \item Constitución Nacional, artículo 75, inciso 24 --- tratados de integración y delegación de competencias.
    \item Ley 1285 (1958) --- inmunidad de jurisdicción de los Estados extranjeros (régimen anterior a Manauta).
    \item Ley 24.488 --- codificación de la inmunidad relativa del Estado extranjero.
\end{itemize}

\subsection*{Jurisprudencia}
\begin{itemize}
    \item \emph{Manauta c/ Embajada de la Federación Rusa} --- Corte Suprema de Justicia de la Nación (Argentina) --- inmunidad relativa del Estado extranjero.
\end{itemize}

\end{document}
